% ============================================================
%  Rasmus Nielsen, "Almindelig Videnskabslære i Grundtræk"
%  "General Theory of Science in Outline"
%  Kjøbenhavn: Gyldendalske Boghandel (F. Hegel & Søn), 1880.
%  TRANSLATION (English)
%
%  German, Latin, French, and Greek quotations are left in the
%  original, with an English gloss in brackets on first occurrence
%  where helpful. Nineteenth-century technical vocabulary is
%  rendered to preserve Nielsen's distinctions: Viden = "knowledge"
%  (as cognition), Videnskab = "science", Aand = "Spirit",
%  Sandsning = "sensation", Anlæg = "endowment/disposition",
%  Realbegreb = "real-concept", Selvhed = "selfhood".
% ============================================================
\documentclass[12pt]{book}
\usepackage[utf8]{inputenc}
\usepackage[T1]{fontenc}
\usepackage[english]{babel}
\usepackage[protrusion=true,expansion=true]{microtype}
\usepackage[margin=1.4in]{geometry}
\usepackage{setspace}
\usepackage{amsmath}
\usepackage{libertinus}
\usepackage{libertinust1math}
\usepackage{textalpha}   % occasional polytonic Greek (e.g. the Homer quotation in \S\,10)
\usepackage{fancyhdr}
\usepackage[colorlinks=true,linkcolor=black,urlcolor=black,
  pdftitle={General Theory of Science in Outline},
  pdfauthor={Rasmus Nielsen}]{hyperref}
\usepackage{chngcntr}
\counterwithout{section}{chapter}
\renewcommand{\thesection}{§\,\arabic{section}}
\setlength{\headheight}{15pt}
\pagestyle{fancy}
\fancyhf{}
\fancyhead[LE,RO]{\thepage}
\fancyhead[RE]{\textit{General Theory of Science in Outline}}
\fancyhead[LO]{\textit{\leftmark}}
\setlength{\parindent}{1.5em}
\setlength{\parskip}{0pt}
\onehalfspacing

\title{General Theory of Science in Outline}
\author{R.\ Nielsen}
\date{Copenhagen: Gyldendal's Bookshop (F.\ Hegel \& Son), 1880\\[6pt]
  \small Translated from the Danish}

\begin{document}
\maketitle
\tableofcontents

\chapter*{Preface}
\addcontentsline{toc}{chapter}{Preface}

\noindent Although the sciences differ among themselves, the difference
between them in principle --- grounded in the relation between sensation
and thought --- cannot be so great that they should not all, the
\textit{a~priori} as well as the empirical, be subsumable under a common
fundamental principle valid for knowledge in general. The question then
becomes: from what does the principle of knowledge essentially derive;
whether it rests originally upon sensation, or also upon a thinking that is
decisive for the things themselves; whether it is grounded upon matter and
force alone, so that its validity must at most be confined to categories
such as cause and effect; or whether it has its root in the idea of
infinite knowledge, so that fundamental determinations such as endowment
and development, means and ends, are seen therein to have objective
validity. If the latter is assumed, the assumption must be justifiable from
the mutual connection of the sciences; for such a connection must
essentially be determined by the way in which the principle develops. That
the general theory of ideas can be cultivated for itself, geology for
itself, history for itself, and so on, cannot be denied; but the sciences
must nonetheless all hang together at bottom: if the one can do without the
other, then there is no connection here at all.

The investigation upon which, at the present time, particular weight must
be laid in the theory of science concerns the relation between the physical
and the biological. This relation may, with special regard to the Darwinian
hypothesis, be conceived in two ways: partly such that everything
teleological disappears, and only the mechanical-physical causes are
acknowledged; partly such that the physical is subject to a natural
teleological form, and self-development shows itself to satisfy the
endowment. As acknowledgement of the Darwinian foundation spreads more and
more, the dispute between the merely physical and the
physical-teleological conception --- not of Darwin's personal view, but of
the applicability of the foundation --- will come to a breakthrough; for
those remnants of teleology that seek shelter and refuge in the modes of
view diverging from Darwinism will scarcely have any hope of being able to
hold their own against a theory which otherwise, in the whole as in the
particular, tests, takes up, and appropriates all the results of research.
The defects of Darwinism, be they now in the ``ontogeny'' or in the
``phylogeny'' or wherever, must be capable of being discovered and remedied
in consequence of the fundamental view itself, not by additions coming from
without.

As regards the cognition of the religious, the spiritual life, in so far as
it develops toward a true understanding of itself, will not fail to submit
to the control of science in everything that depends especially upon the
outward actual; but in relation to its inner essence the religious life
will still, as before, continue to form itself out of its own spirit, after
its own exemplars, and go its way independently of science and of
scientific objections.

This is, in brief, the character of the investigations carried out in the
present work. How far, in these investigations --- setting aside
expressions and turns of phrase such as readily present themselves to one
who, in the nature of the case, cannot possibly be as much at home in the
\textit{termini} of the several special sciences as their actual
practitioners --- there may perhaps be contained something which on closer
examination might prove to be of more general significance, is hereby left
to expert and well-disposed judges to decide.

\bigskip
\noindent Copenhagen, the 23rd of August.
\nopagebreak
\par\hfill R.\ Nielsen.


%%% PART ONE: KNOWLEDGE AND SCIENCE %%%
\part{Knowledge and Science}

\chapter{The Essence of Knowledge in General}

\section{Original Presuppositions}
\label{sec:1}

\noindent Our first knowledge comes, so to speak, of itself. We have
knowledge of many things before we have any inkling of what knowledge
properly is. Only when it dawns upon us that there is something we do not
know, and yet would gladly know --- something about which we are in
uncertainty, although we desire certainty --- do we begin to lay emphasis on
the word knowledge. And yet there is still a long way to scientific
knowledge; we must take preparatory steps. The history of culture teaches
that every people has had to pass through a series of different stages of
development, and to gather a rich material of varied practical knowledge,
before it reached the point where an actual science of the matter could
begin. A knowledge of the course of nature, of heaven and earth, of gods and
men, of family life, state, and moral order may in certain respects be
considerably developed without there corresponding to it any scientific
understanding. What holds of peoples holds also of individuals. Organic
morphology teaches us that individuals of a higher order repeat in the
embryonic state, though with modifications and abridgements, the formations
that are characteristic of the lower and lowest forms from which the species
descends. Something similar holds of the development of consciousness in the
single human being. As a child it is awakened to distinguish outward things
and to feel its own existence; then it learns to use language; it enriches
itself with a store of preliminary experiences, and reaches at last the
point where a regular schooling can begin. Yet the instruction that prepares
a scientific cognition is not itself yet scientific. Items of knowledge are
gathered and ordered; but only little by little does the immediate
consciousness come to reflect upon the original presuppositions of its
knowledge.

Here is required a clear distinction between the circle of images and
representations that belong to consciousness itself, and the things in the
outward world. The individuality, with its thinking I, must be able, with
sure tact according to given usage, to speak of forms such as time and
space; it must be able to distinguish between the thing and its properties,
matter and its forces; it must be able to understand what connection is, what
is meant by expressions such as cause and effect, contingent and necessary,
and so forth. On these conditions the person in question will be able to
begin a real investigation of existing things and their laws; he will, if he
has powers and natural gifts, be able to work up a given material with
scientific skill and arrive at valuable results. But only when he turns his
attention to the conditions themselves, to the original presuppositions of
consciousness, does he come to know what knowledge and science essentially
are.

A remarkable change comes over him. What before was immediately evident now
suddenly becomes enigmatic. Occupied with the mutual relations of things and
their motions in space and time, he was wholly taken up with particular,
exact determinations; what space and time are seemed to him to be understood
of itself. But now questions such as these arise: What properly is space,
what is time as time? Do these determinations of form concern the things
themselves, or are they merely necessary for our apprehension? What is
matter, what is force, what is a thing when we separate it from its
properties, a whole when it is separated from its parts? We seem to see the
material, but the forces we cannot see; we seem to sense the thing, but it is
only its properties; we observe the changes of things, and their succession
with one another, after one another; but what entitles us to assume the
changes as effects of causes lying in the things? What properly is that which
we call cause --- is it something sensible or merely something thinkable,
something that is in existence outside us, or something that is merely in our
own brain? When such questions force themselves upon consciousness, when the
foundation of the first natural and immediate knowledge wavers, there is
awakened in the mind a need to investigate the essence of knowledge and
science.

As the thinker in question begins to test the fundamental propositions of
his knowledge, to doubt everything, in order to work his way through doubt
and uncertainty to grounded certainty, he is in danger of going astray.
Whatever he chooses as a starting-point, he easily finds that every foothold
gives way. It seems to him as though there could be no doubt of the reality
of material objects; but what certainty have we that the objects exist
outside our consciousness? We are referred to our own feeling, our sight, our
hearing --- in short, to our own sensations, our own states of consciousness,
our own subjective, necessary apprehension. So he then seeks in consciousness
itself for something fixed, and supposes perhaps that in the proposition ``I
think, and thinking I am'' he has an unshakable standpoint; but he will soon
find himself disappointed. What we call our I is indeed a unity, an inner
something that accompanies all our representations and thoughts; but what
sort of something is it? In so far as it is a unity of consciousness, it must
be just as changeable as consciousness; but that consciousness, being subject
to organic conditions, can be dimmed, darkened, indeed even wholly
disappear, is an incontrovertible fact. He attempts to grasp the absolutely
certain in thinking and its laws; but he fares no better. Is there any
proposition of thought that contradiction cannot assail? Yes, here is the
fundamental principle of likeness, of identity: A is what A is, $A = A$. But
no --- this too is a deception. The proposition $A = A$ rests upon the
proposition I = I, and this again upon the presupposition consciousness =
consciousness. But consciousness is changeable; if, then, it changes so
rapidly that it is in one state when it posits A the first time, and in
another and different state when it posits A the second time, then it becomes
quite uncertain whether in the A posited twice we really have the selfsame
thing; thus the proposition $A = A$ becomes uncertain.

What now follows from this? That the individuality seeking knowledge and
certainty finds uncertainty everywhere and ends in bottomless skepticism. To
continue on this path and take the matter seriously leads to despair. The
seeking individuality must halt; and it now depends upon how he breaks off.

If he returns to the first immediacy, uses the generally valid
representations by tact, and settles into peace with so-called common sense,
then he can develop his understanding in a practical direction, enrich
himself with knowledge --- indeed, he can even study some field or other with
success, make discoveries, and carry things very far; but what scientific
knowledge properly is, he never comes to know. There have been esteemed men
of science, distinguished by rare sense and acuteness, who, so long as the
talk was of a circle of definite phenomena and laws, could express
themselves with all the assurance of one who knows, but who, when it came to
the general foundation of knowledge, to the validity that may on the whole be
attributed to human knowledge, showed themselves ensnared in skeptical
confusion.

Knowledge and science stand in so essential a relation to one another that
they presuppose each other in the infinite. One cannot become fully at home
in any science whatever, be it empirical or \textit{a~priori}, without
giving oneself an account of what validity human knowledge has, and how the
categories of knowledge that are continually applied originally stand related
to thing and consciousness. Nor will one succeed in procuring beforehand
certainty about what knowledge properly is, and how far there is a true
knowledge, before one engages with science itself; for if the investigations
concerning the worth of knowledge are not carried through in a strictly
scientific manner, one arrives at false results. The blind reliance on fixed
grounds and the loose doubting that hollows out all grounds are examples of
the unscientific.

Knowledge is a light of consciousness; the point is to reach the evident.
Just as in the world of things there are dark and self-luminous bodies, so in
consciousness there is something obscure that must be illuminated by
something else, and something that is in itself evident. We must then begin
with the immediately evident and pass from it to what lies nearest, which can
be illuminated thereby. The connection between the one and the other must
itself be immediately evident. But already this first step is difficult. In
order to investigate what knowledge is, we must begin with knowledge. Now if
our knowledge is evident of itself, then there is nothing here to
investigate; if an investigation of knowledge itself is to have any meaning,
then there must be something else that is in and of itself more evident than
knowledge; but what could this be? A knowing self, an object of knowledge,
and the union of both in one consciousness!

These presuppositions by no means have the evidence of outward things that
present themselves to sight; nevertheless they prove, on closer
investigation, to belong to what is first evident of itself. One convinces
oneself of this by attempts to explain them from something else and more
evident, or from one another. If one tries to derive knowledge from
non-knowledge, one easily sees the deception, for the distinction of
knowledge and non-knowledge must take place within knowledge. Just as
meaningless is it to wish to derive the conscious from the unconscious; for
one must necessarily have the unconscious in consciousness merely in order to
be able to compare it with the conscious. If one tries to hold fast knowledge
for itself without the admixture of any other determination whatever, it ends
with a knowledge of nothing; but a knowledge of nothing is non-knowledge,
and non-knowledge then becomes a negative form for the object of knowledge.
If one tries to explain knowledge from the knower, then one has already
slipped in knowledge; and just as impossible is it to derive the knowing self
from knowledge; for in order to be able to dispense with a self that has
knowledge, knowledge would have to be of itself knowing about itself, and
precisely thereby distinguish knowledge of itself as knowing from knowledge
of everything else. However empty and formal such an original discerning and
distinguishing may seem to be, it is nonetheless not idle; on the contrary,
it is the first and only way to attain scientific insight into what is
evident of itself.

What has here been said of consciousness and knowledge in pure generality
holds also of the particular forms in which consciousness receives its
content --- of the original forms of sensation, space and time, of the forms
of the understanding or of thought, the logical categories, and so forth.
These are not self-subsisting departments of knowledge that allow themselves
to be arranged in certain rubrics, as though one could have sensation in one
drawer and the understanding in another; they are, on the contrary, active
forms that mutually condition one another, and which, in reciprocal
cooperation, enter into so manifold combinations and thereby yield such
peculiar determinations that analysis must everywhere proceed from the
resolution of the composite to the simple: understanding is mediate.

Nevertheless there is in each of these particular universal forms something
peculiar that is evident of itself. In however different ways one may seek to
ascertain the characteristic determinations of space and time and to show how
and under what conditions these determinations of form arise and develop in
consciousness, it is yet incontrovertibly certain that the peculiar element
in each of them can be discerned only from a result of development --- by the
fact that already at the beginning of the development it is present as
something evident of itself. By space we explain space, by time time, by
ground ground, by cause cause; the explanation moves in a circle.

And yet the explanation is by no means idle; for in the course of the
circular movement it happens that the circle for each of the fundamental
forms widens, in that, solely for its own sake, it takes up into itself
determinations of the others. Thus space is determined by time, time by
space, sensation by thought, and thought again by sensation. Thus
consciousness and knowledge are explained by means of particular forms, and
these again by means of consciousness and knowledge. Thus the fundamental
forms of consciousness are explained by a corresponding content, and the
determinations by the corresponding content of the fundamental forms of
consciousness. The scientific demand is satisfied by the fact that no form,
no element of content, no object is allowed to remain standing isolated or to
be joined to the others merely in an outward, contingent manner. All
determinations --- determinations of form and determinations of content,
sensory determinations and thought-determinations, original and derived ---
must enter into the movement in such a way that point by point they
illuminate and shine through one another reciprocally.

\section{The Share of Sensation in Knowledge: Outward Experience}
\label{sec:2}

\noindent Although it is evident of itself that knowledge must have its
object, it by no means follows from this that one can derive the constitution
of the object from the general demand of knowledge. Object and consciousness
are far from resembling each other; the opposition is so prominent that to an
immediate glance it looks vividly as though the object stood independently
outside consciousness. Now is this so? In confirmation of the outward
independence of the object one leans upon the testimony of the senses; but
the testimony is ambiguous. If one conceives the circle of consciousness as a
circle with a definite radius, and lets the things stand in a space of their
own outside the periphery, how then should they be capable of being sensed?
It looks as though they were mirrored in consciousness and related to it as
objects to the mirror. That this is an erroneous conception is easily seen. In
order to show itself as a sensible object, the object must be sensed and
thereby taken up into consciousness itself. The whole sensed object is
situated within the circle of the sensing consciousness; if there is anything
outside, then it must here be the non-sensible, or at any rate the unsensed;
but of that the sensing consciousness can know nothing. If the sensing
subjectivity could really see the objects outside its consciousness in space,
then it would also have to be able to see their different distances from the
boundary of consciousness; the Pleiades, for example, would then have to show
themselves to be farther away --- not from the sensing individual, for that
goes without saying, but from consciousness itself --- than the moon.

Nevertheless we must grant common sense that the presupposition of the
existence of objects outside consciousness has a real meaning. For if one
assumes that the objects themselves are situated within consciousness, then
one must also assume that they are the sensing consciousness's own products.
But if things exist only in the moment we sense them, then they must cease to
exist when we do not sense them; one who is attacked by a robber and is afraid
of being shot need then only close his eyes in order to annihilate the
existence of both the robber and the pistol. Childish. We thus arrive at the
result that the things which are sensed must be both within and outside our
consciousness.

In order to reconcile the two seemingly contradictory determinations, it is
necessary to investigate wherein sensation properly consists. The sensing
consciousness is at once both active and passive: active, in so far as it must
produce the object; passive, in so far as in each and all it must direct its
production according to the object. In consequence of the immediacy of
sensation, the passive has such preponderance that the sensing subject does
not even notice its own contribution, but finds itself wholly lost in the
objectivity of the sensed object. The assumption of an outside rests upon the
intuition of space, upon the irresistible necessity with which the object
obtrudes itself and the striking immediacy with which it shows itself. The
subjectivity can turn away from the object, can close its eyes and stop up its
ears: it does not help; when he looks about him, there it is again.
Different kinds of senses --- sight, feeling, hearing, indeed in certain cases
even smell and taste --- can unite into a concordant testimony of the reality
of the object. Knowledge of the reality of the object is a real knowledge, a
knowledge of outward experience.

Outward experience presupposes, to be sure, not only sensation but also
thought, distinguishing, comparing, judging, and so forth; but sensation and
sensory observation are nonetheless the main thing, since the thing must be
taken exactly as it is given. Respect for the given is respect for facts. It
is not the object that conforms to consciousness; it is, conversely,
consciousness that must conform to the object. From a strictly empirical
standpoint it is simply meaningless to ask whether things are also such as we
apprehend them. How could they be otherwise, when we consider what sensation
and experience properly are? It is not the sense-forms of consciousness that
originally determine what is sensed and decide how the thing is to look; it
is, on the contrary, the sensible thing itself that determines the
sense-forms; if it is not capable of this, then it is not to be sensed --- that
is, it is not a sensible thing. The object which by its natural constitution
cannot be seen does not belong to the visible objects; that which cannot be
heard, not to the audible. It is precisely the sensible determinations of the
object that make it sensible. But what is sensible? Surely that which allows
itself to be sensed. And what is sensible for us except that which corresponds
exactly to our sense-forms and can determine these according to its
constitution? Every thing must be determined by its properties. To assume
things that fall outside all properties is to assume things that fall outside
all determinations; to assume sensible things that have no sensible properties
is to assume things whose sensibility consists in being non-sensible.
Meaningless!

But, one says, the things of experience may perhaps well have sensible
properties and so far be sensible in and for themselves, without our being
able with our senses to apprehend them such as they really are. What we sense
as extended is perhaps not extended at all, or at any rate not extended in the
form in which we sense it; what we take to be hard may perhaps be soft; what
we regard as opaque may perhaps be transparent. --- We answer: quite possibly,
but what concern is that of ours?

The objects of outward experience we must necessarily apprehend such as we
experience them; to draw the validity of experience into doubt is to draw all
facts into doubt and to renounce real knowledge. The more sense for reality
the knower has, the greater therefore will also be his interest in a sharp
observation of things in all their particulars. What is otherwise generally
held to be reasonable or unreasonable does not trouble him; if the thing
really exists, then it must be taken as experience gives it, however
unreasonable its existence may otherwise appear to us. The foundation of
knowledge is, as is seen from this, a foundation of facts; an inviolable
condition for apprehending and securing the outward facts is to get them
fully sensed: this is the share of sensation in knowledge.

That one nonetheless cannot derive all knowledge from sensation and sensible
facts is evident. To sensation there belongs consciousness, but consciousness
is no thing of sense, no object of outward experience. The more one-sidedly
one accustoms oneself to apply the measure of outward experience to all
knowledge, the greater inclination one will also acquire to overlook the
non-sensible essence of consciousness, and to derive sensation itself from
empirical properties of the sense-organs. Instead of explaining sensation
from consciousness, one now sets out to derive it from the activity of the
sense-organs, and this activity again from the peculiar structure that can be
an object of outward experience. So one speaks of ganglion-consciousness, and
makes consciousness a physical property of the ganglion; of spinal-cord
consciousness, of brain-consciousness --- not merely in the sense that
ganglia, spinal cord, and brain are necessary conditions for the sensing
consciousness, but expressly in the meaning that these organs are themselves
sensing subjects. The absurdity of this strikes the eye when one comes to give
an account of how the many different consciousnesses can flow together into
one single unity of consciousness. By placing the connecting element in
connective tissue and nerve-threads, one achieves only --- what is, moreover,
very meritorious --- the demonstration of conditioning intermediate links; but
the matter itself one does not get cleared up. Can consciousness run through a
nerve like electricity through a wire? As the electricity streams through the
wire, the wire becomes electric; is it now admissible to say that a nerve
itself becomes conscious as consciousness streams through the nerve? It does
not help here to break off and declare that we at bottom know just as little
what consciousness is as what electricity is. Set the boundary for human
knowledge where one will, there is yet one thing that science with rigor
demands of the knower, namely this: that he must know what it is he himself is
saying. If we do not know what is meant by the word consciousness, how can one
then assert that there is consciousness in ganglion, spinal cord, and brain,
indeed, over and above this, make consciousness a property of these organs?
The fact is that one stares oneself blind on sensible facts and would derive
the non-sensible from the sensible; but the more one tries to make everything
distinct down to the smallest particulars, the darker, the more enigmatic the
whole becomes. It is the one-sidedness of sensualism that it drowns
consciousness in the unconscious and makes the sensible the original. What in
the eyes of sensualism is to have real worth must be a thing. ``The self, the
I, is not a thing; therefore the self, the I, is properly no thing at all.''

A proposition such as that we know just as little what consciousness is as
what electricity is shows itself, when we remove the ambiguity, to be
fundamentally false. So long as the inner, molecular structure of bodies, the
transposition, deposition, and reciprocal action of the minute parts is not
yet sufficiently investigated, it can rightly be said that electricity has
something enigmatic about it, something we cannot see through; but can the
same be said of consciousness? By no means! We have seen in the foregoing
section that consciousness precisely yields the original presupposition for
all our knowledge and belongs to what is first evident of itself. How
precarious it is to forget or even to overlook this makes itself known even in
the domain of empirical science itself.

It goes without saying that the theory of experience can under no form dispense
with thought; but it is precisely characteristic of the sensualistic theory of
experience that it seeks its proper ground of assurance in sensation, places
reality itself in the sensible, and lets thought be a mere matter of form. When
the sensualist speaks of the thing and its properties, the whole and its
parts, it is by no means thought that assures him that where there are
properties there must also be a thing, that where there is a whole there must
also be parts; no, it is sensation. He supposes himself simply able to sense
things as things, properties as properties, the whole for itself and the
parts for themselves, and likewise by sensation to assure himself of the
connection of the properties with the thing, of the cohesion of the parts in
the whole. But when thought begins to stir freely, then uncertainty arises.
The sensualist becomes irresolute; what is he to do? He believed so surely
that he had sensible guarantees that gold exists as a thing with its
properties, a whole with its parts, and that his knowledge that gold is a
noble metal rested securely upon outward experience. He knows that gold is
yellow, that it can crystallize in octahedra and let itself be rolled out into
leaves of a thickness of $\frac{1}{20000}$ of a line. Sensible observations
have convinced him that gold, of all metals, has the least affinity for
oxygen, that it allows itself to be dissolved in aqua regia, and so forth. But
what certainty has he now properly of all this? He wanted to sense each
property for itself; but the single property does not allow itself to be
sensed immediately, and that for the simple reason that it cannot exist for
itself immediately. When he sensed the property yellow, he had to refer it
partly to the gold, partly to the light, partly to his own visual sensation.
Now is there any immediate sensation that can teach him what share the gold,
the light, and the sensation of light each have in this yellow? He has without
further ado put together different properties into a real unity in thing and
substance, but immediate sense can teach him nothing of the reality of such a
unity. To be yellow, to be ductile, to crystallize in octahedra point in
reality to such disparate things that no sensation can assure us that they
must absolutely belong to one and the same thing. On the contrary! Outward
experience teaches that the same, or at any rate similar, properties can be
found in different things.

For the sake of experience the sensualist is compelled to assume a cohesion of
parts and a certain order in which things must accompany and succeed one
another; but what certainty has he of this? He is obliged to make use of words
such as necessary and contingent, ground and consequent, cause and effect, the
individual and the universal, laws, and so forth. But he can have no proper
confidence in such determinations; he cannot attach a real meaning to the
words. Immediate sensation apprehends the immediately given in its singularity
and contingency, and contains no guarantee either that it will continue to
subsist or that, as a link in a series, it is capable of exercising any
influence on the other links. By subordinating simultaneous phenomena to a
necessary connection, by assuming an antecedent as cause and a consequent as
effect, one goes beyond what can immediately be sensed; and since sensation is
nevertheless supposed to be the only thing that can give sufficient certainty,
one must either reduce the scientific yield of outward experience to a loose
manifold of phenomena, or slip in presuppositions which one yet cannot
acknowledge. By wishing to hold exclusively to a sensible certainty, one is
driven over into uncertainty, and skepticism stands at the door.

The adherents of one-sided empiricism only make the evil worse when, by
referring us to habit, repetition, the association of ideas, and such
palliatives, they try to veil the uncertain by means of the half-certain, and
give it the appearance that it is sensible knowledge which, by introducing
more and more intermediate determinations, little by little works its way
forward and makes the thought-determinations universally valid --- whereas the
truth is that it must, point by point, presuppose something other than
sensation. Sensation is a moment, an indispensable moment in all our knowledge
of sensible things; but when this moment of knowledge is to be extended so as
to hold for everything and to constitute the whole of knowledge, then one
smuggles in the original --- the knowing consciousness, the knowing self ---
while yet denying the original; and so the unscientific arises.

\section{The Share of Imagination in Knowledge: Inward Experience}
\label{sec:3}

\noindent According to ordinary usage one would suppose that imagination and
knowledge had absolutely nothing to do with each other. ``An imagining is
something that has no reality, and an imagined knowledge the same as a false
knowledge.'' The matter will, however, show itself in another light when we
reflect that imagination essentially means an ideal yet outward-directed
formation, a formation produced within us, an image-formation that makes the
active side of consciousness recognizable and corresponds to what we call the
power of imagination. The activity of the power of imagination consists
precisely in forming imaginings; how these imaginings stand related to real
being is to be specially investigated; here we use imagining and product of
the power of imagination as synonymous expressions.

That there must be a close union between sensation and imagination is easy to
see. When, after having sensed an object, we turn away from it, its image
still remains standing before the eye of consciousness, though fainter and
paler. This image bears witness that the power of imagination was already
active in the immediate sensation of the object itself. If the sensation of
the real object did not take place by means of an image-forming activity,
remembering and recollection would be impossible. During sensation itself the
image-formation is indeed so dependent upon impressions from the object that
the imagining is not noticed at all; but if the present object could be
apprehended as it shows itself to consciousness without an unconscious
image-formation, then its after-image could not form either; consciousness
would retain nothing: the object would exist for it only in the moment of
sensation. And what would the consequence of this be? That a consciousness of
several different objects would be impossible. Indeed, sensation itself would
be impossible. The sensible object must, as a whole, consist of parts. In
order to sense a whole whose parts are $a$, $b$, $c$, and so forth, we must
begin by sensing one part, for example $a$; if now, as we sense $b$, we have
lost $a$, and, when we come to $c$, both $a$ and $b$, then we never reach the
point of getting the parts comprehended together. But when we cannot
comprehend the parts in a unity of sensation, how then should we be able to
sense the whole?

A striking proof of the indissoluble unity of sensation and imagination lies
in intuition. It stands with intuition just as with space: we must explain
intuition by intuition, just as space by space --- that is to say, intuition
belongs to what is first evident of itself. Nonetheless it is the sensation of
real objects that gives us conditions for a clarification of what intuition
properly determines, and of what it is that determines intuition. When it is
said that, in order to sense a whole, we must begin piecemeal by sensing its
parts, we can conversely say that, in order to sense a part, we must begin by
sensing a whole --- that is, that for the sake of sensation we must begin with
intuition. Each smallest sensible part has an extension, consists of parts,
and is thus a whole: we cannot sense differentials. It is the immediate union
of intuiting and sensing that brings it about that we seem at once to sense
both the whole and its parts.

Without imagination it is impossible to sense any change in the outward. In
order to come to a consciousness that a thing changes, we must at the least
sense it twice, and in sensing it the second time be able to remember how it
looked the first time. Comparison is conditioned by sensation's being
converted into imagination, and imagination into sensation. In the manner of
sensation the change is apprehended by our comparing thing with thing, whereby
it comes to look as though the things themselves showed now thus, now
otherwise --- in other words: we find the things changed because they really
change. Time, the form in which the change takes place, then holds as well for
the things as for our apprehension of the things. But only by means of
imagination can the change be apprehended in such a way that not merely are the
objects compared with their images, but the different object-images are
compared with one another. Under this comparison the liberated images are
again changed into representations; the manifold of representations is the
object of inward experience.

The question whence such categories as inner, possibility, cause, and so forth
have their origin can in this connection be provisionally answered. When the
inner of things is spoken of, it is evident that this inner cannot be referred
to any outward experience. Take the immediate outer away, and there is again
for sensation a new outer. Whence then do we get the inner? From imagination.
We invent it into the thing, and come by opposition to the determination of an
outer. The same holds of possibility. What is to affect the outward sense must
really exist; the possible can therefore not come from without, but must come
from within --- that is, from imagination. Now if inner and possibility are an
addition of imagination to what must by the senses be determined as outer and
actuality, then the point is to discern the connection between sensation and
imagination. Without imagination, no sensation. Every experience of the
outward is indeed conditioned by sheer sense-impressions; but no
sense-impression can in pure, immediate singularity be apprehended for itself
alone; the manifold single impressions, be they simultaneous or successive,
must be gathered and determined by a corresponding manifold of feelings and
representations that spring from the unity of imagination in the self. If we
denote the sense-impression by $i$, the representation by $f$, then the series
$i$, $i$, $i$ gives us no experience, and the series $f$, $f$, $f$ none
either. The series $i$, $f$, $i$, $f$, on the other hand, makes an experience
possible; for now $i$ can be determined by $f$, and conversely, a comparison
can take place, and the differences show themselves: $i'$, $f'$, $i''$,
$f''$. The comprehension rests upon the unity of the intuiting,
representing consciousness.

The sensation of the one continually awakens a representation of the other.
The sensation of A awakens the representation of B, and this again prepares an
apprehension of C. The sight of a leaf awakens the representation of a stalk.
In the representation there lies an expectation, partly of something that
already is, though it has not yet been sensed --- as when the sight of the
front awakens the representation of the back; partly of something that is not
yet, but will come --- as when the sight of a flash of lightning awakens the
representation of a clap of thunder. The expectation is involuntary, and the
imagination can deceive. That something must follow is certain, but whether
precisely the expected follows is uncertain; something else may follow. Now
possibility has come to be in imagination. So long as the vicissitudes of
things are merely seen to take place in space and time, the representation of
an inner is not yet given. The links A, B, C can change their position to one
another without any of them changing; but if during the transposition a change
takes place either with one of the links or with all of them, then imagination
receives a new task. B is changed when A comes to it; B and C both become
changed when they meet: here an inner must be presupposed. Before and after the
change B is the same, and yet not the same. Were it not the same, then one
could not say that the change has taken place with B, that it is precisely B
that has been changed. Were B after the change still quite the same, then it
would have undergone no change. The objectionable element in this is removed by
imagination through referring the different outer to a common inner. Before, B
had one appearance, now it has another, but at bottom --- that is, according to
its inner --- it is yet the same. The one who senses becomes conscious that B,
according to circumstances, can show itself with a different outer in different
states: water can become ice, and ice water; that must lie in the inner of the
water, of the ice. What the inner is is provisionally determined by a
possibility, but what sort of possibility it is must be decided by the
ensuing actuality. The sight of a seed awakens the representation of a plant;
the seed is the plant's possibility, but only a possibility, for it is also
possible that a bird may eat the seed; the sight of a seed can also awaken the
representation of a bird. The representation of possibility is so far
subjective, but by laying the possibility into the seed imagination has
renounced its right over it and made the expectation dependent upon what is
sensed; imagination has, in other words, made the possibility objective ---
that is, laid it into the actual itself.

The union of actual and possible confirmed by experience --- a union that is
itself conditioning for experience's own possibility --- leads to the imagining
of cause. That A is cause and B effect can by no means be sensed immediately.
By making sensation the be-all and end-all of experience one can easily defend
the assertion that the representation of cause is an imagining to which nothing
actual corresponds. We cannot see in the fire that it will burn us, and yet
there is no one, not even the skeptic, who has any desire to stick his finger
into the flame. Why not? The sight of fire awakens the representation of pain.
Every time we try to hold our finger in the flame, we experience that it
smarts; why then, despite continual repetitions, do we not dare to regard it
as certain that the fire is cause and the pain effect? Because, it is said, one
cannot see the cause sitting in the fire. One can indeed feel the pain, but in
order to feel it as effect one would have to feel the fire as cause, and that
is just what one cannot do. The representation of cause is an imagining. This
mode of view is, in the real sense, irrefutable. One can just as little
maintain the reality of the cause by mere sensation as by mere thought.
Thought can establish that there must be a cause, as surely as there is an
effect; but if now there is no effect, then there is no cause either. The
scientific does not rest upon excluding imagination, but rests upon exact
determinations of the share of inner experience in the outer --- that is, upon
exact determinations of the share of imagination in knowledge.

That the power of imagination, even when it leans upon outward experiences,
is not able to ground any truly true or self-valid knowledge, lies in the
nature of the case and has shown itself in the history of culture. A people's
culture begins, as indicated above, not with science and scientific
investigations, but with myths and mythology. In mythology the power of
imagination has a free playground without yet grasping its own freedom: it
invents the gods and yet imagines that it is outward, sensible figures that it
is describing; it invents the supernatural and is not aware that it is sheer
natural constituents it has employed. Little by little there is indeed
awakened an inkling that there is a difference between poetic invention and
actuality, between the world of imagination and the palpable world of things
that is the object of outward sensation. But the inner experience of the
states of the mind, of want and longing, of hope and fear, is so
overwhelmingly strong that the outward actual, colored by imagination, is
everywhere transformed into a means for a conscious inner. Even when the
consciousness of antiquity had come to reflection and become attentive to
something that is otherwise more trustworthy than either sensible things or
imaginings and changing representations --- namely thought, which sees through
the phenomenon, holds fast the essence, and grasps the universally valid ---
even then one was not able to secure for oneself a knowledge of the real
content of existence. Either one broke the staff over the whole outward,
sensible world, denied all manifoldness, becoming, and motion, and held, in
order to avoid imaginings, so one-sidedly to the thinkable, with exclusion of
everything represented, that the pure thought-determinations, the empty
universal forms, at last swallowed each other up in contradiction, and the one
being that remained was transformed into an absolutely determined that was
devoid of all determinations. Or, when one saw that the sensible, the
changeable, and the perishable must yet have a kind of being, one let this
being of appearances dissolve into concepts, made the concepts self-subsistent
as ideas existing in themselves, and conceived the world of things as a
shadowing-forth of ideas. So a science formed itself, but the science became
speculative, a knowledge of essential imaginings. However much thought
exerted itself to maintain the eternally valid, one could not reach reality.
Reality was misjudged because sensation and the significance of the sensible
were misjudged. --- Speculation is idealism, a transformation of what is
foreign to consciousness into likeness with the proper essence of
consciousness, a knowledge in which the object and the knower are reflections
of one and the same thing, a knowledge in imagination.

And yet imagination is just as necessary for science as sensation, speculation
just as necessary as empiricism, ideality just as necessary as reality. The
assumption that by exclusively looking at facts and holding to the real one
can free oneself from imagination is itself an imagining. One cannot, for
example, dispense with language. Language is of idealistic origin; an imagining
underlies every word. The word is to correspond to the object, and yet there
is no likeness between the word ``tree'' and the real tree, the word ``house''
and the real house. The word is a sign; the bond between the sign and the
signified is an imagining; but this imagining is utterly indispensable. What is
nothing? A word! The word ``nothing'' and nothing itself have no likeness with
each other --- that is, the likeness is an imagining. And yet this imagining is
utterly indispensable for thought. If thought were not, by the help of
imagination, able to grasp nothing, then it would also not be able to grasp
negation as an essential boundary-determination, not able to maintain the
difference between possible and impossible, not able to apprehend laws.

Imagination is a false leader, an unreliable guide; but precisely in its
unreliability must we seek its advantage, its indispensability. For the
unreliability of imagination is one with its flexibility, its compliance, its
infinite suppleness. It does with consciousness all that it will, for it is
sheer possibility; it shows consciousness out beyond itself and into itself,
makes it a slave of all things and a lord over all things; it plays with
difference and laughs at contradiction, sets boundaries everywhere and effaces
all boundaries, makes the sensible thinkable, the thinkable sensible; it is,
in other words, the infinite freedom of the knowing consciousness.

Formal freedom is an inviolable condition for cognizing necessity; knowledge
must, for the sake of necessity, be free. For scientific insight it is not
enough that the knower finds before him a world of things outside his own
consciousness; he must know what this outside properly means; but in order to
come to know this, he must, by the help of imagination, be able to climb over
the wall of consciousness, must be able to represent to himself how it would
look were he to stand on the other side. When the unreasonable, the confused,
then shows itself to him, he returns to the immediacy of sensation, lays due
weight upon the testimony of the senses, and acknowledges the reality of
things. It is, in the scientific sense, not sufficient that the knower
observes things as they in fact accompany and succeed one another, and notes
the difference between usual and unusual cases; he must on the wings of
imagination raise himself high above the cases of experience, and only when he
has reflected that the whole would be completely meaningless if in all these
outward changes there were nothing unchangeable, a law according to which the
one proceeds from the other and follows upon the other --- then only does he
learn to bow to the necessity of thought and to grant its determinations of
laws, grounds, and causes an original validity.

\section{A Priori Knowledge}
\label{sec:4}

\noindent The difference between \textit{a~priori} and \textit{a~posteriori}
knowledge seems in the simplest way to be designatable thus: all the
knowledge we draw from our own consciousness, and have prior to a particular
knowledge of the thing, is \textit{a~priori}; whereas the knowledge that
presupposes a sensation of existing objects is \textit{a~posteriori}.
Fundamental representations such as unity and plurality, likeness and
unlikeness, magnitude and smallness, beautiful and unbeautiful, good and
evil, and others, are \textit{a~priori}. When we say of one action that it is
good, of another that it is evil, we do not draw our representations of good
and evil from the actions, but we refer the actions to fundamental
representations that we carry in ourselves. When we call one object beautiful,
another ugly, we do not borrow the measure for beautiful and ugly from the
objects, but from exemplars in our own consciousness. To apprehend two objects
partly as like, partly as unlike, presupposes a comparison; but a comparison
would be impossible if we did not have the fundamental representations likeness
and unlikeness beforehand --- that is, \textit{a~priori}. If, on the other
hand, the talk is of such objects as stones, strata, plants, animals, and the
like, then the relation reverses itself: then the object must precede our
representation, our concept of it. The representations and concepts that are
thus drawn from things and therefore come after the things, we call
\textit{a~posteriori}, empirical.

But when the boundary between \textit{a~priori} and empirical is to be
apprehended definitely and exactly, difficulties arise. To lay the emphasis on
the \textit{a~priori} is characteristic of idealism; to make the empirical the
be-all and end-all is characteristic of sensualism. By subjecting the
idealist's ``innate ideas'' to a sharp critique, the sensualist can easily
persuade himself that the \textit{a~priori} is only a semblance, that all our
knowledge allows itself to be derived from sense-impressions, sensations,
reflections of the sensible --- that is, in an empirical manner.

``Were our knowledge,'' says the sensualist, ``really to stem from innate
ideas, then these ideas would have to be not merely in the soul as soul, but,
for the sake of knowledge, expressly in consciousness. The human soul that has
reached the point of having become conscious of itself must then at the same
time have become conscious of these ideas; for that unconscious ideas are in
consciousness is a contradiction. But now our inner experience teaches us that
the so-called first ideas have formed themselves little by little, that a
multitude of sense-impressions, feelings, and single representations have
preceded before we could form ideas which, though they seem single, are yet at
bottom highly composite. How many like and unlike impressions, feelings of
pleasant and unpleasant, of white and black, of hard and soft, must there not
be in a consciousness before it can reach the common names: likeness and
unlikeness, hardness and softness. Common names are abstractions; in order for
consciousness to abstract, it must have something from which it can abstract.
It can therefore not be abstraction, but sensation, that is the first. The
secret of the innate ideas consists simply in this, that the sensing soul,
receptive to sense-impressions, abstracts the common or universal from the
individual, makes the abstracted self-subsistent, and imagines that in the
derived it possesses the original. --- Outward experiences teach us the same.
Were there innate ideas, then we should find them in all who have any
consciousness; but that children and savages have no inkling of such ideas is
notorious. An Australian aborigine, for example, may have a sharp eye for like
and unlike things in his surroundings, but for likeness and unlikeness
\textit{in abstracto} he has no perception. Tell him that two magnitudes, each
of which is equal to one and the same third, are equal to one another: he will
not understand it. If we choose ideas such as beautiful and ugly, good and
evil, of which most people do have a representation, then it becomes quite
evident to us that these ideas cannot be innate. Were they innate, then they
would have to be the same in all, have the same validity for all; there could
then arise no dispute about what is in and for itself beautiful and what is
ugly, what is in and for itself good and evil. But that on this point great
disagreement prevails, not only among the ignorant and half-knowing, but
especially among those who ought best to know the ideas, is an indisputable
fact.''

If sensualism has in this way robbed the soul of all innate ideas and swept
the \textit{a~priori} out of consciousness, the investigation must begin
afresh; and what then becomes the original? Experience! But in order to
experience, one must sense. Sensation! But sensation presupposes a soul, a
consciousness, that can sense; sensation is, as shown above, inseparable from
imagination, and the discerning imagination again inseparable from thought.
But when all experience must necessarily presuppose a sensing, representing,
thinking consciousness, then experience cannot possibly itself be the first.
We get the \textit{a~priori} back; yet not in the form of innate ideas, but
as capacity, as endowment. Everything now depends upon ascertaining wherein
the capacity consists, and upon showing the different modes of consciousness
in which the endowment makes itself known.

The half-idealism that is called the critical lays particular stress upon a
sharp distinction between the forms that necessarily belong to consciousness
and the content that is due to the thing; experience comes about by a content
from the thing filling the forms of consciousness. The existentially concrete
is always empirical, but the forms are abstract; the point is to become
clearly conscious of the difference between the concrete and the abstract, and
moreover to establish which and how many forms of consciousness are
incontrovertibly \textit{a~priori}.

The following procedure may pass for characteristic. ``We apprehend space
concretely, in that we see it filled by a countless multitude of objects of
experience. But the concrete apprehension is ambiguous; we must ask: is space,
as form for the existence of things, given to us in and with the things, or is
it a form of consciousness that lies in our apprehension and does not concern
the thing in itself at all? In order to get this question answered, we must
apply abstraction. The interval between A and B is filled, for example, by the
objects a, b, and c; if we now think of these objects as gone, then the space
between A and B would become empty; but it will not become gone, no part will
be lost: the distance from A to B will be the same as before. If space were a
determination belonging to the object as object, then each object would have
to take its space with it when it became gone; there would have to be, so to
speak, holes in space, a hole for each of the vanished objects. Since this
cannot possibly be the case, inasmuch as holes in space are themselves spaces
with imagined boundaries, it follows from this that all things in space can
vanish without space vanishing: and from this it again follows that it is not
the things themselves, but consciousness, that bears the form of space in
itself. This can be expressed thus: we can abstract from the things in space,
but from space as space we cannot abstract; when all things have vanished,
consciousness retains its own \textit{a~priori} sense-form, the boundless,
infinite space. --- We apprehend time concretely when we regard it in unity
with its content --- that is, with the changes and events that take place in
it; in accordance with this concrete apprehension it is said: time hastens,
the times change, and so forth. But if we abstract from the content, then time
stands still. All that has shown itself in experience as an antecedent, a
present, and a consequent can be thought away; but as regards the possibility
of experience, time as time cannot be thought away. That in the pure time
separated from all real or imagined changes there can be no changes goes
without saying; but that time can also not lie in the things or adhere to the
things is likewise evident. We divide time into present, past, and future;
were time now to belong to what is in itself, then it would have to be an
infinitely continued present; for the past has vanished and corresponds to a
non-being; the future, which has not yet come, likewise points to a non-being.
But how long is the present? It has, when we look more closely, no extension;
it dwindles to a now, the most fleeting of all, and cannot possibly concern
what is in itself. This ambiguous fact, that time is at once standing still and
vanishing, proves that time as time is not reality, but a form of imagination,
a form necessary for all changes, an \textit{a~priori} form of consciousness.
We apprehend the understanding concretely, in that we apprehend the things we
understand in immediate union with our understanding of them. It is a
proposition of experience that our understanding must conform to the
understanding expressed in the things: we gain understanding of things because
there is understanding in things. Are there then two kinds of understanding,
an understanding in us that can err, and an understanding in the things
themselves that cannot err? Can the things through experience really be
assumed to bring the infallible understanding into our consciousness?

In order to get these questions answered, we must not merely look away from
the manifold content of the understanding and by abstraction emphasize the
pure forms of understanding; the abstraction must even redouble itself into
reflection, for the forms of understanding are forms of reflection: the
understanding works in reflecting and thinking. There is understanding in the
thing in so far as the thing has properties and by these properties stands in
relation to other things; the properties are abstract determinations and
become intelligible only by the fact that the one abstraction is mirrored ---
that is, reflected --- in the other. Light, for example, has its properties,
and the bodies that are illuminated have, each for itself, their properties.
But in the abstract determinations for themselves there is no understanding.
The intelligible shows itself only when the properties are reflected in one
another reciprocally, those of light in those of the bodies and those of the
bodies in those of light. The leaf is green; that is a proposition of
experience. The intelligible element in this proposition of experience appears
from the fact that the proposition can be transformed into a judgment in which
the thing is subject and the property predicate. The same holds of the
proposition: the cloud is grey, and countless others.

If we now abstract from the different content of the propositions of
experience, then we are convinced that the intelligible lies neither in the
things nor in the properties, neither in the particular subjects, which can be
anything whatever, nor in the predicates, of which the same holds, but solely
and only in the reflected relation of thing and property, subject and
predicate: all determinations of understanding are determinations of relation.
Every understanding of particular relations must necessarily presuppose an
original perception of relation. It is hence clear that things in their
manifold relations cannot possibly teach us what relation is, or impart to us
any understanding other than that which we originally have in ourselves.
Instead of saying that it is the understanding in things that is by experience
carried over into our consciousness, we must conversely say that it is our
consciousness that lays the understanding into things. The universally valid
determinations of reflection and categories of the understanding therefore
belong exclusively to the \textit{a~priori}.''

The boundary-line which critical half-idealism in this way sets up between
\textit{a~priori} and empirical knowledge, between what comes from
consciousness, from within, and what comes from things, from without, has
nevertheless the precarious feature that the distinction is only critical.
The critical distinction is always recognizable by a fixed, immovable
boundary; but a fixed, immovable boundary-line between the \textit{a~priori}
and the \textit{a~posteriori} cannot be drawn. This appears already from what
has been said (\S\,1) about knowledge and object. If one runs knowledge
together with the knowing individual, then it is true enough that the
individual with its consciousness has the objects outside it --- the individual
is itself an outward object; but who is able to demonstrate a fixed boundary
between knowledge and the object of knowledge? Between countries, provinces,
plots of land, and the like, definite boundaries can be drawn. These
boundaries may at a certain time be fixed and then again be moved; but it is
not so with the boundary between what lies within and what lies outside the
original property of consciousness. It is not true that by abstracting from the
content of space and time one can retain pure space and pure time; for when the
last shadows of things and changes vanish, space and time vanish too. It is not
true that by abstracting from all particular relations and terms of relation
one can become conscious of the general determinations of relation as the
\textit{a~priori} forms of the pure understanding. Relations without terms of
relation are unthinkable; even in $\frac{0}{0}$ there must be a reflection of
vanished terms if there is to be any relation. These reflections, these
reflexes necessary for reflection, are in the last instance reflexes of
existing things. Imagination may have transformed them; but imagination can
take nothing of itself, it must draw its material from immediate experience;
abstraction can thin the reflexes to transparency, but abstraction must have
something from which it can abstract.

The outward must reflect itself in an inward, and conversely; the objects must
reflect themselves in consciousness, and consciousness by the help of the
objects; the empirical must reflect itself in the \textit{a~priori} and this
in the former. Immediate boundaries of existence, boundaries of existence, are
boundaries of finitude, and such boundaries signify that A falls on the one
side, B on the other side of the boundary; but the boundaries of reflection
have the peculiar feature that A and B exclude each other only by
interpenetrating one another in an infinite reciprocal play. The soul does not
first, even in an unconscious manner, have its \textit{a~priori} sense-forms
and forms of understanding, so that impressions from things can come to it and
experience begin. Here there is in time no first and no last. The
\textit{a~priori} forms must just as originally rest upon impressions from
things as these impressions upon the forms. The \textit{a~priori} is not
confined to a certain number of abstract determinations of form which, once
brought to consciousness, prove to be finished. No; from the fundamental
relation it follows that there must be just as many modifications and
ramifications of the \textit{a~priori} determinations of form as there are in
the world of experience modifications and ramifications of corresponding
determinations of objects. The growing manifold of cognized object-determinations
gradually calls forth a corresponding manifold of \textit{a~priori}
determinations of reflection. All knowledge and science aims at an
\textit{a~priori}-zation of what is given from without; for knowledge and
science strive for insight, and all insight consists in an
\textit{a~priori}-zation. The false \textit{a~priori}-zation, recognizable by
the fact that one partly overlooks, partly distorts what is given in
existence, forgets the decisive significance of the empirical
object-determinations for the determining reflection, misjudges the
\textit{a~priori}, and must therefore make do with an imagining instead of true
insight.

What physiology shows --- that in the course of millennia a progressive
development of the brain and of cerebral activity can be traced --- finds its
confirmation in the history of the sciences. Consciousness of the
\textit{a~priori} must be acquired just as well as consciousness of the
empirical. A one-sided cultivation of \textit{a~priori} sciences can lead to a
misunderstanding of the \textit{a~priori}, and a one-sided cultivation of
empirical sciences to a misunderstanding of the empirical; for development
rests upon reciprocal action. True insight into the empirical is conditioned
by a corresponding insight into the \textit{a~priori}, and conversely.

True knowledge is insight; progressive insight rests upon a steadily
progressive \textit{a~priori}-zation: knowledge is in its essence
\textit{a~priori}.

\chapter{Subjective and Objective Knowledge}

\section{The Problem of Objectification}
\label{sec:5}

\noindent In so far as the content of knowledge lies in the knower's own
interior, knowledge is subjective; in so far as the content is drawn from the
world of things, knowledge is objective. Yet on this presupposition no
critical --- that is, fixed, abiding --- boundary can be set up between
subjective and objective knowledge. All subjective knowledge must necessarily
have an element of the objective, for a knowledge without object, an absolutely
objectless knowledge, is just as meaningless as a circle-center without a
periphery. Nonetheless there is an essential difference between outer and inner
objects, and, as regards the latter, an essential difference between the
subjective, which belongs exclusively to the single individuality, and what is
common to all knowing individualities, the ideally objective, which is
synonymous with the universally valid. Knowledge of a stone is outwardly
objective; knowledge of the law of falling bodies is essentially objective; but
the individual's knowledge of his own state, his desire and longing in any
given moment, is essentially subjective, although longing and desire are ideal
objects of knowledge.

The fundamental condition for all knowledge is that the knowing self must
apprehend and determine its object, whether this object obtrudes involuntarily
and with blind necessity from without, or is voluntarily brought forward from
within. Only when the object of knowledge is apprehended and determined by the
knower --- that is, objectified --- does it become an object of knowledge. No
object without a corresponding objectification; it is an \textit{a~priori} law
that bears all empiricism, a fundamental law that in science is, if possible,
even more unshakable than Newton's law of gravitation. From this it is seen
that a critical boundary, a boundary-line on whose one side one holds fast the
objectifying subjectivity for itself, while the object remains standing on the
other side, is confusing and meaningless. The opposition of subjective and
objective must be determined in such a way that in the object at every point
there is a moment of the subjective, and conversely: objectification itself is
the unity of the subjective and the objective, but a unity of opposite terms,
a unity of opposition.

Only when we apprehend the opposition of subjective and objective knowledge
\textit{in concreto} is there ground to inquire after a critical boundary. The
content of knowledge that is drawn exclusively from the peculiar interior of
the knowing individuality cannot possibly take on any objective, intrinsically
universally valid form or be scientifically grounded. What subjectivity within
this horizon objectifies concerns only its own specific being, its thoughts,
feelings, inner experiences. However unreservedly the subject may speak of it
to others, it does not thereby lose its subjective character. What is thus
objectified may belong to natural idiosyncrasies, defects, advantages,
exceptional peculiarities, and the like; and it may comprehend what is most
essential of all for subjectivity --- its personal view of life, its individual
apprehension of its original destination, its final goal, and the conditions
for attaining it.

The objectively knowing person must be aware that human subjectivity, judging
by objective experiences, can attribute to itself unequal worth, can conceive
itself as merely temporal, can posit itself as eternal, and on the whole
apprehend its destination in different ways; but he must also know that the
ground of the different self-valuation cannot be drawn from anything
intrinsically objective, but lies originally in subjectivity itself. If the
person in question is not yet clear about himself and oriented in his own
consciousness, he is continually liable to overstep the boundary, to chase
after objective grounds, to seek certainty about himself outside himself, and
to confuse subjective knowledge with objective.

Here several cases may show themselves. If the individuality is weak, and the
lack of clarity has its ground in weakness, then he gets no rest until he has
found an authority who can tell him what he properly is and to what he is
destined; this authority is then the objective: now he supposes himself to have
an objective knowledge of his highest destination. If the individuality is
strong, and the lack of clarity has its root in unconstrained egoism, then he
finds no peace until he has imposed his individual view upon other individuals,
the more individuals the better: the submission of the many is to him a proof
that he has succeeded in making his subjective knowledge objective. If the
individuality is profound --- that is, so utterly profound that what is in
principle disparate runs into one before his gaze --- then he will overstep the
boundary with great dignity. He will become objective to himself in the
following imagining: ``What is truth for me personally must necessarily be
objective truth, truth for all, for my essence and the essence of humanity are
at bottom the same essence. When nevertheless not all assume what I assume,
then it is either from lack of good will or from lack of insight into the
objective, for my subjective knowledge is objective: I personally know the
objective truth.'' --- If we regard these three individualities --- the weak,
the egoistic, and the strangely profound --- working together in history, then
we understand what strife, what hatred, what persecution and subjugation it
costs to get such a subjective knowledge thoroughly transposed into the
objective.

From an opposite extreme we find objective knowledge willing, in the name of
science, to take subjective knowledge under its guardianship. First, it is said,
science must have human nature thoroughly investigated before the individual
dares to pass any judgment on the proper worth of the personal; only when, by
physiological investigations, one has ascertained what that which we call soul
properly is, can the individual come to know whether he really has a soul or
not. That it is not science, but a misunderstanding of what science is, that
from this side tempts one to overstep the boundary, lies in the nature of the
case. The scientific consciousness is indeed turned toward the objective, even
when subjectivity is at issue, for only the universally valid has validity in
science; but it is therefore far from denying the subjective, in the
peculiarity of the self, its original right and significance. In order to get
the objective scientifically illuminated, the knower must have the sun of
objectivity at his back.

How the objective takes the power from the objectifying subjectivity, strips it
of all individual idiosyncrasies, and compels it to submit to universally valid
determinations and laws, is seen most clearly from such examples where the
objectified does not depend upon outward experience, but has its presuppositions
in consciousness itself. We presuppose that all men are square; we presuppose
that the individual inquired about, Caius, is a man. These are arbitrary
presuppositions, and we should do best not to assume them; but once we have
assumed them, we are bound by the assumption and cannot withdraw from the
consequence that Caius must be square. We set a proportion such as $a : b = c :
x$; we set it quite arbitrarily; but, bound to the presupposition, we are
obliged to assume $x = \frac{bc}{a}$. From subjective presuppositions follow
objective consequences, and from objective consequences objective knowledge.
This knowledge is indeed as yet only formal; but the formal objectification is
nonetheless intellectual, and intellectual objectification is an inviolable
condition for all scientific knowledge. Should anyone here remark that a
subjective content can also be objectified intellectually, that what is thus
objectified must take on scientific character, as surely as it is the
intellectual objectification that gives knowledge its objective stamp: then it
is not difficult to assure oneself that the boundary between subjective and
objective knowledge still stands firm.

Real knowledge, knowledge of a positive content, requires not merely a formal
but also a real objectification. A couple of examples may elucidate this. If
the human self is destined for immortality, then every man must prepare himself
to be judged according to the laws of eternity. --- If all bodies are heavy,
then every body which, raised above the surface of the earth, is deprived of
its support, must fall. --- In formal respect both propositions are equally
acceptable; but in relation to knowledge they are disparate. The premise in the
first is a proposition of subjectivity, whose general validity wavers because it
is apprehended differently by different individualities; its lack of objectivity
shows itself in a general uncertainty, a constant hovering between reality and
imagining. The latter premise, on the other hand, is a proposition of
objectivity. The understanding, which leans upon outward experience, can be made
evident to all who have ordinary sense and understanding. The inner peculiarity
of the subjectivities does not here come into consideration.

Since an objectification of outward objects must just as well presuppose the
contribution of subjectivity as an objectification of inward ones, it follows
undeniably from this that in the former, subjective differences must just as
well be able to show themselves as in the latter. The knowledgeable person
objectifies things to himself otherwise than does the unknowledgeable. The
astronomer objectifies the structure of the world to himself in a quite
different way than one who, without acquaintance with astronomy, immediately
contemplates the starry sky. The mineralogist sees in the mineral a different
object than the ignorant amateur who collects rare stones. The botanist's
objectification of a plant is essentially different from the painter's and the
poet's. We grant it. But these differences do not efface the boundary between
subjective and objective knowledge; on the contrary, they serve only to sharpen
it and make it recognizable.

Subjective knowledge points us to differences in subjectivity which can only
approximately be removed in a subjective manner, by discipline, awakening,
admonition, and so forth. Objective knowledge, on the other hand, takes account
only of such differences of subjectivity as it is itself able to remove or at
any rate equalize in an objective manner --- namely by making the
unknowledgeable knowledgeable, the ignorant knowing. The possibility of such an
equalization rests upon the fact that the consciousness of individuals, however
divergent the preliminary modes of view may be, contains fundamental
conditions, elementary presuppositions, that are common to all --- namely
sensation and thought. The ignorant person can see that the sun moves while the
earth stands still, and that is precisely a condition for his being able to come
to see that it is the earth that moves while the sun stands still. The
unknowledgeable person can see that the leaf is green, and this is just a
condition for his being able to learn to see that it is properly not the leaf,
but its chlorophyll, that is green.

Subjective knowledge of existing things gives us only an opinion; and the
opinions are different, indeed conflicting with one another, inasmuch as they
rest upon superficial objectification. Objective knowledge removes the
difference of opinions, in that it gives insight, and insight is conditioned by
essential, thorough objectification.

Imagination, however, has its share in objective knowledge as well as in
subjective, and the influence of this is in reality of so questionable a kind
that it threatens to deprive us of objectivity itself and to reduce all our
knowledge to a merely subjective knowledge. The object is, according to the
fundamental law, the objectified; so far it does indeed become meaningless to
wish to draw a critical boundary between the object itself and our apprehension
of it. ``Is, for example, this table, this lamp in itself such an object as the
one I can objectify?'' What an unreasonable question. What sort of object are
you speaking of? ``This table!'' Is it the table you objectify, or another that
you do not objectify? ``By the object I naturally understand this objectified
table.'' And then you can ask whether what you objectify is also really what you
objectify.

But objectification encloses two disparate moments, the subjective and the
objective, and when each of these is brought forward for itself, there can
indeed be a question of how far the object and the objectified are in each and
all congruent. We can test it by alternately setting the variable and the
constant in the two factors of objectification.

If we set the constant in subjectivity and let the variable fall in the
objective, then the different results of objectification will be dependent upon
the objective in the objects: the results vary because the objects vary. The
table shows itself square when it is square, and round when it is round.
Imagination makes this manifest by letting each object stand with such
definiteness and independence outside consciousness that no objectification at
all is needed: the stone is stone and the leaf leaf, even if no one senses it;
the sun shines, whether anyone sees it or not.

If we set the constant in objectivity and let the variable fall in the
subjective, then the different results of objectification will make every
objective determination unrecognizable. Imagination makes it intuitable by
letting the objectified have the object behind it in the infinite. Behind the
objectified lamp stands the proper lamp, the object as it is in itself; if
another subjectivity, a subjectivity with more perfect conditions, objectified
the lamp, then it would become another lamp. The objectified lamps would be
different; but the lamp in itself, the proper object, always stands behind; the
objectified lamp varies with every varying objectification, and the lamp in
itself, the constant object, recedes into the infinite.

The result of these investigations does not, however, lead to a breach of the
fundamental law, but to a closer understanding of the different turns under
which imagination can deceive the objectifying consciousness. As doubt arises
about the validity of objectification, imagination at once has two objects for
comparison, namely the objectified and the non-objectified. The objectified
lamp, the object that corresponds to the word, has the lamp in itself --- the
unnameable that does not correspond to the word --- behind it. How then do we
come to consciousness of the non-objectified object? Imagination has of course
itself objectified it, but this is supposed to be a great secret. If imagination
is to come out with the secret, it shows itself at once to be a deception. Here
one would have to distinguish between imaginary and real objects; but the
distinction cannot be carried through. First it is said that the objectified
objects, the table and the lamp, are imaginary, and that the unnameable objects,
the table in itself, the lamp in itself, to which the words apply only as
imaginary designations, are real. But when one then reflects that the unnameable
objects, to which imagination on its own responsibility attributes objectivity,
are precisely the imagined objects --- that is, the imaginary ones --- then we
are tempted to take the objectified for the real. Here is an uncertainty that
must be removed.

\section{The Problem of Cognition}
\label{sec:6}

\noindent The law of objectification says that only that is an object which
becomes objectified; but from this it by no means follows that our human
objectification is the only one that originally determines the objective, or
that the objects are in each and all such as we are able to objectify them. The
question how far our faculties reach when the talk is of a true, a thorough,
comprehensive, and exhaustive objectification of conditions of existence and
existing things is the fundamental question of cognition. In a general sense
the question can be answered thus: we cognize what we objectify, in proportion
as the objectified can become transparent to us; but it becomes transparent to
us only in proportion as our consciousness is able to produce it. Bounded forms
of intuition, objects of intellectual objectification, such as lines,
triangles, circles, tetrahedra, hexahedra, and so forth, we cognize fully, in
so far as we are able to produce them in the form of intuition and by
corresponding equations to subsume them under certain laws. To ask whether
behind the objectified circle there should not stand an unnameable object, a
circle in itself, which we cannot cognize, is idle profundity.

It is otherwise with the objects that have in themselves something we cannot
produce, an element that obtrudes upon intuition and determines the
representation, but remains standing for consciousness as something outward.
This is the case with matter and the material. We know matter, we acknowledge
its reality; but its essence, what it is in itself, we cannot cognize. Our
thought can, by the help of intuition, produce a circle-element, a
hexahedron-element, but a mass-element, a molecule, an atom we are not able to
produce. Therefore the objective in matter presents itself as something dark,
enigmatic, intrinsically incognizable.

The material in matter is sensible, but not thinkable, even though it is
determined by thought. If we investigate the significance of sensation for
objectification, and thereby for cognition, we are referred to the physiology
of the senses. We learn how the sense-organs are built; we learn how
manifold is the activity these organs exercise, how different they are both in
number and in form and development in the lower and higher animal species,
indeed that a single organ in man, the eye for example, can be so one-sidedly
endowed that it, as in the color-blind, makes an exception to the general rule.
Now as we come to see how dependent the sensible objectification is upon the
peculiar activity of the senses, how dependent this activity is upon the
structure of the sense-organs, and how dependent our cognition is upon the
testimony of the senses, it is no wonder that doubt arises whether we are
indeed able to cognize things as they are in themselves. To place the cause in
the thing, the effect in our cognizing, and to appeal to the likeness between
cause and effect, is extremely precarious. What likeness is there between an
ether-vibration and a sensation of light, between an air-wave and a sound?
What transformations must the influences coming from things not undergo before
they reach consciousness through the brain?

If we now assume that on other globes in the universe there may be beings with
senses of another and higher kind than ours, then these beings would
objectify things to themselves in a far more perfect manner than we, and in
the material would perhaps cognize an object of which we have no inkling. Nor
is this all. When we cannot even rely upon the objective validity of our
sense-forms --- for what assures us that space has precisely three and only
three dimensions? --- then we can also not fully rely upon our categories of
understanding. We are indeed, in all our cognizing, obliged to discern and
distinguish unity and plurality, premise and consequence, whole and parts,
cause and effect; but who can know whether such a distinction has validity for
the perfect understanding? It is possible that there is an intellectual
intuition which, by seeing infinitely many things in space at once, is exalted
above space, and by seeing countless motions in one view, is raised above
time; it is possible that an understanding with such an intuition does not
cognize the discerning, distinguishing, and dividing that are indispensable
for our cognition, but sees the unity in the plurality, the consequence in the
premise, the parts in the whole, the effect in the cause; thus: it is possible
that there is an understanding so qualitatively different from the human that
we cannot even form any representation, let alone any concept, of it.

But what then follows from this? That our so-called objective cognition is only
a cognition of phenomena and corresponding laws, but by no means any cognition
of the objective itself, of the essence of things. It is a vulgar opinion that
much knowledge and deep-going research should lead to such high notions of the
infallibility of the faculty of cognition that the researcher may easily be
tempted to arrogance. This is by no means the case. We find as a rule that the
most distinguished researchers, those who have enriched science with the
greatest discoveries, speak in a very modest language of the reach of our
cognition: ``what we know is as nothing reckoned against what we do not know.''

This admission is, however, misunderstood when the unscientific person, content
with a practical judgment, would instruct the researcher and lead him onto the
right path. ``Why make such a fuss of science when, by your own confession, it
is not able to lead us to cognition of the essence of things? You stand with
the riches of your investigations like a poor man on the barren heath of
objective cognitions; every time you suppose you are holding the object fast, it
transforms itself into a mirage, a phenomenon of something you do not know. You
measure the strength of an electric current, but do not know what electricity
is; you observe the structure of the body, but do not know what life is, or
whence it comes; you grasp after the fruits of cognition and would slake your
thirst at the springs of truth, but the fruits elude you, the springs deceive
you: you have it almost like a Tantalus in his underworld. No, then the
subjective is otherwise refreshing. Regard existence in the mirror of freedom;
use the light of life, let the miracle help you, and do not bind yourself by
petty calculation to blind laws.'' Such well-meant counsels cannot, however,
help the honest researcher. He knows that if he but lets go the thread of the
laws in a single point, all his research is in vain.

Scientific tasks require scientific solution; those who suppose they can help
science along by introducing cognitions of an overscientific or unscientific
kind only betray that they have never grasped a scientific problem and scarcely
suspect what science essentially is. Science must help itself.

The difficulties adduced are, then, not so daunting as a theory of cognition of
one-sided critique would make them out to be. It is indeed, as regards the
senses and the sense-organs, undeniable that one and the same substratum must,
by being objectified in different systems, show itself as different objects;
but this confirms precisely what the fundamental law says. If the substratum is
denoted by $X$, the determinations lying in it by $X_0$, $X_1$, $X_2$, then
these determinations, when sensed in the $A$-system, will give $a_0$, $a_1$,
$a_2$, \ldots\ and the whole corresponding object will be $O_a$; in the
$B$-system it will be $O_b$; in the $C$-system $O_c$, and so forth. But is this
a proof that the objective cannot be cognized? By no means; it is precisely a
proof of reciprocal action between objective and subjective. If the same
substratum, on the contrary, could, by being objectified in different systems,
give occasion to the formation of the same object, $O_a = O_b = O_c$, then only
would the cooperation of the objective be transformed into an imagining. That
different objects must form when the substratum is objectified in different
systems is just as intelligible as that different bodies must form when carbon
enters into combination with oxygen, with hydrogen, with nitrogen, and so
forth. A subjectivity that was able to see through the different activities of
the different sense-systems would by no means either declare certain objects,
$O_b$ and $O_c$, false and only $O_a$ true, or all the objects equally false;
but would much rather inculcate that each object in its system has full
validity on account of the reciprocal action. That the objectification in the
one system can stand lower than in the other makes a difference only in so far
as the effect must always correspond to the cause.

It has been maintained that cognition cannot possibly reach the objective,
since the cause in the object must be highly unlike the effect in the
subjective. This difficulty is more artificially contrived than real.
According to reciprocal action the cause cannot belong exclusively to the
objective and the effect to the subjective, but, when the conditions are taken
along, there is cause on both sides, in the subjective as well as in the
objective, and effect from both sides; there can then, according to the law of
reciprocal action, be no doubt that the effect corresponds to the cause. Only
abstraction in union with imagination can conjure up the semblance that there
are objective determinations for themselves without any regard to the
subjective, and subjective determinations without relation to the objective; but
one need only make an attempt to carry through such a distinction, and one will
discover in fact how empty it is. When every objective determination is a
determination of reciprocal action, and every subjective determination
likewise, how can one then doubt that the cognition resulting from the
reciprocal action must give a real unity, an essential likeness, of the
objective itself and the objectified.

Yes, one grants, a likeness there indeed is; but the likeness concerns not the
essence of the object, but only its phenomenon. --- Yet what does this properly
mean? That imagination again plays its game with abstraction! One surely does
not wish to set objects of things in a row with objects of hallucination? If the
phenomenon has nothing at all to do with the essence, then it is wholly
without significance; then the essence of the lamp may perhaps be hidden under
the phenomenon elephant, and the essence of the elephant under the phenomenon
inkwell. What help is it that we cognize the phenomena of things, when the
phenomena do not concern the essence at all? If a definite phenomenon is to
stand in a necessary relation to a definite essence, as semblance to reality,
then reality and semblance must correspond exactly to each other; the semblance
that corresponds to reality is a real semblance, and the semblance to which no
reality corresponds, a false semblance. Cognition must therefore aim at
discerning between the false and the real semblance. But of the real semblance
we have only one mark, namely this, that it corresponds to a definite reality.
Is there now any critique that can demonstrate a real difference between the
reality of the thing and its essence, between the real semblance and the
truthful phenomenon? The real semblance is the definite appearance of the
reality of the essence, and the reality a surety that the real semblance does
not deceive. Science can therefore console itself that, in the same proportion
as it is able to cognize the phenomena in their connection, it will also be
able to cognize the essential determinations and to have in these a relative
expression of the essence of things.

In the investigations of space and time, of the understanding and its
categories, the abstractions are sometimes driven to such an extreme that the
illusion plainly makes a fool of the critique. Time still gets off tolerably;
but space is continually exposed to the rack of the critical inquisition. That
space has precisely three dimensions, neither fewer nor more, is extremely
doubtful; for there are supposed to be sensing beings whose apprehension is
confined to one, at most to two dimensions, and there have already appeared in
science cognizers who could not make do with less than four dimensions. It
looks ominous. If one can really sense in a space with a single dimension, why
should one not, by proper economy, be able to bring it to sensing in a space
with no dimension at all? If one has already reached the point of discovering
the fourth dimension, then it is not unreasonable that one will in time be able
to discover a fifth, a sixth, perhaps an eleventh dimension; and what then is
space?

Still more ominous does it look with the categories of understanding. That the
understanding of one man can surpass that of another is a well-known matter;
but here it must expressly be emphasized that, although the inferior
understanding does not understand the superior, the latter is nevertheless
always able to understand the former and to make it clear that the laws and the
categories are the selfsame. Analyses, formations of concepts, and results may
on the higher stages of science be so inaccessible to those who as yet know
only the first rudiments that it almost looks as though the understanding
itself changed its nature at every step upward. What the knowledgeable person
sees with a single glance can cause the unknowledgeable incredible trouble
before he gets it divided up, pieced together, and understood. But when one
allows oneself to be misled by such a difference into assuming an understanding
so exalted that the categories which the understanding in us is obliged to
apply lose their meaning, then one finds oneself in the midst of the
meaningless. We can think of our categories as integrated by intermediate
determinations that we have not yet brought to consciousness; we can think of a
comprehensiveness in survey, a sharpness in discernment, a rapidity in
comprehension, a simultaneity in cognition of wholes, that infinitely surpasses
every human development of understanding; but we must constantly presuppose
that the categories and fundamental laws are originally the same.

What sense is there in the assumption that a perfect understanding has no use
at all for our categories? Here there is no scientific cognition, but, plainly
spoken, a \textit{salto mortale} of imagination that entertains the critique. In
order to see what such a perfect understanding is to mean, we must not merely
think of the boundaries of our faculties of understanding as infinitely widened;
no, we must go about it otherwise, we must think of our own understanding as
abolished and a wholly new understanding, an understanding with another essence
and other laws, set in its place. What then helps us to such an exchange?
Imagination! For only imagination can hit upon a counsel such as this: ``Give
up the understanding you have, then only do you cognize the understanding you do
not have.'' But the critique of imagination is an unreliable guide; in order to
grasp the understanding we do not have, we must let go the one we have; we are
thus to cognize the perfect understanding by departing from the understanding:
that is highly unscientific. The critique assures us that it is not meant so,
and we readily believe it, for imagination veils the meaning; but when the veil
is lifted, what then? Only by holding fast the fundamental law of
objectification and freeing itself from the illusions, the critical as well as
the speculative, will science gradually come to determine its essential
boundary and work its way forward to the solution of the problem of cognition.

\section{Complete and Incomplete Knowledge}
\label{sec:7}

\noindent As the circles of cognition widen, our knowledge becomes more
perfect; it might then seem that the perfect knowledge, which comprehends all
and sees through all, must overstep all boundaries and be wholly unbounded. Yet
it cannot be so. No knowledge is without an object, no objectification without
an object, no cognizing without a something that is cognized. Object, the thing
objectified, the something that is cognized, are different expressions for
boundary-determinations, in that they designate that in knowledge which is not
knowledge, but by which knowledge is determined. A knowledge without
opposition, without boundary, a knowledge of knowledge in the infinite, is like
a circle with an infinite radius. The difference between the perfect and the
incomplete knowledge rests upon an essential difference between inherent
boundaries and barriers.

A knowledge that has barriers is restricted, and restricted is all the
knowledge that is conditioned by sensation. Every sensible object is for the
sensing consciousness a barrier, the moon and the star as well as the wall and
the door. It is a fundamental law of sensation that consciousness must either
do without the realities or put up with their presenting themselves as
barriers; its imperfection shows itself precisely in this, that it cannot
itself produce them. If, then, there is a perfect knowledge, it cannot be
sensing; it must then have the realities in a quite peculiar way, namely by
originally producing them. We strike here upon a pervasive difference between
sensing and really producing objectification. The sensing objectification is an
objectification in impotence, corresponding to an impotent knowledge; for the
passivity in the one who senses is precisely an expression of impotence. The
really producing objectification is an objectification in omnipotence; the
production itself, the creative activity, is the expression of omnipotence. That
a creative activity cannot be seen through by sensing beings, though they stood
a thousand stages above man, follows of itself. But does there not lie in the
concept of a creative production an ensnaring contradiction, a serpent that
imagination has hidden in the grass, and that the critique will discover? If
this were so, we should have to halt at the result that only a restricted,
impotent knowledge can be true, while a perfect, omnipotent knowledge must
necessarily be something in itself untrue. Undeniably there is here a
contradiction, yet it lies, when we look more closely, not in the perfect
knowledge, but in the incomplete.

We have seen in the foregoing that the sensing objectification, though it take
thought to its aid and put all conditions, both physiological and
psychological, into requisition, can never come to terms with the object
itself. Every time we would let the non-objectified object stand, it vanishes,
and when we have brought it to vanish, behold, there it stands again. It is
really a contradiction, and this contradiction cannot be removed unless we
presuppose a perfect knowledge. Nothing is objective that is not objectified;
and yet the sensing consciousness gets a substratum left over which is
objective without being objectified: this enigmatic substratum of
objectification is reality itself, the reality that the powerful knowledge
posits and maintains. By means of this matter-at-issue foreign to us, our
original barrier, the restricted knowledge gets everything objective grounded,
not indeed directly upon itself, but yet indirectly upon knowledge, so that the
cycle of cognition can be closed.

It is impossible to get to the bottom of the incomplete knowledge and to give a
comprehensive account of its constitution and worth without constantly having
the perfect as an ideal \textit{in mente}. The rejoinder that the assumption of
a perfect knowledge is a blind representation, a loose hypothesis, says nothing,
provided one will but grant that glimpses of the perfect are necessary in order
to illuminate the imperfect.

Although in our knowing consciousness it sometimes seems to swarm with images,
representations, and thoughts, almost as with bees in a hive, yet it is, when we
look more closely, by no means so. On the contrary, it is only a few images,
exceedingly few representations, that the determining thought can hold back,
while the swarm vanishes. Of three representations only one stands at any time
in full illumination, while the two others are more or less darkened. When
several things at once attract our attention, and it is as though we were
simultaneously occupying ourselves in equal degree with them all, then it is an
imagining called forth by the rapidly shifting attention. Exactly speaking,
there is only one element that in the given moment can find place in the
foreground of consciousness, for the place is narrow. In proportion as the
whole is lit, the parts are dark, and when a single part steps forward, the
others must step back. Our knowledge is restricted by space and time; that is
an imperfection; but it is properly not space and time, it is the barriers in
space and time, that cause the imperfection.

It is an imagining that forms such as space and time should be excluded from the
perfect knowledge, and that our knowledge would become the more perfect the more
insignificant it found these forms. Perfecting is, on the contrary, conditioned
by the knower's capacity to render \textit{a~priori} and to appropriate the real
content of space and time. This can, in the case of sensing beings, naturally
happen only in approximation. We move in our consciousness from the one object
to the other, from the one representation, the one thought-determination to the
other, and the movement requires time. The more slowly and laboriously it goes,
the more imperfectly is time used; the more quickly and easily, the more
perfectly. The fuller the appropriation of content, the more will time change
from being the empty form that waits for filling --- the infinitely
long-drawn-out, in which the past and the present and the future are found in
sheer dividing-up --- to a continuity of a filled moment, in which the gathered
moments of the filling content reciprocally illuminate one another: time will,
in other words, approach eternity. How near the knower, absorbed in his tasks,
can bring the unrest of time to eternal rest, of that we are reminded by that
saying of Archimedes: \textit{noli turbare circulos meos!} [do not disturb my
circles!] One who is to orient himself in the heavens must begin piecemeal and
take each region for itself. Division in space is also division in time. It
takes time before, by the help of planes, axes, great circles, and spherical
coordinates, one can find one's way with the simultaneous and the successive;
but the astronomer sees the whole as with one glance: space enters into
limitation, and its immeasurable extension does not overwhelm the knower. Here
in the imperfect is a reflection of the perfect. The approximation proves the
validity of the perfect knowledge; but for sensing beings there is here a
barrier.

In so far as the sensing beings are at the same time thinking beings, one has
supposed it possible to be rid of finitude, to surmount the barrier, and to
reach the perfect by holding to thought as thought. With an intuiting freed from
sensible heaviness, thought can indeed race through time and space more quickly
than the lightning, faster than the light. In the pure forms of thought we
should thus be able to attain an unrestricted, absolutely perfect knowledge.
Unfortunately an illusion! The ``absolute knowledge'' that is conditioned by
continued abstraction from all realities is, despite its rich content of
determinations of form, without any content of actuality, without reality. By
taking even but a glimpse of reality along into a definite representation, one
will easily convince oneself that it does not stand with the much-praised
rapidity of thought quite as one usually assumes. Let an ell, for example, pass
as the measure of an immediately intuitable length of road. One can now, it is
said, easily think of a mile. But when a Danish mile is really to be traversed
in thought, it must take place by laying an ell to itself 12,000 times in a
straight line and letting intuition cancel the dividing-up. However much thought
hurries on, one still notices that it takes time. One may need a couple of days
to get the task practiced, and then thought, with the definite representation,
will still only approximately go together in intuition. And yet thought was
supposed to follow with ease the pace of light through infinite space?

The imperfection of our knowledge is, says the ``critique,'' to show itself in
this, that we cognize only the laws of things, but not their essence; of course
we have no concept of a perfect knowledge. --- Yes, we answer, so long as one
regards the laws only as universally valid determinations of the order and
connection in which things come to meet us, so long as one is still uncertain
how far these determinations lie in the things themselves or merely rest upon
our apprehension of the things, it is quite natural that a cognition of laws
and a cognition of essence must fall far apart from each other. It is
undeniably an expedient, when one would be rid of the endless doubting about
the general validity of the laws, to lay them into our own consciousness and
determine them as forms of understanding that are necessary for all experience.
When no experience is possible without laws, then one can be sure that the
things of experience will conform to the laws. But the conception suffers from
ominous weaknesses. Either the laws must, according to this mode of view, be
younger than the things or older than consciousness; but both assumptions are
equally unreasonable. One appeals to the fact that the things of experience are
not things in themselves, but phenomena for the consciousness that makes
experiences. In this there lies indeed the truth that phenomena and laws lose
all significance when one abstracts from an objectifying knowledge; but how it
comes about that the things as things, which must yet have a certain share in
the determination of the phenomena, should have no share at all in the
determination of the laws, is not cleared up. Have things in their essence
perhaps no laws at all, because there is no understanding in them? Or are the
laws of essence in themselves expressions of so perfect an understanding that
they do not at all suit our understanding? Both opinions are equally
meaningless.

Those who assume that the laws lie neither in the essence of things nor in our
consciousness, but are general formulas for the order we abstract from
relations that take place, degrade our knowledge to something so essence-less
that it can consist only in demonstrating, gathering, and ordering a swarm of
facts. The facts themselves cannot be understood, they must be taken
immediately as they are given; the order and connection upon which understanding
depends cannot be understood either, for the order found is itself a fact: thus
everything becomes unintelligible. One can indeed, by observations,
calculations, and inductions, convince oneself that something is more usual than
something else, something more general than something else, that something so
often precedes that it can be called cause, and that something else so often,
indeed without any exception known to us, follows that it can be called effect;
but that causes and effects really exist, or that the laws are essential,
intrinsically necessary, of this one cannot convince oneself. One uses the
understanding without knowing the understanding, and makes do with an assurance
without insight.

There is indeed an insurmountable gulf between the perfect knowledge and the
incomplete, and as a consequence an insurmountable gulf between the perfect and
incomplete understanding; but the gulf makes no breach in the laws themselves,
the eternal laws. Were science steadily to advance for billions of years, and
were we to live to see the last results: we should yet never reach a
simultaneous, both intuitive and distinctive knowledge of the real content of
the universe; but already on the stage where we now are we can cognize that our
subjective understanding in conscious form must be essentially one with the
unconscious understanding that has its objective expression in things, as surely
as we can cognize that our brains, our organs of consciousness, are subject to
the laws of nature. We can see that the objective understanding in things must
have its original ground in a knowing understanding, in a perfect knowledge; for
an understanding that no one understands, and that does not understand itself,
is impossible.

The incomplete knowledge must, so to speak, make itself a disciple of the
perfect, which, according to the essential unity of the laws of things and the
laws of thought, is also possible. Just as no thought has significance for
itself, but only in reciprocal relation to other thoughts, so too no existing
thing has significance exclusively for itself, but only in relation to other
things. The reciprocal determination of terms and relations has just as full a
validity for the perfect knowledge as for the incomplete. That the incomplete
knowledge must spell its way forward, while the perfect sees all as with a
single glance, makes no difference in relation to the essence of the laws. The
glance that sees the consequence in the premise, the derived in the original,
the parts in the whole, sees at the same time the essential difference; if the
difference vanished, everything would run out into one, and this one would
become nothing. But that unity has manifoldness and manifoldness unity, that
identity has difference and difference identity, that the parts are the
subsistence of the whole and the whole that of the parts, are precisely
essential truths which the incomplete knowledge must continually practice, in
order not to perish in sheer dividing-up. Each single law, apprehended for
itself, is an abstractum and may so far be said to lack reality; but the more
perfect knowledge becomes, the more clearly it shows itself that the different
laws in thing and actuality integrate one another. In order to secure a single
law, it is necessary to emphasize a common property of several things and to
regard their relation with respect to this property; but no thing can exist with
a single property: the body that falls, the magnet that attracts, are, while
fulfilling each respective law, subject in relation to light, heat, electricity,
and so forth, to other laws. The heavy body is more than heavy; it has a
multitude of other properties and is, by each particular property, implicated in
a system of corresponding laws.

Science sets out to reduce all laws to one fundamental law, and proves
precisely thereby that it has the perfect knowledge for its ideal. It is the
ideal of an unrestricted understanding, a comprehension of all things and all
laws under the unity of an all-seeing knowledge, that as a tacit presupposition
gives science its ideal, \textit{a~priori} foundation. Take this foundation
away, and science, with all its doctrinal structures, must lose its imposing
elevation, and the investigations sink down to serving a craving for knowledge
that is never sated; the goddess with the helm of Hades and the piercing gaze
must then lose her virginal bearing, bow her head, and resign herself to
treading the treadmill of the useful.

\chapter{Systematic Knowledge}

\section{System and Method}
\label{sec:8}

\noindent That all scientific knowledge, essentially seen, must be systematic
follows simply from the fact that the single links of cognition acquire
significance only when they are seen in mutual connection. The connected links
will then form a whole, a unity-of-manifoldness, consisting in this, that each
single link takes its place in relation both to the other links and to the
whole itself; such a whole is a system. How many links are necessarily needed
to form a system can easily be reckoned: their number must be greater than $1$
and less than $\infty$. That there must be more than one link goes without
saying, for there must be relation and connection; that no definite, existing
system can have infinitely many links is seen from the fact that a definite and
yet infinite number does not exist. When an infinite system is spoken of, one
can only be thinking of unsurveyable magnitude or of unceasing becoming.

A system of several links can be resolved into smaller systems, and several
smaller ones can again be composed into a larger. The manner in which this
happens depends upon the method, this upon the constitution of the connection,
and this again upon the constitution of the links. If the system consists of
existing links, and its composition stands in our power, then the method will
be: partly to form a whole of corresponding parts, partly to resolve both the
whole and its parts into a real unity, a unity of endowment, from which both
whole and parts must again proceed, partly finally to come to insight into the
laws of formation upon whose fulfillment the system essentially rests. The more
the conditionally necessary connection of the links has the character of
outward and finite combinations, the more the system will be marked on the one
side by contingency, on the other by arbitrariness. The more intimate and
essential the connection is, the more an all-governing necessity in the
formation of the parts out of the whole will impress upon the system its
peculiar stamp. If the connection is finally so pervasive that the parts become
moments in the whole, the links moments in the relations, then the development
of the whole will be the decisive thing, and the partial results of the
formation will at all points point back to the endowment; necessity will go
together with freedom, and we then have a teleological system.

That a teleological system must be the ideal of knowledge can be seen from the
scientific method. Method is a roundabout way, the way of objectification, of
cognition, from the first necessary presuppositions to the sought result. The
attainment of the result is the goal. The roundabout way and procedure that
characterizes the method is thus: the way to the goal. Every scientific task is
teleological, for the task sets an end, and the method indicates the means ---
that is, the conditions for its solution: the proof is teleological, for it is
methodical. But on the stage of finitude the objective relations in the matter
by no means correspond to the subjective movement in the method. In the
question of the relation between an exterior angle on the prolonged side of a
triangle and the two interior opposite ones, the procedure in the proof has
indeed the result for its end, but the magnitude-relation of the angles does
not conform to the course of the proof; this relation does not move, it does
not become, it is. It may be a task to investigate how water arises from its
constituents; but although the procedure in the investigation is teleological,
the movement in the matter cannot be said to be teleological; no chemical
experiment can establish that the oxygen combines with the hydrogen in order to
form water. A genetic system which from the objective side satisfies the
subjective demands of the method may indeed in part arise by formal
constructions, \textit{a~priori} productions; but in the world of realities it
does not exist. From this standpoint the critique of understanding, in seeking
to combat the so-called speculative or immanent method, which identifies the
movement of the concept and the movement of the matter, can beforehand be said
to have the victory in hand. Schelling's words, ``Ueber die Natur philosophiren heisst die Natur schaffen''
[To philosophize about nature means to create nature], do indeed yield a meaning, in so far as a
penetrating understanding of the thing is conditioned by a re-formation, an
intellectual re-production, in the course of which the object arises from the
ground up, link by link, and comes to be in consciousness. But the coming-to-be
is only ideal, not real; the Schellingian prophecy that ``the time will come
when the sciences will more and more cease, in that no longer man sees, but the
eternal seeing itself is seeing in man,'' will not be fulfilled: what he calls
``science itself'' cannot possibly outlive ``the sciences.''

It must not, however, be overlooked that the critical conception also has its
weak sides. By holding to fragments one may very well investigate facts and
laws without teleological considerations; provided one only presupposes a causal
relation between the facts, one has all one needs. If A (the premise with cause
and conditions) is, then B (the consequence, the result) must also be. If a
planetary mass contains the constituents of water, then under the conditions
that come about, water must form. One moves unceasingly between premise and
result, result and premise, and gets every whole grounded as a product of the
reciprocal action of the parts.

In what extent the scientific method can bring about a genetic system without
teleology, of this the Kant--Laplace theory of the formation of the solar
system is a grand example. But the genetic system is only half-genetic so long
as one confines oneself to demonstrating the coming-to-be of the whole without
demonstrating the coming-to-be of the parts. If the coming-to-be of the parts is
to be grounded just as well as that of the whole, then the system must
necessarily become teleological. It must then show itself that the parts as
well as the whole are developed out of a common endowment, that the real
possibility of the parts is determined precisely with regard to the whole ---
that is, that the result is preformed in the premise, and its attainment the
end of the development.

But, says the critique, such a doubly-genetic system does not exist, and cannot
exist. When we ask: why? we get two different answers. There must, in the first
place, be something that stands fast, a given to which one can hold and from
which, in any grounding, one can set out; but that is precisely what teleology
misjudges. First one assumes the premise in order to get the result grounded,
then one again makes the result a premise in order to get the premise grounded;
that is a subjective circular movement. In reality one achieves only making
premise and result equally hypothetical. In order to hide the hypothetical and
make the movement look like a progress, one slips in an all-determining knowing
and willing --- that is, one slips in the arbitrary; therefore all teleology
must be rejected by science.

When, against this, it is urged how unscientific it is to divide existing things
into two classes: such as must have a ground, namely the composite, and such as
have no ground, but are because they are, namely the single elements, the
constituents; and how inconsistent it is to make the existence of the
fundamental parts contingent in order to make the formation of the wholes
necessary --- then one gets another explanation. A teleological system is
impossible for sensing beings; all research sets out from something factually
given, and science must arrange its method accordingly. In this conception we
can acquiesce, for by it the whole mode of view is changed. --- When one will
not smuggle in the contingent in order to avoid the arbitrary, the teleological
system must necessarily continue to be the ideal of science. It then becomes the
business of the critique to watch over the respecting of the conditions of our
cognition, that one does not, by the flight of imagination out beyond the
barriers, lay hands on the ideal.

The dogmatic system, whose essential distinguishing mark is a naive assurance of
the perfect agreement of real being and human knowledge, has not seldom made
itself guilty of teleological blunders. For the teleological dogmatist it is a
settled matter that the parts exist for the sake of the whole, that an
understanding of their connection can be attained only when one has seen what
the meaning of the whole is. As cognition of ends becomes the decisive thing, it
lies near at hand that the knowing consciousness seeks the measure in its own
subjective essence and, without knowing it, subjoins to existence a meaning that
best corresponds to what human subjectivity must regard as the true, the good,
the beautiful. The dogmatist's naive confidence in his ``total view'' is seen
clearly from the way in which he treats the single phenomena and facts. No fact,
however refractory, embarrasses him; he does not give way before any case that
comes up. By re-explanation, abstraction, and explaining-away he gets the better
of everything that otherwise will not go up into the ``total view.''

The dogmatic system is on the whole predominantly synthetic. From a general
principle, be it teleological or non-teleological, it advances through particular
determinations, defining, reasoning, demonstrating, so that the whole conduct of
the proof gets a mathematical tinge (\textit{methodus mathematica}). With a
tinge of incontrovertibility the dogmatic system appears as a system of
authority, usually in such a way that the authority is outward, in so far as it
is due to a ruling tradition, but yet also in such a way that it rests upon the
superiority with which the total view is carried through, and so far is inward.

The old Ptolemaic structure of the world had the Aristotelian dogmatism for its
foundation. How hollow the foundation was became manifest when Copernicus broke
with the authority and got the system reduced to a hypothesis. It was a
turning-point not merely in the history of astronomy, but in the history of the
sciences. The more the critique gained power over the authority, the sharper the
observations of the real facts in all their particulars became, the more
analysis stepped into the foreground. The critique is no system, it is a method:
the critical method is essentially analytic. Synthesis goes from the whole to
the parts, from the universal to the concrete; analysis separates the whole in
order to investigate the parts. By a peculiar union of induction and abstraction
it makes its way through the concrete to the universal. Occupied with
investigations of the single and the laws, of the elements and their both real
and possible combinations, the critique has no particular sense for total views;
but it is vigilant at its post where it comes to testing the reality of the
presuppositions, the consistency of the procedures, and the strength of the
proofs.

Since synthesis and analysis mutually condition each other, it is seen from this
that no system can be either exclusively synthetic or exclusively analytic; the
opposition rests upon the fact that now the one, now the other of the two
moments forms the foundation. In recognition of this there has developed a
system which will be neither analytic nor synthetic, but analytic-synthetic from
first to last, a method which will be neither dogmatic nor critical, but
dogmatic-critical from beginning to end: it is the speculative system with the
dialectical method.

The strife that hence arises between critique and dialectic is of pervasive
significance; neither of the contending parties can either absolutely conquer or
absolutely lose, for victory and defeat depend upon the subject of contention.
One can rightly say that it is a strife about ground and existence, when one ---
be it well noted --- refers all determinations of essence to the ground and by
existence does not think of the concept, but of existing things. In the latter
case the critique will always keep the upper hand; the boundaries of the
critique are sheer boundaries of existence, the principle of contradiction
reigns, and opposite terms exclude one another. As regards the ground, the
non-existent that bears the existence, the fundamental principle of
contradiction will, on the other hand, not be able to prevent the difference
from collapsing into unity, in order again to proceed from unity. For the sake
of the methods we shall elucidate the strife by a few examples. That things have
properties, that they change, arise, perish, and so forth, on this critique and
dialectic are mutually agreed; but when it comes to determining wherein change
properly consists, how a thing stands related to its properties, and what the
ground is from which the existent arises and into which it again resolves
itself, then there arises a strife as for life and death. The critique asserts
that every thing is what it is, and cannot possibly be what it is not; dialectic
objects that if this is to hold unconditionally, then no thing can change,
nothing arise and nothing perish. In order to understand what change essentially
is, one must see that the thing as thing has its existence in sheer
contradiction. The thing changes because it is at once both something and other;
it changes because it is both one with and yet different from its properties; it
changes because it has a ground and is both one with and yet different from its
ground. That the thing changes means that its existence is its concept, and the
change its method; its determinations of being are sheer determinations of
thought. The movements in its change, movement of the matter and movement of the
concept, are in their difference identical; that is, the dialectical method is
immanent.

This identifying of the different, this mixing of the disparate, this manifest
breach of the fundamental principle of contradiction offends and outrages the
critique. If something is the same as another, then black is white, and white is
black, then the lamp is table and the table lamp, the stone man, and man stone!
Then common understanding is invoked, and common understanding must side with
the critique. Nevertheless the victory is not so brilliant as at first glance it
looks. If we hold to the pure, general determinations of thought, then dialectic
by no means breaches the fundamental principle of contradiction; on the
contrary, the immanent method sets out precisely to get the contradiction
thoroughly resolved. Its thesis is that something is not other; its antithesis
is that something is other; but it remains standing just as little at the latter
proposition as at the former, for the synthesis is that something is continually
in transition to other, that the transition, the change, is precisely the
resolution of the contradiction. In the analysis of pure concepts the critique
always comes off short.

There can be a subtle dialectic, but there can also be a subtle critique. In the
analysis of pure concepts the critique is always sophistical. It strikes
everywhere upon the contradiction, but instead of seriously grappling with it
and resolving it, it goes, by the help of all sorts of considerations
extraneous to the matter itself (\textit{diversi respectus}), around the task,
and entangles itself in contradictions. The fact is that the mutual relations of
thoughts, of concepts, are relations of infinity, the critical boundaries on the
other hand sheer boundaries of finitude. Dialectic is, in the sphere of
abstractions and of infinity, nothing else at all than critique carried through.

On the stage of existence --- that is, of finitude --- the critique, on the other
hand, is completely right. All boundaries of reality are boundaries of
finitude, critical boundaries. The moment the immanent dialectic ventures to
tread the ground of reality, the critique can take it captive. Where the
relations are finite, the finite considerations (\textit{diversi respectus})
must hold. What in one respect may be a right may in another respect be a left;
but therefore right is yet not left, and left not right. The critique teaches us
that to disparate spheres there belongs disparate cognition; the critique
teaches that disparate sciences give different systems and require different
methods.

\section{The Systematic Connection of the Sciences}
\label{sec:9}

\noindent With respect to certainty, grounding, insight, and connected
cognition, there is something common to all sciences, something that under
every form makes knowledge scientific. On account of the common-scientific
element, all sciences are as one science; one blood runs in their veins, they
are of one family, just as all the Muses are sisters. But with respect to stuff
and material, to task and procedure, they are highly different. This difference
does not merely touch the surface of the phenomena, it penetrates into the
essence of knowledge, it modifies the insight. Investigation of a kind of soil
gives another sort of insight than investigation of a historical event or of a
geometrical curve. One must therefore regard the connection of the sciences
partly from the standpoint of plurality and difference, partly from that of
likeness and unity.

Both modes of view have come forward one-sidedly in the treatment of what is
called the encyclopedia of the sciences. Already in antiquity attempts were
made to comprehend systematically, according to single disciplines, the circle
of science and art that every free-born Greek and Roman had to possess; from
the Greek division into two, music and gymnastics, there afterward proceeded
the so-called seven liberal arts: ``grammar, dialectic, rhetoric, arithmetic,
geometry, astronomy, and music.'' When thought stood above experience,
philosophy took up the separate sciences into its doctrinal structures and ended
with the Stoic tripartition: logic, physics, ethics. In the Middle Ages there
were compendia, more or less complete surveys of the entire sum of human
knowledge at the time, under titles such as: \textit{Summa, Speculum, Theatrum
vitae humanae}, and the like. Later one abandoned the systematic consideration
even to such a degree that one made use of an alphabetical arrangement of the
material and set up the encyclopedias as dictionaries of real subjects. The
French Encyclopedia (begun in 1749), edited by Diderot, was a grand literary
work, by which the spirit of the eighteenth century, with critical
understanding, combated in clear forms everything it regarded as blind
superstition and unwarranted traditions. But however valuable encyclopedic
dictionaries may be, it is yet only in an indirect and rather outward manner
that they promote any insight into the essential connection of the sciences.

For a general theory of science the encyclopedic attempts that are guided by
general, comprehensive principles are, on the other hand, of greater value;
although the one-sided is here too difficult to avoid. The first
philosophically grounded survey of the sciences was furnished by Bacon of
Verulam in the \textit{Novum Organon} (1620) and \textit{De dignitate et
augmentis scientiarum} (1623); but despite the ingenious fusion of thought and
image, speculation and empiricism, its scientific specific gravity is yet
considerably diminished when it is subjected to a critical examination. Of
deeply penetrating significance, on the other hand, are the presentations that
have either Kant's or Comte's or Hegel's principles for their foundation.

There are, according to the Kantian critique of reason, two different sources
from which our cognitions spring, sense and reason. From the former comes
sense-cognition, from the latter the pure cognitions of reason. The latter fall
into mathematics and philosophy. We thus get three different kinds of cognition:
empirical, mathematical, philosophical. Each of these modes of cognition has its
peculiar character: contingency and intuitability belong to the empirical,
intuitability and necessity to the mathematical, necessity without
intuitability to the philosophical.

``In the mathematical cognitions we become conscious of the
composition-forms of things, in the philosophical of their connection-forms
\textit{in abstracto}. Kant called the former: forms of the figurative
synthesis, that is, intuitable forms of combination; the latter: forms of the
intellectual synthesis, that is, merely thinkable forms of combination, laws.
The consideration of the different kinds of cognition leads to a scientific
architectonic, that is, to a doctrine of a system of all human sciences. The
three modes of cognition correspond to three different forms of system. The
pure philosophical form of cognition belongs under the categorical system; the
pure mathematical form under the hypothetical system, and the pure empirical
content under the conjunctive system. Philosophy is a cognition of mere
concepts. With such a mode of cognition no other systematic derivation is
possible than by a categorical subordination of particular cases under the
generality of the rule. Mathematics is a cognition by concepts in pure
intuition. In this intuition the least composite is the most original and the
ground from which new consequences are derived by ever new compositions. The
derivation of the composite from the non-composite happens here by hypothetical
combination of thought. Finally, pure empiricism has to do only with the
apprehension of the single given in sense-intuition. Each fact rests upon
itself, and the one stands beside the other in a conjunctive system that
co-orders parts in the whole.'' (Apelt: \textit{Die Theorie der Induktion}.)

After this scheme it will scarcely be possible to form any concept of the
peculiar character of the particular sciences or of their mutual connection.
Already the critical distinction between the three modes of cognition shows
itself to be highly questionable. How is one to conceive a distinction between
composition-forms and connection-forms carried through in such a way that in
reality two different sciences form, of which the one has sheer
composition-forms, forms of the figurative synthesis, sheer intuitable forms of
combination, while the other has sheer connection-forms, thinkable forms of
combination, laws? Can there then be sheer laws, relations of necessity without
intuitable terms, in the one science, and sheer compositions of intuitable
terms without laws in the other? Let one try whether so simple a geometrical
theorem as: two adjacent angles together make two right angles, allows itself to
be established by intuitable compositions without intellectual synthesis, and one
will instantly discover the impossibility. Two definite and so far intuitable
angles that one can lay off give only an example, a single case that can be
exchanged for any other cases whatever; now can one lay off all possible angles
with the same vertex on the same side of a straight line? In order to get the
theorem proved, one must necessarily resolve the figurative synthesis into the
intellectual. If we conversely choose an example among the connection-forms, the
merely thinkable forms of combination, can we then dispense with the
intuitable? That every effect must have its cause may indeed pass for an example
of an intellectual synthesis; but what is effect, what is cause? Unintelligible
words, when definite representations are not attached to them. But whence should
the representations come, if there were nothing intuitable from which they could
be drawn, no example upon which they could be supported? The truth is that in
every mathematical proposition there must be something that belongs to the
merely thinkable, just as in every philosophical proposition there must be
something that points to the merely intuitable. The only difference on which the
critique of reason can here rely is that the intuitable is far more prominent in
certain parts of mathematics, especially in synthetic geometry, than in the more
abstract forms of philosophy. Yet even if we let that distinction pass for a
distinction in principle, which is nevertheless grounded only upon an \textit{a
potiori fit denominatio} [the designation is made from the stronger part], how
does it stand then with the empirical? If the empirical sciences can carry it no
further than to a combination of facts, supported by ``contingency and
intuitability,'' where the one stands beside the other in a conjunctive system
without necessity --- for mathematics and philosophy lay claim to everything
that belongs to necessity --- what is one then to judge of their scientific
worth?

Auguste Comte, the founder of positivism, that is, of ``positive philosophy,''
sets out from the impossibility of reaching absolute concepts, and therefore
gives up every attempt to find the origin of the universe, the end of
existence, and the innermost cause of phenomena, but seeks instead, by combining
observation and the activity of understanding, to ascertain the real laws and
unchangeable relations in given phenomena. Positivism, which in all facts sees
appearances of unchangeable laws of nature, sets out to reduce the number of the
laws to the least possible, without, however, making any pretense of being able
to lead them back to a single fundamental law; it seeks, as far as possible, to
present the observable phenomena as particular cases of an unknown fundamental
phenomenon. From this mode of view there follows consistently an endeavor to
gather all sciences into one system. First, then, a distinction must be made
between theory and practice: ``positive philosophy occupies itself only with the
theoretical knowledge of nature, that is, only with the fundamental concepts, not
with their practical application. Within the domain of the natural sciences a
distinction must moreover be made between the general or fundamental and the
particular, derived sciences. Astronomy is fundamental, geology derived;
chemistry is fundamental, mineralogy derived, and so forth. The fundamental
natural sciences positivism now seeks to set up in an ordered sequence, such as
their inner nature entails. The sequence is developed according to an
encyclopedic law in such a way that mathematics is at once foundation for and
integrating part of philosophy; thereafter follow astronomy, physics, chemistry,
physiology, and sociology or social physics. These six fundamental sciences
constitute ``the hierarchy of the positive sciences,'' which ought to be laid as
the foundation for all future scientific development, because their inner
connection and unity comprehends the whole of human knowledge.

It must be granted that positivism has a sharp eye for the peculiar character of
the natural sciences and their mutual connection; but the peculiarity of the
sciences of spirit it does not know. It rejects the teleological principle, and
must already for that reason become inconsistent in its treatment of sociology.

From the standpoint of absolute knowledge Hegel's systematic survey is the most
grand, and yet in the deduction itself the most mistaken. ``In an ordinary
encyclopedia,'' he says, ``the sciences are taken empirically, as they are found.
They are therein to be completely enumerated and ordered by setting together the
similar, that which coincides under common determinations. The philosophical
encyclopedia, on the other hand, is the science of the necessary connection
determined by the concept, and of the philosophical coming-to-be of the
fundamental concepts and fundamental propositions of the sciences. The sciences
are, according to their mode of cognition, either empirical or purely rational.
Absolutely regarded, both are to have the same content. It is the goal of
scientific striving to raise that which is known empirically to the ever-true,
to the concept, to make it rational and thereby incorporate it into the rational
science. The whole of science then divides into three principal parts: 1) Logic.
2) Natural science. 3) Science of spirit. The sciences of nature and of spirit
can then be regarded as the system of the real or separate sciences, in
distinction from the pure science or logic, inasmuch as they are the system of
the pure science in relation to nature and spirit.'' To this then corresponds the
following scheme: The transition from logic to natural science is formed by
mathematics on account of space and time. Natural science falls into three
parts: a) Mechanics with regard to matter and motion. b) The physics of the
inorganic with regard to light, magnetism, electricity, the physical bodies, and
the chemical process. c) The physics of the organic, under which then again
belong geology, which treats the formation of the earth; botany and plant
physiology, whose object is vegetable nature; physiology, comparative anatomy,
zoology, which treat the healthy, animal nature, just as medicine treats the
sick. The science of spirit also falls into three sections: a) Spirit in its
concept, anthropology and psychology (language: philology). b) The practical
spirit, which expresses itself in deed and action, under which fall the doctrine
of right, political science, history. c) The completion of spirit in art,
religion, and science, to which correspond aesthetics, the doctrine of religion,
philosophy.

Here no far-reaching investigations are needed of the manner in which, in the
system, natural science is deduced from logic, and the science of spirit again
from natural science: a single utterance of the author is enough to clear up the
whole misunderstanding. ``The encyclopedia,'' says Hegel, ``is properly a
presentation of the general content of philosophy; for what in the sciences is
grounded upon reason depends on philosophy, whereas what is in them of that
which rests upon arbitrary and outward determinations, or, as it is said, is
positive and statutory, lies outside it.'' However reasonable and sensible such a
declaration may at first glance seem to be, it shows itself, on closer scrutiny,
to be both ambiguous and misleading. That it is ambiguous strikes the eye when
we compare it with the fundamental thought of positivism. Comte and Hegel seem
in this matter to be in perfect agreement, and yet what a difference! With Comte
it is, in the name of reason, mathematics that is both a science for itself and
the foundation of all sciences; with Hegel it is logic. What the philosophy of
the absolute posits in logic as the content of absolute knowledge, positivism
must, according to its principle, reject as absolute nonsense --- from both sides
out of regard for reason and for the rational element in the sciences. How
misleading the utterance is, is clearly seen when we go into the concrete.

That every particular science must have its peculiar starting-point, conditioned
by stuff and subject, follows from the nature of the case. An immanent logic can
by abstraction begin with ``pure being,'' but neither astronomy nor mineralogy is
able to begin with pure being. At the starting-point it is at once seen how far a
science is empirically observing or logically abstracting. Every researching,
observing science must necessarily have its own method, its own goal. All
sciences indeed have a highest common goal, namely the cognition of truth; but
real truth is not in the merely ideal and general, it is in the real, in concrete
actuality. There are in reality many starting-points and many ways toward the
final goal. Every particular science must therefore pursue its goal and guard
against being dissolved into an abstract common goal. The ideal of natural
science is empirically-researching natural science, not natural philosophy; the
ideal of history is historical, not philosophical. All sciences are willing to
serve and help one another reciprocally; they are willing to serve, enrich, and
help philosophy in the best manner; but to empty into philosophy like rivers
into the sea, to make themselves mere auxiliary sciences, ``handymen'' under the
guardianship of philosophy, this they will not, cannot do, and that for good
teleological reasons.

\section{The Division of the Sciences}
\label{sec:10}

\noindent Every science has its own circle of investigation, its own method, and
its own end; its highest task is to reach its peculiar goal, its greatest
perfection to be rightly scientific. It is in this respect impossible to
introduce either a caste-system or an order of rank into the republic of the
sciences. If one asks: which are the most important sciences? the knowledgeable
will in general answer that they are all equally important; but if we ask each
single cultivator of science which science he, for his part, must regard as the
most important, the mathematician will undoubtedly answer: mathematics! the
botanist: botany! the numismatist: numismatics! and each of them with equal
right. When nevertheless a distinction is made between higher and lower, more and
less comprehensive sciences, between principal subjects and subsidiary subjects,
and so forth, then the distinction has its ground not in a higher or lower
degree of scientific character, for unscientific or half-scientific disciplines
cannot be reckoned among the sciences.

The division into higher and lower, superordinate and subordinate sciences rests
essentially upon a consideration of totality, in that something must be central
while something else is peripheral. The organism of the sciences must have a
center of gravity. But the center of gravity shifts in the course of time:
antiquity set it in philosophy, the Middle Ages in theology, the present finds
it in natural science. Nevertheless the matter is questionable. To make a single
trend-setting science the highest and ruling one is about as precarious as to
wish to attribute to the central atom of the sphere the real weight of the
sphere and to weigh it for itself, as proof that the whole weight is concentrated
in the center. The highest, ruling element in the circle of the sciences is not a
science for itself, it is the spirit of the sciences, the scientific element
itself, that animates them all. When this highest, this superordinate element is
to step forward in the form of a particular science, a science for itself, it
must necessarily get the other sciences outside it, and cannot be the
all-penetrating, all-animating.

What advantages and perfections would not that science have to possess which in
truth were to be fitted to be unconditionally sole-ruling and to transform the
republic of the sciences into an absolute monarchy? It would have to be not
merely the most comprehensive, but at the same time the most concrete, not merely
the most spirit-filled, but at the same time the most scientific of all. But what
science dares to undertake to satisfy these conditions? Philosophy it cannot be;
for philosophy knows with itself that in the pure cognition of essence elements
of actuality must always perish; it is clearly conscious that the real moments it
assimilates into its ideal content are bought and paid for with loss of reality,
so that its most concrete determinations, even when it does its utmost, must be
conditioned by abstractions and fall within the abstract. Natural science it
cannot be; for natural science in its generality is an abstraction. In reality
natural science consists in many sciences. The most spirit-filled and most
scientific of these sciences would then have to be able, by its observations,
calculations, and inductions, to instruct the other sciences about what nature
originally is, and how the relation in principle between the things in nature and
nature itself must necessarily be determined. But the doctrine of nature, in
whatever principal subject we meet it, is far too soberly strict, far too sure of
its tasks and its method, to confuse itself with metaphysics and engage with a
problem that lies outside its domain. Where then is the highest, governing
science? It surely could never be theology! Yes, suppose that theology, with its
unfathomable fullness of spirit, awakened from medieval slumber and dogmatically
reborn to new life in the baptism of speculation, seated itself on the throne and
began to sing in Homeric fashion: οὐκ ἀγαθὸν πολυκοιρανίη· εἷς
κοίρανος ἔστω! [the rule of many is not good; let there be one ruler!] --- it
would be an enlivening sight. All the sciences would surely bow deeply to the
favored one; yet, be it well noted, on the inviolable condition that it, instead
of singing, proved in deed that it among all the sciences really is --- the most
scientific.

The positivist division is not satisfactory. The distinction between fundamental
and derived sciences is indeed of significance, in so far as the former furnish
the foundation for the latter; but already the circumstance that here, besides
mathematics, the common foundation, five other foundations are needed, shows how
precarious it stands with the deduction. If the five particular foundations are
to be derived from the common first one, then a distinction must be made between
the original and the derived foundations; but if one assumes derived foundations,
then the ``hierarchy'' will come to consist of more than six sciences. Wherein
consists the deduction? The derived must have its ground and its general
conditions in the original, must show itself as the original's own development,
modified by synthetic additions; thus the formulas for the different conic
sections are modified developments of their fundamental formula, thus the crystal
forms in a system are derived forms of the system's fundamental form. How will
one now be able to ground the assertion that chemistry, physiology, and sociology
are only modifications of mathematics? One cannot even derive astronomy from
mathematics. It must indeed be granted that neither astronomy nor physics will be
able to get its scientific grounding and complete development without the help of
mathematics; but this is something quite different. That mathematics is an
indispensable auxiliary science for astronomy and physics is incontrovertible;
but in these empirical sciences the mathematical propositions are, essentially
seen, only borrowed propositions. In whatever extent they are applied, and
however systematically they are appropriated, they yet do not constitute the
foundation itself. It is the observation of the real phenomena, the real stars,
the physical properties of the real bodies, and the relations that really take
place, that give the named sciences their peculiar character and real
foundation. Only in a highly imperfect manner can the sciences be derived from
one another.

The general foundation, \textit{fundamentum divisionis}, must be sought in the
general presuppositions of knowledge: the knower, the object of knowledge, and
the unity of consciousness of both. The principle will then divide itself into
different principles. If the emphasis is laid upon the unity of knowledge with
the knower, then we get the \textit{a~priori}; if the emphasis is laid upon the
relation of knowledge to existing objects, that is, upon sensation and
sense-impression, then we get the empirical. The sciences must then, with respect
to this, divide into \textit{a~priori} and empirical. The \textit{a~priori}
again takes on a different character, according as the representations either
resolve themselves into general thought-determinations, or furnish exact terms of
relation, or enrich intuition with imaginary masses and forces; to these
correspond: logic, mathematics, and rational mechanics. Likewise the empirical
will take on a different character; it will be modified peculiarly, not only by
the fact that the objects present different principal aspects from which their
law-bound reciprocal action is to be regarded, but in principle by the fact that
the real cognition of law can partly acquiesce in the relation between given
premises and the results proceeding from them, partly must understand a
self-activity which through the fulfillment of the laws strives toward a definite
goal. With respect to this we divide the empirical sciences into the
real-empirical --- mechanical and chemical physics, chemistry, and geology ---
and the teleological-empirical --- botany and zoology, anthropology and history.

The fundamental relation between the \textit{a~priori} and the empirical points
in the last instance to a fundamental relation between necessity and freedom,
which must be treated in the doctrine of spirit. If freedom itself fell outside
necessity, then every system of cognition that sets out to maintain and
illuminate the ideal of freedom would without further ado be declared
unscientific, for necessity is, so to speak, the only fixed ground upon which a
scientific structure can rest. But from the \textit{a~priori} in the relation
between the content of knowledge and the knowing self it must already be possible
to see that freedom dwells in necessity and works its way forward out of the
ground. That all knowledge originally rests upon a self-activity, in so far as it
rests upon sensing and thinking, has already been shown in our foregoing; that
this self-activity is by no means interrupted by the necessary, the objective,
but on the contrary is conditioned thereby, has likewise been emphasized. If the
perfect could be reached, it would end with the knowing subjectivity, in
objectifying everything else, really objectifying itself and the demands of its
own essence: we should then have a freedom identical with necessity.

But on the stage of imperfection, on the conditions of finitude there must, as we
have seen, remain a difference between the sensing, objectifying subjectivity and
the reality of the objects that correspond to its objectification. In this lies
the ground that science, with its grounds of necessity, can cognize only the
relativity of freedom, and must dissolve itself in proportion as the knower is
absorbed in the ideal freedom. The doctrine of spirit shall then show how the
objective science, step by step --- through aesthetics, ethics, and the
philosophy of religion --- accomplishes its own dissolution in the ideal itself.
As regards, finally, the relation between plurality and unity, there is ground to
divide the sciences into special sciences and the science of spirit, yet in such
a way that the ground of division is not fixed critically, but apprehended
dialectically. Critically and outright the division does not hold; for all
sciences are sciences of spirit and must go up into the unity, as surely as the
scientific spirit penetrates them all. Conversely, the unity of the science of
spirit is only an abstraction when it does not divide itself into a plurality of
sciences of spirit. The division is motivated by the fact that the boundary of
the special sciences is critical, that of the science of spirit dialectical.
Every science, philosophy, ``the science of the sciences,'' not excepted, must as
sharply as possible bound its compass, follow its method, and hold fast its goal;
precisely thereby it becomes a special science, and the one special science
cannot be dialecticized out of the other. With spirit and the spiritual moments
it is the reverse. Since the development of spirit is the common goal of all
sciences, and the method depends upon the goal, all the particular content that
the science of spirit takes up from the special sciences into its ideal
concretions must at the same time go in under the unity. Even as the science of
spirit divides itself into separate disciplines, it will yet show itself that the
one discipline can be transformed into a moment that is taken up into the other,
so that the boundary becomes dialectical, and the unity of totality decisive.


%%% PART TWO: THE SPECIAL SCIENCES %%%
\part{The Special Sciences}

\chapter{A Priori Sciences}

\section{Logic}
\label{sec:11}

\noindent We know no thought that is not a product of thinking, and no thinking
that is not a function of a thinking being. Since, however, one can find
thoughts expressed, not merely in words, but in real products, it becomes
possible to enter into a system of thought without asking who the thinker is.
This alone stands fast, that thoughts cannot think themselves, but must
presuppose a thinking subjectivity. Logic, the doctrine of thinking, can with
respect to this be presented from three standpoints. One may hold to man as the
thinking subject and abstract laws and forms from the acts of thought; one may,
without attending to the self from which the thoughts proceed, analyze the
universally necessary forms of thought and develop them in a system; one may
resolve the finite thinking into the infinite in such a way that the thinking
self-unity shines through in thinking idea-forms. Hereby is determined the
relation between formal or subjective logic, immanent or objective logic, and the
logic of the fundamental ideas.

Formal logic divides into the doctrine of concept, judgment, and syllogism.

The concept is formed in a human manner, in that thinking, which originally,
though unconsciously, is present along with sensation, specially raises itself
and, by comparison and abstraction, resolves the immediate sense-images into a
hovering universal image that must be held fast in the word. The universal image
of a tree, for example, cannot be drawn; the indefinite outlines show that it
finds itself on the transition from representation to concept. The sum of the
determinations belonging to a concept constitutes the concept's content; the
number of objects to which it finds application, its compass. By noting the
inverse relation between content and compass, one comes to a distinction between
superordinate and subordinate concepts, genus- and species-concepts, in that the
genus-concepts are seen to have greater compass but less content, while the
reverse holds of the species-concepts.

The judgment is formed in a human manner, and the concept finds its logical
expression in language: words are combined into sentences, and in sentences the
judgment is expressed. The judgment must then formally state: the subject, the
representation that is to be determined; the predicate, the determining
representation; and the copula, the determining connective word. From the
subject's connection with the predicate by the copula it follows that the
judgment can be divided into a plurality of logically different judgments; and
the judgments are classified according to quality: into affirmative, negative,
and infinite; according to quantity: into singular, particular, and general;
according to relation: into categorical, hypothetical, and disjunctive; according
to modality: into assertoric, problematic, and apodictic. If the predicate
proceeds by formal development of the subject, then the judgment is analytic; if
the predicate can be joined to the subject only in consequence of a particular
intuition, be it \textit{a~priori} or empirical, then the judgment is synthetic.

The syllogism is formed in a human manner; for the judgment must be grounded: the
act of thought by which the validity of one judgment is derived from one or more
other judgments is a logical syllogism. If the syllogism has only one premise, it
is called immediate; if it has several premises, then mediate. The mediate
syllogisms are classified into: the categorical, conditioned by the subordinate
concept's distinguishing mark (\textit{nota notae}) being subsumed under the
superordinate; the hypothetical, which rest upon the dependence of the
conditioned upon its condition (\textit{modus ponens} and \textit{modus
tollens}); and the disjunctive, which are grounded upon the fact that the terms
of a whole are mutually so related that, when one is expressly posited, the others
are excluded (\textit{modus ponendo tollens}), and that, when the others are
excluded, the one remaining must be posited (\textit{modus tollendo ponens}). The
restriction of the thinking subjectivity is, in concept-formation as in judgment
and syllogism, recognizable by the fact that thinking and being everywhere fall
apart from each other.

Immanent logic, founded by Hegel, abstracts from the thinking subjectivity,
identifies the real determinations of being with the general determinations of
thought, and develops the more concrete concepts out of the more abstract so
dialectically that the movement of thought from ``pure being'' to ``the absolute
Idea'' takes on the appearance of a continuously advancing movement of the
matter. The system divides into the doctrine of Being, Essence, and Concept.

\emph{Being.} General being is one with general thinking; but the unity of
thought and being is in its more abstract form contentless: pure being is the
same as pure nothing. And yet being, the positive, is different from nothing, the
negative. That being both is and is not nothing means that being passes
unceasingly over into nothing, and nothing into being: this transition is
becoming. In becoming, nothing is a boundary for being, and being for nothing.
The being determined by its nothing is a being-thing, a something, and the
nothing determined by being, an other. Something is different from other, for
other is its boundary; and yet the same as other, for other is itself something,
and something again other; something passes, in consequence of its boundary,
over into other, and other into something: this transition is change. In change
something is one with its other and so far unchangeable. It is thus the
unchangeable that changes. All the unrest seen in being and existence is
intelligible from the fact that boundaries such as nothing against being, other
against something, many against one, changeable against unchangeable, are for
thought pure negations, while they yet show themselves in representation as
immediate terms. This contradiction is removed in that the negation and the
immediate shine through each other in reflection.

\emph{Essence.} Reflected being is at once essence and semblance. The
unchangeable that reflects itself in the changeable is essence; the changeable,
fleeting, that reflects the unchangeable is semblance: the reciprocal play of
semblance and essence is reflection. In reflection identity has difference, and
difference identity. In its identity with the essence the semblance is
phenomenon; in its difference from the phenomenon the essence is ground. The
phenomenon goes to its ground; but precisely by getting the essence for its
ground it becomes grounded and posited as existence. The existence reflected into
itself is an existing thing, a thing; but the thing can be reflected into itself
only by reflecting its existence in other existing things: the determinations
corresponding to these reflections are properties. Under the reciprocal
reflecting the things resolve themselves into properties, and everything becomes
unstable. But the unstable presupposes a subsisting that cannot be resolved,
because its subsistence infinitely confirms itself in the changing of the
unstable; this subsisting is the substance, and the modes of appearance changing
in instability are the accidents. The relation between the substance and its
accidents is a relation of causality: the substance is cause, the original
matter; the accidents are effects. But the cause has its substantial expression
in the effect: the accidents are themselves finite substances. The effect is
thus an effect of the one substance upon the other, which is to say: an effect of
one cause upon another cause, that is, reciprocal action. In reciprocal action
the effect is cause and the cause effect; by the fact that all the moments of
actuality contained in the essence are, under all-sided reciprocal action,
developed to full totality, the essence developed in totality is concept.

\emph{Concept.} The self-comprehending totality is, in its identity with itself,
the universal; in its difference from itself, the particular; in its totalized
determinacy by itself, the singular. The self-development of the concept is its
fundamental division and mediation, that is, the judgment and the syllogism. By
means of the fundamental division the totality is, in its first immediate unity,
the singular, posited as the subject, and the totality in the reflected form, the
universal, posited as the predicate: the subject is what the predicate is, the
singular is the universal. But by the fact that the singular is comprehended as
the universal, and the universal as the singular, the opposite terms of the
conceptual division become mutually comprehended in reciprocal determination as
the particular. The outer terms of the fundamental division, the extremes, have
thus, in the course of the concept's self-development, set down a middle term, a
medium, in which they can be mediated. This mediation is the syllogism. The
concept's self-mediation through judgment and syllogism is its subjectivity; the
immediacy regained by the mediation is its objectivity. As immediately existing
totality the concept is object, and, in consequence of the fundamental division,
an infinite manifold of existing objects. The objects in their outward,
indifferent reciprocity are mechanical, acting on one another by pressure and
impact. The distinction hereby arising between motion and rest is a relative
appearance of the absolute mechanics in which motion is rest, and rest motion.
--- Under the indifference of the mechanical objects there hides a difference
that comes to light in chemical attraction. The difference of the objects entails
a tension that is equalized in the chemical product. But the product, in which
the difference is neutralized, must as the neutral again resolve itself into
difference, and the process begin afresh. In the chemical process, however, the
concept proceeds from its own presuppositions, cancels the presuppositions in
their result, in order again to cancel the result in the presuppositions. The
concept thus developed is now master over the stuff: the result anticipated in
the presuppositions is end, the presuppositions with their intermediate
determinations are means, the movement is teleological. The teleological concept,
comprehending itself in its elements, is a form that gives itself content, a
content that gives itself form; it is in absolute unity of form and content, of
subjective and objective, of end and means; in one word: the absolute Idea.

The logic of the fundamental ideas sets out for a further development and closer
determination of the content of the absolute Idea; the particular fundamental
forms under which the content is apprehended are Knowledge, Power, and Truth.

\subsubsection*{a) The Idea of Knowledge.} A thought that no one has thought, a
concept that no one comprehends, an idea that no one grasps, a knowledge that no
one knows, are absolutely meaningless imaginings. If one tries to veil the
meaningless under turns of phrase such as these: the thoughts of being think
themselves, the objective concept comprehends itself, the Idea moves itself with
absolute knowledge --- then the contradiction shines through; for this ``itself''
and ``itself'' points at all points to a self-reflecting unity, a unity that in
all thoughts thinks itself, in all concepts comprehends itself, a self-unity that
in infinite knowledge grasps, bears, and moves the Idea with its whole content.
Why such an absolutely knowing subjectivity must necessarily be, and why it
cannot be the human subjectivity, is logically recognizable from the resolution
of reflection, judgment, and syllogism into the Idea of Knowledge.

\emph{$\alpha$) The Resolution of Reflection.} For the finite subjectivity the
sensible is an immediate other for sensation, just as the thinkable is an other
for thought. If now this sensible and thinkable other is apprehended as an
outward boundary, as a barrier, then the sensing-thinking subjectivity gets a
world of sensible-thinkable objects outside it. Dazzled by this outside of
immediacy, the finite knowledge then assumes that the sensible stands
beforehand fully finished and merely waits to be sensed. This waiting then
signifies no defect or essential lack. If sensation comes, then it is well; the
sensible takes no harm from it, suffers nothing thereunder: it gladly permits the
one who senses to apply his apprehension, for it is there to be sensed, it is
sensible. But if sensation does not come, then that too is well; the sensible has
nonetheless its complete determinations of being, its marked forms of actuality
in itself. As it goes with the sensible in the thing, so it goes also with the
thinkable. Therefore one dreams of objective nature-thoughts that are without
being thought: a telling proof of the restriction of finite knowledge. What is to
be done about this? The critique senses mischief, but does not come out beyond
the finite.

The resolution of the finite reflection into the Idea of Knowledge is the
resolution of the contradiction. Sensible and sensation, thinkable and thought
are correlative determinations. What then is the sensible for itself? An
indefinite that must allow itself to be determined --- by sensation. If the
sensible is held fast in all its indefiniteness, then it is not an existing
thing, for every existing thing is a definite being: the as-yet indefinite is not
in actuality, it is only a possibility. The sensible is a possibility, namely for
sensation, a possibility of sensation, just as the thinkable is a possibility of
thought. The reverse also holds, that sensation for itself is only a possibility
for the sensible, and thought for itself only a possibility for the thinkable.
Where there is nothing sensible, there can be no sensation, and where there is
nothing thinkable, there can be no thought. To sense a nothing is impossible, and
when thought thinks its nothing, its negation, this nothing is itself, by means
of the word, a thinkable and acquires existence for thought. In the Idea of
Knowledge the finite distinction of determinations of being and determinations of
thought is resolved, as surely as there is a knowledge whose reflection is
infinite.

\emph{$\beta$) The Resolution of Judgment.} From the standpoint of the finite
subjectivity the subject in a judgment will, despite all mediations, reach full
identity with its predicate. That the subject categorically is its predicate gets
its fullest expression in the concrete form: the individual is its species. All
the determinations belonging to the individual are sensible, all the marks
corresponding to the species as species are thinkable: the individual is sensed,
the species is thought; thus the judgment expresses at once the difference between
sensible and thinkable, and also the unity. The finite knowledge, restricted by
its outside, cannot, however, come to clarity with itself concerning the relation
between subject and predicate. It looks as though the individual, the sensible,
were properly the first and original, the species, the thinkable, on the other
hand the derived; for the sensing and judging subjectivity must set out from a
sensible apprehension of the individuals and abstract the species from the
individuals. Thus the existence of the individuals is outside, the abstracted
species-concept on the other hand inside consciousness. Sensation (the sensing
consciousness) asserts that it is the individuals that determine the species,
while thought for its part holds fast that it is the species that determines the
individuals; sensation asserts that the individual is the original, the species
the derived; thought asserts that the species is the original, the individual the
derived: here is tension and strife between sensation and thought. The resolution
of the finite judgment is the resolution of the contradiction. The two factors,
the sensible and the thinkable, are namely both taken up as moments in the
infinite act of selfhood by which the absolutely knowing subjectivity originally
posits the content of judgments identical with the form of judgment.

\emph{$\gamma$) The Resolution of the Syllogism.} If it were merely a matter of
setting up a syllogistic scheme such as this: the individual, I, is the genus, G,
by means of the species, S: I\,--\,S\,--\,G; the species is the genus by means of
the individual (S\,--\,I\,--\,G); the individual is the species by means of the
genus (I\,--\,G\,--\,S) --- then idea-cognition would belong not only to the
easiest, but also to the most superficial of all scientific cognitions. The
essential task comes to light only when one reflects on what such a syllogistic
scheme properly means. It means namely that of the syllogism's three terms ---
individual, species, and genus --- none is the immediately first, and none the
last, but that all three must have the same originality. The sum of
sense-determinations that is gathered in the singularity of the individual, I, is
determined beforehand by the totality of thought-determinations contained in the
particularity of the species, S; but this content of thought is again a peculiar
concretion of the universal-concrete content of the genus-concept. So the
formation of the genus-concept would then be the unconditioned first. But since
this concept in the Idea must contain all the determinations for the species, the
formation of the genus-concept must necessarily presuppose the species-concepts.
So the species-concept would then be the absolutely first; but that is
impossible, since every species-concept presupposes partly its genus-concept,
partly the individuals whose species it is. Here is the contradiction. The
individual cannot be formed before the species and the genus are formed, for, seen
in the Idea, the individual is the immediate syllogistic unity of species and
genus. The species cannot be formed before the individual and the genus are
formed; for, seen in the Idea, the species is the reflected syllogistic unity of
genus and individual, and must expressly have the two extremes for its premises.
So it is surely evident that the syllogistic system, which is to express the
content of the Idea in the form of Knowledge, must spring from a single
trilogical syllogistic act, and that such a syllogistic act can be accomplished
only by an absolute Knower.

\subsubsection*{b) The Idea of Power.} That immanent logic is not able to make
intelligible the identity of thinking and being, because its absolute Idea lacks
an absolutely knowing subjectivity, must by now be clear from our foregoing; but
that it is even less able to do so when its general-logical being is to be
transposed into real existence, because its absolute Idea lacks an absolutely
powerful subjectivity, the following will show. While positivism immediately holds
to facts and builds upon realities, the problem of reality has not even come to be
in the philosophy of the absolute. This problem contains such hard oppositions,
indeed such cutting contradictions, that mere thinking, be it ever so dialectical,
cannot possibly solve it in the form of thought.

\emph{$\alpha$) Thing and Concept.} The contradiction is that thing and concept
both are and yet are not the same. They are the same; for one cannot think any
idiosyncrasy decisive for the thing that should not belong to its concept, and
one cannot find in the thing's concept any determination that should not
essentially belong to the thing; and yet they are not the same, for the thing is
sensible, but the concept non-sensible. They are the same, for concept and thing
have the same content, and yet they are not the same, for the concept, the
universal, has an infinitely greater compass than the thing, the existing
singular: one concept can be in many things. They are the same, for the existing
thing is the existence of the concept, and yet not the same, for the existing
thing is no concept, and the concept no existing thing. A stone is no concept, and
the concept no stone; one can crush a stone, but not the concept of the stone; one
can boil a plant, but not the concept of the plant; one can ride on a horse, but
not on the concept of the horse. The concept can be thought, and one can think
about the thing, but the thing as existing thing cannot be thought; the concept
can be expressed in words, but the existential determinacy of the thing is
unutterable: the concept is logical, the thing is anti-logical.

\emph{$\beta$) The Logical and the Anti-logical.} When it is assumed without
further ado that it is immediately the thing that is thought, and that the word
exactly expresses what the thing is, then this is a deception that the logical
judgment occasions. What now is This? This is --- This! And yet the pointing
designation can expressly show that This is not This; for This, for example, is a
house, but This is a horse, and This is again another This. By pointing at the
object the word points away from the word, the thought away from the thought:
this is the anti-logical. The immediacy that separates thing and concept cannot be
logical; for although the logical can point out beyond itself, it always happens
within the sphere of the logical. The immediacy can also not be illogical; if the
determinations of the object were illogical or stood in contradiction with the
logical, then the object would have to be either concept-less or stand in
contradiction with its concept. Nor can the immediate in the object be something
that is added or joined from without to the logical, as though one could separate
the object itself from the concept expressed in the object. What in the object
could that be which remained when the concept departed? The material, the matter?
But if the matter with all its determinations falls outside the concept, how then
can the concept be expressed in the material object?

If we seek the logical root from which the anti-logical springs, then we find it
in the category Power. Power in its immediacy is the expression of reality, and
the expression for the immediacy of reality is again the anti-logical. The ground
of our regarding objects as realities and attributing to them an immediate
objective independence is precisely the impression of the power with which they
each come to meet us. If the world of things were only an objectification of an
absolutely knowing subjectivity, it would be resolvable into a system of concepts,
of laws; but that things are realities means, in other words, that they are
objects of an absolutely powerful objectification.

\emph{$\gamma$) The Outburst of Power.} The transition corresponding to
objectification, from ideal to real, from concept to thing, has no intermediate
determinations and can have none. How would these intermediate determinations be
constituted? Logical they cannot be, for one can enrich a concept with as many
logical determinations as one will, it remains as concept within the logical.
Anti-logical they cannot be either, for the anti-logical first comes to be in the
transition itself. The transition must then be a movement of power, a transposition
of the ideal reality of power into real power, a leap by which nothing is leaped
over, a transition from possibility to actuality, understood in such a way that the
possibility lies in the concept, that is, in the logical, the actuality on the
other hand in the anti-logical reality, that is, out beyond the logical. The
conditions that are mediating in the transition must themselves belong to the
real; thus: either the ideal or the real, either concept or reality; an
intermediate matter does not exist. To wish to condense concepts into realities or
to sublimate realities into concepts is a speculative botching, equally at odds
with experience and with sound logic. The being that mediates infinitely between
ideality and reality is the self-being absolutely powerful in its knowledge, the
general, all-objectifying subjectivity in the Idea of Power.

\subsubsection*{c) The Idea of Truth.} What is Truth? In the form of reflection
truth is the unity of thinking and being: thinking has its truth in being, and
being in thinking; untruth is thus discord between thinking and being. In the form
of judgment truth is the identity of subject and predicate: the object, the
subject, has its truth in the concept, the fully developed predicate, and
conversely; untruth, the contradiction in the judgment, is thus strife between
subject and predicate. In the form of the syllogism truth is the middle of the
extremes: the more independently the extremes come forward, the fuller is the
middle, and the deeper is the truth. The infinitely content-rich middle, the
eternal ground from which the intellectual world of Knowledge and the real world
of Power proceed as extremes in the finite and join into articulated unity in the
infinite, is the Idea of Truth.

\emph{$\alpha$) Idealism.} The scientific distinguishing mark of idealism is that
the whole content of truth goes up into the form of science. If the emphasis falls
methodically upon the origin of concepts, that is, upon the comprehending
subjectivity, a finite I with barriers (Fichte), then idealism is subjective; if
the method, on the other hand, is to let the universally valid conceptual content
unfold itself in the development of the conceptual determinations (Hegel), then
idealism is objective. The untruth in subjective idealism is that it is not able to
develop the objective forms of existence out of the subjective I-forms; its lack of
formal development is the essential ground of its lack of content. --- For the sake
of content objective idealism begins with the very most abstract forms of being:
Being, Becoming, Existence \ldots, integrates the abstract by ever more concrete
determinations, and behold! the objective self-unfolding of the ideal content takes
place as by an enchantment. But what is the truth in this movement? An unceasing
struggle of thinking with untruth! The abstract is in its one-sidedness the
untrue, but as a moment in a corresponding concretion it always has its truth.
Truth is a transition from the one-sided to the all-sided truth; the movement
itself is the truth, or more exactly: a constant coming to the truth. And yet
neither is this the truth. The movement is true only so long as the full truth is
not reached, thus true only on account of a relative untruth. --- In order for
idealism to be absolute, the truth, the absolute content of absolute knowledge,
must be in the absolute Knower. With this presupposition idealism points out beyond
itself, for the absolute Knower is the absolutely Powerful: absolute knowledge is,
in the Idea of Truth, one with absolute power.

\emph{$\beta$) Realism.} The true and the untrue in realism are seen most clearly
from ``positive philosophy.'' Positivism presents in this respect three
characteristic principal aspects: it rejects every carried-through development of
concepts under the name of metaphysics; it degrades real causality to a merely
subjective condition for cognizing the laws; it denies ideal causality and breaks
the staff over all teleology. --- From the standpoint of facts, of the finite
realities, it is quite true that everything is relative, that absolute concepts do
not exist; but that the relative should therefore not have its ground in the
absolute, and that it should not lie in the interest of science to get the
conceptual relation between the relative and the absolute sufficiently illuminated,
is an arbitrary postulate, an untruth. When the concept of the relative is thought
through, then the absolute shows itself, and when the concept of the absolute is
thought through, then the relative shows itself; when the finite is analyzed, then
the infinite shows itself, and when the infinite is analyzed, the finite shows
itself.

Similar one-sidednesses come to light in the positivist treatment of causality. By
degrading the categories cause and effect to nominal values, subjective
designations for the connection of things, and depriving real causality of its
objective truth, one at the same time deprives the laws of all objective truth.
But if real causality has objective truth, then ideal causality must also have
original, subjective truth. That every thing acts according to its nature must
surely mean that it acts in consequence of its real-concept; the real-concept is
thus ideal cause. It is indeed true that the concept as concept is no acting cause:
the concept of gravity does not weigh, the concept of fire does not burn, the
concept of the stone strikes no dog dead. But the anti-logical existence of things
is an outburst of power from the logic of the concept. If no real-concepts were
realized in the existing things, then all objective possibilities, and with them
all acting and reciprocal action, would cease. In concept-less facticity the
things would exist beside one another, excluding one another, but for the rest not
have the least to do with one another. It is real-concepts that determine the
acting causes; but the system of real-concepts is again determined by the
all-comprehending, absolutely powerful self-unity, which in infinite determining,
comprehending, and acting is absolute causality.

\emph{$\gamma$) Dualism of Identity.} The philosophy of the absolute and that of
the relative, the representatives of idealism and realism, set themselves up as two
extremes, each on its side of the Idea of Truth; that they are both afflicted with
untruth in principle can be seen from the fact that the one must consistently
annihilate the other. Idealism presupposes that the universal is the objectively
all-comprehending that comprehends itself. When this proposition is made clear, a
mystification at once shows itself. If the proposition is to be taken strictly
objectively, then it must be the general being, filled with all that is --- in
other words: the universe --- that under the form of actuality comprehends itself.
That the real universe now expresses a system of concepts must be granted; but it
is one thing that the objective expresses a system of concepts, another that it
comprehends itself. The self-comprehending universal must necessarily be
apprehended as self-being, but such a self-being is not objective; it is, on the
contrary, subjective.

Realism makes the true immediately one with the actual; that is an untruth, for it
is a contradiction. Just as little as there can be anything objective that is not
objectified, just as little can there be anything in actuality true that a Knower
has not known. The actual, held fast for itself, is neither true nor untrue; truth
rests upon a judgment, and the truth of the judgment rests upon a syllogism. Truth
demands a dualism of Knowledge and Being, and in this dualism absolute identity.
The identity of the dualism is the self-unity that masters the extremes, the
syllogistic unity that from eternity is the absolute unity of beginning.

\section{Mathematics}
\label{sec:12}

\noindent Although mathematics and formal logic are closely related sciences, it
is yet a misjudgment of their peculiar characters when one wishes, as some have
attempted, to reduce logic to mathematics, or, what is still more unreasonable, to
deduce the being of mathematics from concepts of quantity in immanent logic. How
different the foundation of the two sciences is can be seen from the original
relation in which each of them sets the singular to the universal. However much
weight logic must, for the sake of concretion, lay upon the determination of the
peculiar relation of the single terms, it cannot, in the course of its development
of concepts, avoid letting the universal consume the singular, just as thought
consumes the representation, whereby everything in the last instance comes to rest
upon transparency in reflection. It is otherwise in mathematics. It begins with the
abstract universal, in so far as it begins by reducing all terms to magnitudes and
the relations to pure relations of magnitude; but it at the same time secures the
singular so punctually, articulates the determinations of relation so specifically,
that precisely thereby it acquires the character of an exact science. On this
account it cannot make do with the generality of reflection; it requires the
operation. For reflection spoken language is sufficient, for in the word the
general has the preponderance; but for the punctuality in the exact analysis the
word is not sufficient; mathematics must, for the sake of the operation, have its
own pasigraphic sign-language.

The relation between reflection and operation can be clearly seen in mathematical
problems. The condition for solving a scientific problem is above all to get it
sharply determined and rightly posed. Several problems remain standing in science so
obscure that one does not rightly know whether they are soluble or insoluble,
because they have not been apprehended with due sharpness and seen in full
illumination. The exact science can in this respect be said to go before with a good
example, when it trains its pupils to set up in an equation. The problem is in
general presented in the form of reflection; but on closer consideration it shows
itself that it must be posed once again, namely in the form of operation, in order
to be soluble: the problem consists essentially in getting the problem posed, namely
as a problem for the operation. The following problem, which gives two equations of
the first degree, may serve as an example. In the form of reflection it is presented
thus: ``When a number of $m$ oxen graze off $p$ acres of pasture in $n$ days,
likewise $m'$ oxen $p'$ acres in $n'$ days, how many oxen ($m''$) will then graze off
$p''$ acres in $n''$ days, the grass growing uniformly and having on all three
pastures the same height from the beginning?'' Everyone who understands the language
and knows that $m$, $n$, $p$, and so forth mean different numbers, understands the
problem so far, and can make his reflections upon it; but the problem is to pose the
problem, that is, to pose it in an equation.

From the relation between reflection and operation is understood the peculiar
emphasis the exact science lays upon proof. It can use the image, as in geometry; the
representation, as in arithmetic; reflection, as in the analysis of the infinite; it
has fixed fundamental principles, but it heeds no authority and takes nothing on
estimation: it demands proofs everywhere, and its proofs are incontrovertible. One
would suppose that a theorem such as the Pythagorean, so very ancient and proved so
many thousand times, from different standpoints in different ways, could at last be
exalted into a standing proposition of authority. Yes, were it merely a matter of
certainty, then this theorem is indeed raised above every doubt; but the exact
science still demands the proof, for the proof is not merely for the sake of
certainty, but essentially for the sake of grounding, of insight. Nor can anything be
assumed by mere inspection: that a curve parallel to a circle is itself a circle is
evident; but he who, by immediate estimation, assumes that a curve parallel to an
ellipse must necessarily be an ellipse, is remarkably mistaken. If one were to judge
by what one can simply see, then the asymptote, by continued prolongation, would have
to come to touch its curve; but the proof says no, and intuition must bow to thought.
That a straight line of equal length always signifies the same thing surely no one
will doubt --- except the mathematician. He alone knows how disparate the elements can
be of which a line consists; he knows that an ellipse or parabola can lie resolved in
the line; he knows that two straight lines intersecting each other can represent a
hyperbola.

The reflected unity of reflection and operation has its most general expression in the
mathematical function. Without any restriction whatever in the determination of the
general element in the mutual dependence of pure magnitudes upon one another, the
exact science can specify the relation of dependence so punctually that the
articulation goes into the infinite. In the equation $y = f(x) = ax$ let $a$ be the
constant, $y$ and $x$ the variable magnitudes. One sees at once from this what
difference there is between fixed and flowing boundaries. The constant magnitude is
arbitrary, it can be as great and as small as one will; but once it is posited, its
boundaries are fixed. The magnitudes $y$ and $x$, on the other hand, have boundaries
so movable that these cannot be fixed without the magnitudes in question changing
their character. From this one sees the difference between magnitude and number: every
number is a magnitude, but not every magnitude is a number: $x$ and $y$ are magnitudes,
$a$ is a number. The relation of dependence between $y$ and $x$ is reciprocal; when
$y$'s dependence on $x$ is denoted by $y = f(x) = ax$, then $x$'s dependence on $y$ can
be denoted by $x = \Phi(y) = \dfrac{y}{a}$. The notations $f$ and $\Phi$ concern not
the magnitudes themselves, but their manner of dependence, and in their manner of
dependence the law is seen.

By the help of the function the exact science becomes able to articulate the relation
between finite and infinite in such a way that the reciprocal determination becomes
evident. If $y = ax$, then also $dy = a\,dx$, $dy$ and $dx$ denoting elements,
vanishingly small parts, that is, differentials of $y$ and $x$. But from this it again
follows that one must get $\dfrac{dy}{dx} = a$. This result, simple as it is, is
infinitely fruitful. For it proves that a constant relation can be preserved even
though the terms of the relation dwindle into the infinite, and that the otherwise
undeterminable can be determined. According to the stated function $\dfrac{0}{0}$ must
be equal to $a$; but the same function must then also be able to give
$\dfrac{\infty}{\infty} = a$; the relation is preserved, in that the increment to $y$,
the dependent magnitude, is constantly measured by the increment to $x$, which here is
the independent one. If one measures $a$ by $dx$, then $\dfrac{a}{dx} = \infty$; if one
measures $dx$ by $0$, then also $\dfrac{dx}{0} = \infty$. The infinite seems so far to
be the immeasurable, the unattainable; but if we set out from the equation $dy =
a\,dx$, and denote the infinite summation by $\int$, the integral sign, then
$\int_{x=c}^{x=b} a\,dx = y = ab - ac$ gives a finite magnitude, which proves that the
infinite, the infinitely small, $dx$, as well as the infinitely great, the number
denoted by $\int$, falls within finite boundaries. In this again lies the possibility
of an articulated determination of the relation between discreteness and continuity.

Logically understood, discreteness is the form of finitude, continuity the form of
infinity. Existence is discrete, ``pure being'' is continuous; the atoms are in
discreteness, but space is a continuum; events are discrete, time is continuity. How
now do we come to the finitude of discreteness? By a breach of continuity, that is, by
a dividing-up. But how then do we come back to infinity? Not by a piecing-together;
for the piecing-together belongs to finitude just as well as the dividing-up. By a leap,
then? Quite so, when the difference between finite and infinite is held fast
qualitatively. But the exact science, which in everything qualitative looks to the
quantitative, knows yet another transition, approximation, the infinite approximation.
Separated from the concept of operation, reflection can very easily demonstrate the
contradiction in a concept such as infinite approximation. Approximation means that one
constantly comes nearer to the sought; thus $2$ is nearer to $100$ than $1$, $10$ still
nearer; but when the infinite is always infinitely far from the finite, then
approximation means nothing. An infinite approximation, an approximation to that which
one can yet never come nearer to, is an approximation that is yet not an approximation,
hence a contradiction.

The exact science masters the task, in that it knows how to give one and the same
magnitude both a finite and an infinite form, and thereby to prove in the operation that
the unattainable is attainable, and the attainable unattainable. In the equation
\[
  \frac{1}{1-x} = 1 + x + x^2 + \cdots x^n + \frac{x^{n+1}}{1-x}
\]
we have an example. If $x > 1$, then the series is divergent; by thinking of it as
continued into the infinite, one would come infinitely far from the goal; but if one
adds the remainder, the infinite removal is cancelled: one is at the goal. If $x < 1$,
then the series is convergent, and the approximation shows that, however much the series
is infinite, and the goal so far unattainable, the remainder yet becomes smaller and
smaller: the infinite approximation is thus after all a real approximation. If one gives
the number $2$ the forms $\dfrac{1}{1-\frac12}$ and $1 + \frac12 +
\left(\tfrac12\right)^2 + \cdots \left(\tfrac12\right)^n$, then it is indeed certain that
$\left(\tfrac12\right)^n$ is infinitely far from $0$ so long as $n$ is finite; but the
series has at the same time the property of showing us the transition from finite to
infinite discreteness. The smaller the difference between $\left(\tfrac12\right)^n$ and
$\left(\tfrac12\right)^{n+1}$ becomes, the more discreteness approaches continuity, and
in continuity the goal is reached, for $0 + 0 = 0$.

The infinite discreteness is thus the intermediate determination between continuity and
the dividing-up. For the sake of analysis the dividing-up, the finite discreteness, is
necessary; but mathematics cancels the dividing-up and gets the infinite discreteness
articulated in consequence of the function. It does not piece curved lines together out
of polygons, it cancels the dividing-up by an infinite summation of elements. If $y =
\sqrt{a^2 - x^2}$ is an equation for the circle, $a$ denoting the radius, and the point
of intersection of the coordinates lying in the center, then $dy = \dfrac{-x\,dx}{\sqrt{a^2
- x^2}}$ is an element of the periphery, and $\displaystyle\int_0^{2\pi} \frac{-x}{\sqrt{a^2
- x^2}}\,dx$ the whole circle-periphery as a flowing, continued curve.

So long as reflection was not yet methodically reconciled with operation, it was
constantly liable to be led astray. Many of the old thinkers' paradoxes, which look like
sophisms, are in reality, when we look more closely, scientific problems; thus the Eleatic
sophism Achilles. It was not dialectic that the Eleatics needed; a finite calculation could
not help them either. It was a distinct insight into the relation between continuity and
discreteness that they lacked; but the undisciplined understanding cannot grasp it, and so
it is at once said: the absolute, the infinite, is utterly incomprehensible to us; to think
about it is only idle brooding. But, remarkably enough, while one, in a realistic,
positivist manner, wishes absolutely to reject the absolute, one yet has scruples about
rejecting the infinite. This is highly inconsistent. If the finite cannot be understood
without the infinite, then the relative cannot be understood without the absolute either.
But the understanding is different. The analysis of the absolute is an analysis of ideas,
which belongs under logic; the analysis of the infinite is an analysis of operation, which
belongs under mathematics.

\section{Rational Mechanics}
\label{sec:13}

\noindent Although rational mechanics can in a way be regarded as a branch of
mathematics, it must in the present connection be brought forward as a particular
science, because it takes up into its foundation peculiar presuppositions, namely
matter and its forces. That it is nonetheless an \textit{a~priori} science lies in
the fact that it is not the real, but the represented --- represented matter,
represented forces --- upon which the emphasis essentially falls in the
investigation of the general laws. Force is the moving, matter that which is
moved; characteristic of matter is its inertia: it cannot, when at rest, set
itself in motion, and when in motion, cannot set itself at rest. As matter is
moved by forces, it is moved by outward causes, for it is always the force in the
one part that has the other for its point of attack. The manner in which rational
mechanics has secured for itself the conditions for the analysis of forces is
characteristic of its whole method.

Although the forces are invisible, they can yet, as the represented causes of
motions, be made visible by lines; a straight line, the represented track of the
motion, thus indicates both the strength and the direction of a single force. The
direction of the line corresponds to the direction of the force, and its length to
the strength of the force. For it is evident that of two forces which alternately
act on the same mass under the same conditions in equal time, that one will be the
stronger which causes the mass to traverse the longer way. If the forces $P$ and
$Q$ act on the same point, their result of strength, the resultant, $R$, will be
equal to their sum when they act in the same direction, and in the opposite case
equal to their difference: $R = P + Q$. If the directions are opposite, and $P =
Q$, then $R = 0$; there is thus equilibrium, and no motion. If the directions are
opposite, but $P > Q$, then $R = P - Q$ is not $0$; here is motion in the
direction of $P$. Cooperating forces with a common point of attack, whatever their
finite magnitude-relation may otherwise be, can form infinitely many different
angles with one another, two and two; it becomes therefore a task to bring the
countless cases under a comprehensive, universally valid form: this form is the
parallelogram of forces. If the angle the two forces $P$ and $Q$ form with each
other is denoted by $(P.Q)$, their resultant by $R$, then the equation $R =
\sqrt{P^2 + Q^2 + 2PQ\cos(P.Q)}$ holds for all possible cases. If the angle $(P.Q)
= 0^\circ$, then $\cos = 1$, and $R = P + Q$; if the angle $(P.Q) = 180^\circ$,
then $\cos = -1$, and $R = P - Q$. The diagonal of the parallelogram is the
resultant of the two component forces, and these its components; but each resultant
can again become a component in another parallelogram, and each component a
resultant of other components. The law of the parallelogram is a fundamental law
for the composition and resolution of mechanical forces, from which an infinite
system of particular laws allows itself to be developed.

In the parallelogram of forces rational mechanics gives us an intuitable image of
the unity of real-concept and law. But the exact science does not develop the
concept as concept; it must, for the sake of punctuality, hold fast the
represented, the terms of relation in their outward relations, and therefore
present the concepts in the form of laws.

The difference between law and real-concept rests upon the fact that a sum of
connected determinations can be apprehended partly in such a way that the terms are
represented each for itself, and their mutual dependence shows itself then in the
law; partly in such a way that the terms are transformed into reflexes in the
determinate totality, and the totality reflected in its reflexes shows itself then
as the concept. The quotient $\dfrac{y}{x} = \dfrac{dy}{dx} = a$ touched upon in
the previous section stands precisely at the transition; $y$ and $x$ are finite
terms, $dy$ and $dx$ are reflexes. The exact science must hold fast the
represented and comes in its analyses, first and last, to dwell upon the
articulation of relation, that is, upon the law. Even in the calculus of
variations, where maxima and minima depend not merely upon varying terms, but where
these terms themselves depend upon arbitrary functions, the terms cannot yet be
transformed into pure reflexes, but must, for the sake of the laws, be preserved in
the punctuality of representation. This punctuality steps forward especially in
rational mechanics.

We see it clearly in the reciprocal determination of force and cause. The cause with
its conditions belongs essentially under reflection and goes up into the
real-concept, but the forces with the represented masses and material terms fall
under the laws. Mechanics distinguishes between accelerating, moving, and living
force, and gets, by holding most strictly to the energy of matter, a series of laws
developed; the exact science divides up the concept and proves by this dividing-up
that it has its strength precisely in its one-sidedness.

\emph{Accelerating Force.} As the cause of motion the force can either act
suddenly, as in impact, and is then called instantaneous force, or continuously
change the velocity, and is then called, positively or negatively, accelerating.
But in the concept itself the difference is vanishing. The instantaneous force
requires a finite time, however small, in order to produce a finite velocity, and
is so far essentially accelerating; conversely, the accelerating force can be
resolved into infinitely many impacts and is in each impact to be regarded as
instantaneous. If we now ask about the cause, then the force in and for itself as
force is not the cause of the acceleration; it is so only by means of inertia. By
finite separation of the constant force, the inertia present in matter, and the
result of both, we get the exact expression, for example, of the law of falling
bodies. A constant force $k$ acts on a material point; in the relative moment $dt$,
$k\,dt = dv$, $t$ being the time, $v$ the velocity. The material point receives in
each moment an infinitely small, though constant, increment of velocity. That which
makes the increment of velocity, which is only $dv$, be $2\,dv$ in the next moment,
in the next again $3\,dv$, and so forth, so that in the $n$th moment it has grown
up to $n\,dv = v$ for $n = \infty$, is precisely the inertia, which preserves every
increment of velocity and thereby makes it possible that the velocity can grow with
time, be a function of time: $v = f(t) = kt$. But only when force and inertia are
seen as reflexes in the same acting do we get the concept cause.

\emph{Moving Force.} Here the force is apprehended as cause in relation to the
represented mass. If the mass $M$ is greater than the mass $m$, and $g$ the
acceleration, then the moving force in $Mg$ has undeniably a greater strength than
in $mg$; but while between $Mg$ and $mg$ there is only a quantitative difference,
since they both belong under the same representation of cause, there is on the other
hand a qualitative difference between $g$ and $mg$, for here the representation of
cause is changed. For the moving force as represented cause it is, strictly
speaking, not necessary that it cause a real motion. The cause is indeed not without
effect, but the motion may be cancelled by a counter-action, a pressure, and the
cause has nonetheless its reality. Thus with weight, a product of mass and
acceleration without perceptible motion: $Mg = P$. From the weight $P$ it is seen
that the mass has reality only in relation to gravity. If gravity changes its
intensity, then the mass gets another weight and, as the equation $M =
\dfrac{P}{g}$ shows, another reality. But however satisfactory the finite
dividing-up of the terms of relation may be for an exact determination of relations
and laws, just as unsatisfactory is it when one asks after the real-concepts
themselves and the essential causes.

\emph{Living Force.} By going from representation to representation, from premise to
premise, altering the given conditions and introducing new ones, mechanics can
systematically advance from simpler to more complicated relations, to new laws and
new formulas. If one represents the resistance of the moving force in the pressure
as cancelled, then the motion sets in, time gets its significance, and the motion
lays itself before the day in the product of mass and velocity, that is, in the
quantity of motion. If we call the moving force $mk$, the time $t$, and the velocity
$v$, then the equation $mkt = mv$, hence $mk = \dfrac{mv}{t}$, shows that any
constant force whatever can be measured by the quantity of motion it produces in a
unit of time. Nonetheless the intervention of time is inconvenient for the measure.
When two bodies have the same velocity, but the one has obtained it by a sudden
impact, the other by falling for a longer time, will one then, on account of the
difference of time, be able to say that their forces are different? That would, to
use an expression of Leibniz, be the same as to say that a man was richer because he
had taken a longer time to win his money. The force $mk$ can in the time $t$ awaken a
quantity of motion in $m$ that is $mV$, and in $M$ that is $Mv$: hence $mV = Mv$. If
we now ask what such a force can accomplish, the time being applied to move the mass,
or rather the weight, forward in space, then we come to the living force. If $mkt =
mV$, the activity in the first time-element will be $mk\,dt = m\,dV$. What now is
accomplished in the running time? The force moves the burden forward in space, but as
it drags on the burden, the velocity decreases, and the decreasing velocity is an
expression for the decreasing force. How long the force can hold out must then be
seen from the length of way, $s$, through which it is able to carry its burden. The
equation $m\,dV = mk\,dt$ must thus be transformed into $m\,dV = mk\dfrac{ds}{V}$, so
that it is known from the space what is accomplished in the time. The old saying that
time is equally long but not equally useful finds here a peculiar confirmation; $dt$,
the relative moment, is constant, but $V$ varies, whence it follows that the
space-element $ds$ also varies: $ds > ds^1 > ds^2 \ldots$ since $V > V^1 > V^2
\ldots$ From $m\,dV = mk\dfrac{ds}{V}$ follows $mV\,dV = mk\,ds$. By integrating this
equation between boundaries we obtain $mV^2 = 2mks$ and at the same time a measure
for the work. But when to the quantity of motion $mV$ there corresponds a living
force $mV^2$, then to the quantity of motion $Mv$ there corresponds a living force
$Mv^2$. Since now $mV^2$, in consequence of the presupposition $mV = Mv$, has
another value than $Mv^2$, it follows from this that the force in $mV$, when set to
work, accomplishes something other than the force in $Mv$. The living force, which is
measured by the work, is thus qualitatively different from the force moving in time,
which is measured by the quantity of motion.

If we now place ourselves at the standpoint of logic, that is, outside mechanics,
then we must recognize that the exact science, by its analyses, precisely procures
for us the requisites that belong to the real-concept when reflection is carried
through. The concept of cause is a real-concept, a concept of actuality, and a
concept of actuality always has reflexes of real things in its content. Reflexes of
things are reflexes of finitude, reflexes of dividing-up; but in the concept the
dividing-up must be cancelled. Reflection must therefore at once both posit and
cancel the represented. The dividing-up reflexes of the represented cause are: the
condition and the conditioned. Under these all determinations enter. Not even force
can in this light be held fast, for force for itself is only a representation; the
represented that we call force sinks down to being a condition. The conditioned, the
dependent, encloses at bottom the condition, the independent, and yet they set
themselves against each other as term against term. The reflection of conditions is
a reflection of represented terms, an outward reflection, hovering between essence
and existence, concept and actuality, characteristic of mechanical activity. Thus
force and inertia are conditions for acceleration; they are represented each for
itself, and yet point unceasingly to one another. Mass and acceleration are again
conditions for weight; the same mass can be subjected to another acceleration, and
the same acceleration applied to another mass; so far the conditions must be
separated. And yet they cannot be separated, for the altered product of mass and
acceleration alters the weight. If the weight is moved, time is drawn into the
circle of conditions, the velocity conditioned thereby becomes co-conditioning, and
the moving force shows itself as quantity of motion. This is again a condition for
the quantity of motion's being measurable by the space the resistance traverses.
Thus we have the living force; yet not yet as cause, for in the dividing-up reflexes
the cause is resolved: it is transformed into a sum of conditions.

The real-concept now forms itself logically by the fact that a sum of conditions is
reflected in one another; the dividing-up vanishes, and the totality-unity reflected
in its reflexes, the concept, shows itself as the determining element. This
reflection of infinity has tempted ``immanent logic'' to raise itself above
mathematics and imagine that it, without the intervention of rational mechanics,
could itself master the laws and the finite conditions. A regrettable
misunderstanding in science! Immanence-logic is by no means wrong in wishing to
cancel the dividing-up; its wrong consists, on the contrary, in this, that with all
its reflecting it does not sufficiently see through the double-sided essence of
reflection. In the cancellation of dividing-up all moments of the divided-up must,
be it well noted, clearly shine through; if this does not happen, then imagination
steps into the place of fact, that is, the speculation becomes imagined. One does
not reach the absolute by overlooking the relative; one does not comprehend the
unconditioned by a general resolution of all conditions; one must expressly be able
to recognize the relative terms in their reflexes, the finite conditions in the
resolving form: from this is seen the dependence of logic upon the exact sciences.

Rational mechanics, which has its mastery in the systematic treatment of law-bound
relations between the conditioned and its conditions, then also makes clear how
particular conditions can be referred to one fundamental condition, and this again
develop itself in a manifold of single conditions. In order for forces acting on a
point to be in equilibrium, their result must be zero; this is the general condition
of equilibrium, its fundamental condition. But what a swarm of particular conditions
and varying single conditions can enter into this! The equilibrium of forces is the
cause of rest; the cause of motion is the cancellation of equilibrium: it is at
bottom the same cause and yet, under the changing of conditions, a different cause.
If we reflect on the difference, then statics and dynamics each have their circle of
tasks, and so far fall apart from each other; if we reflect on the unity, then the
static task becomes dynamic, and the dynamic static. The reciprocal determination is
grounded upon the reflection of two corresponding principles, the principle of
virtual velocities and D'Alembert's principle for the equilibrium of lost forces.

The deep gulf between the logical in the general real-concepts and the anti-logical
in the world of things can indeed no human knowledge fill, but rational mechanics
can build a bridge over it.

\chapter{Real-Empirical Sciences}

\section{Mechanics and Chemical Physics}
\label{sec:14}

\noindent The concepts realized in nature are more than concepts; they are existing
things; here physics begins, and logic can, as formal logic, especially as the logic
of induction, merely serve as a guide. The things are indeed subject to general
laws; but in their real fulfillment the laws of nature deviate from abstract laws;
the planets do not, after all, move in ellipses, the ballistic curve is no parabola,
and the law for the physical pendulum has a different appearance than that for the
mathematical. Rational mechanics treads a foreign ground and becomes an auxiliary
science for mechanical physics. Here is a new method, because here is a new
foundation, that is, the foundation of outward experiences. Everyday experiences are
too incoherent and superficial. The sensible thing with its properties stands in such
many-sided and disparate relations to other things with corresponding properties that
a peculiar art is needed to loosen the knot of relations. What, amid the
complications, must from a chosen standpoint be held fast as essential may, when the
standpoint is changed, be regarded as something unessential, indeed as something
contingent that intervenes disturbingly in the understanding of what matters. The
analysis is more than logical, it is real; the disturbing element cannot be removed by
mere abstraction, it must, as far as possible, be removed by experiment. The task of
research is to direct each question to real nature in so exact a form that the fact
itself must necessarily answer either yes or no.

As a shining example of reciprocal action between mechanical physics and rational
mechanics stands Newton's discovery of the general law of gravitation. The eye that
saw in the falling body an expression of the force that also draws the moon down
toward the earth belongs not to sensation alone, but also to thought. It was the
thinker's intimation of a general essence inherent in all bodies that gave the impulse
to the great discovery. That the moon's fall must be related to the apple's as the
force that draws the moon to the force that draws the apple is a matter of course; but
in what relation do these forces stand to one another? The force that draws the apple
is known by the law of falling bodies, $15\tfrac58$ feet in the first second; the
force that draws the moon can, since the moon's mean distance and mean period of
revolution are known magnitudes, likewise be determined. If the earth's distance from
the moon's center is denoted by $R$, the angle corresponding to the arc (in one second
of time) by $\Theta$, then $\rho$, the segment that denotes the moon's fall toward the
earth, is expressed by the equation $\rho = 2R\sin^2\tfrac12\Theta$. The relation of
the forces is thus determined by the quotient
$\dfrac{2R\sin^2\tfrac12\Theta}{15\tfrac58}$; but, in order for $R$ to be determined,
the earth's radius $r$ must, be it noted, also be known. Here is a circle of
presupposed observations. And yet Newton could not, without special preparation, get
it enriched by new observations. It lay near at hand, to be sure, to seek the ground
of the difference of attraction in the difference of distance; but with this general
surmise no experience could yet be made; only by the assumption that the attracting
forces stand in inverse relation to the squares of the distance was the equation
determined:
\[
  \frac{2R\sin^2\tfrac12\Theta}{15\tfrac58} = \frac{\dfrac{\mu}{R^2}}{\dfrac{\mu}{r^2}} = \frac{r^2}{R^2}.
\]
Is the relation $\dfrac{\mu}{r^2}$ the right one? That was Newton's great question. The
answer let itself be awaited, but it did not fail to come.

That the same fundamental mechanical laws that hold for the globes in the heavens must
also hold for the bodies on earth is in science an incontrovertible presupposition;
but since the fulfillment varies with the conditions, since the bodies on earth lie far
nearer to us and can be made an object of observations, experiments, and
investigations in a quite different manner than those distant bodies, it goes without
saying that particular cases with particular tasks and specific conditions here cross
and interrupt the systematic continuity otherwise than in the observation of the
heavenly bodies. No wonder that astronomy, on account of its more abstract character,
though not yet completed, yet stands with a stamp of completion that distinguishes it
above every other empirical science. How intimately the rational can fuse with the
empirical, Laplace has proved in his \textit{Mécanique céleste}.

In mechanical physics the active terms --- atoms, molecules, masses --- preserve, in
their mutual reciprocal action, a certain outward independence, so that the changes
essentially rest upon a change of place in space. Empty space is between them, in so
far as they do not touch one another immediately. They all have the common property of
being heavy. The force of gravity acts at a distance, and the distance is the empty.
Such is the impression when we one-sidedly hold to the merely mechanical. But
penetrating investigations concerning the remarkable peculiarities of light, heat,
magnetism, and electricity have led to the recognition that the opposition between the
heavy masses and the empty space is not sufficient to explain the phenomena, that
besides the ponderable matter there must be an imponderable one, and that there are
attracting forces different from gravity. Physics then divides, with respect to this,
into mechanical and chemical physics.

Since the imponderable matter cannot be an object of immediate experience, its
assumption is naturally only a hypothesis. But with the hypotheses of physics it
stands in a peculiar way; we see it clearly from the two well-known hypotheses about
light, the corpuscular hypothesis and the vibration hypothesis.

Newton assumed that the luminous substance consisted of small bodies which the
light-givers cast out with an extraordinarily great velocity, and which, by striking
the retina, produced the impression of light. Huygens, on the other hand, assumed that
the universe was filled with a substance that unites with an extraordinarily low
density an extraordinarily high degree of elasticity. This substance is the ether,
which is set in wave-motion by the luminous bodies; it is ether-waves which, when they
strike the retina, produce the impression of light. Neither light-corpuscle nor
ether-wave is affected by gravity. --- Were the dispute about such hypotheses to be
decided by general grounds of reasoning, then it would look ominous for physics. Those
who attack the theory of light-corpuscles object that one would every moment have to
get stray little bodies into the eye and see flashes from all sides. Those who espouse
the theory answer that this difficulty is easily avoided. For when a luminous point
sent 100 light-corpuscles into the eye each second, the eye would receive a steady
impression of light; now light traverses 40,000 miles in a second, so there would be
400 miles between two successive light-corpuscles, so that there is no danger of
collision. The ether-waves, on the other hand, would in their opinion have to be
weakened on the way, since the ether cannot possibly be perfectly elastic. The
opponents of the corpuscle-theory can again object that the impact of the little
bodies would have to knock the eye out. A light-corpuscle of a mass like
$\frac{1}{2\,000\,000\,000}$ grain would give an impact like a ball of 25 grains'
weight shot from a rifle, and so forth. Such a dispute for and against would become
just as endless as a dispute about the condition in the other world. But here it shows
itself what difference it makes to dispute about a supernatural whose connection with
the natural is uncertain, and to dispute about something non-sensible that yet belongs
under the sensible, and whose significance must therefore be testable on facts. Thus
there came at last a turning-point that became decisive for Newton's hypothesis. For
in consequence of the hypothesis light must go $\tfrac43$ times faster in water than
in air; since it is now established by experiment that the relation is precisely the
reverse, it follows from this that the corpuscular hypothesis must be given up. The
vibration hypothesis, on the other hand, according to which the strength of light
rests upon the amplitude, the color upon the period of vibration, has hitherto shown
itself satisfactory in the explanation both of the propagation of light and of
particular phenomena corresponding to reflection, refraction, interference,
polarization, and so forth.

In the logical sense thing and concept are commensurable; but in the empirical sense
they are incommensurable, for the thing is anti-logical. From this it follows that not
merely hypothesis, but also analogy, must in the application of induction play a
prominent role in physical observations. Analogy can guide; for two phenomena that are
alike in several respects awaken the surmise of likeness in all respects; but analogy
can also mislead, for under the apparent likeness an essential unlikeness may hide,
and conversely. In the doctrine of the relation of heat to light we have an example. It
is a well-known matter that heat can spread in rays just like light, that the
heat-rays are refracted, and refracted differently like the light-rays, that they are
reflected according to the same law: angle of reflection $=$ angle of incidence, and so
forth. But heat is not light; the unlikeness comes to light, among other things, in
this, that dark bodies, bodies that do not radiate light, can radiate heat. Analogy,
however, takes up the likeness in the unlikeness, and the hypothetical explanation is
that the mere heat-rays, though they may pass through the watery medium of the eye, yet
cannot affect the optic nerve, it being assumed that the waves are too broad and the
vibrations too slow. The unlikeness grows when we remember that heat spreads, not only
by radiation, but also by conduction, in that by immediate contact it passes over from
the one body to the other and spreads through the whole mass. To the highly
characteristic difference between good and bad conductors of heat there corresponds no
such difference between good and bad conductors of light. But the likeness again
conquers the unlikeness; for an explanation is possible when one assumes that the
heat-waves, the vibrating ether-parts that produce heat in the body, at the same time
bring the body's own minute parts into vibration. The assumption finds confirmation in
the doctrine of specific heat and atomic heat. The thorough testing of an analogy here
coincides with the testing of a hypothesis.

In physics law is the highest instance, the last ground to which everything must be
referred; what then is the ground that the bodies attract one another at a distance?
Gravity. But what is gravity? A law! What is the cause of light and heat? Ether-activity.
But what is the ground of the ether-activity? A law! What is the cause of the polar
relations in the magnetic and electric phenomena? Magnetism and electricity. But what
is the ground of magnetism and electricity? A law! From this standpoint there shows
itself a remarkable agreement and at the same time a striking opposition between physics
and rational mechanics. Both point unanimously to law as foundation, both have the
hypothetical judgment: if A is B, then C must be D, for their guiding thread. But the
use of the given guiding thread must become different, according as cognition goes in
an \textit{a~priori} or empirical direction. Understood \textit{a~priori}, the terms of
relation are themselves moments in the relation; one can be assured that in the
hypothetical assumption of the premise there is contained not the least bit more than
what the logical consequence of the conclusion states. If it is assumed, as Aristotle
teaches, that immediate contact is an inviolable condition for the bodies' acting upon
one another, then the rejection of Newton's ``Principia'' must be an irresistible
consequence thereof. In the domain of physics, on the other hand, such argumentations
have no validity. And why not? Because the relation between term and relation is in the
empirical quite other than in the \textit{a~priori}. In physics, especially in
mechanical and chemical physics, the terms have such independence that the relations by
which the laws are known must be secured by careful induction before the hypothetical
judgment can come about, let alone lead to a reliable conclusion. If the stone falls
every time it is deprived of its support, then --- yes, then what? Then it must be
because it is passive and sluggish and cannot bear itself. No, then it must be because
it belongs to the bodies whose nature it is to seek downward toward the center, toward
the world's eternal center. No, then it must be because, in reciprocal action with the
earth, it is attracted by the earth. If, already in order to secure the premise, a
system of inductions must be presupposed, then it is understood from this that only
careful research can lead to a reliable conclusion. The apparently analogous cases are
often so different that the conclusion may deceive, whether one applies \textit{modus
ponens} or \textit{modus tollens}. The incomplete induction is namely a conclusion of
probability, in which the probability that the predicate holds for the general case
(the genus) grows the greater the number of particular cases (species) is for which it
holds. Since such an induction may, however, have its ``instances,'' the laws found by
induction may also have their exceptions; thus the law that heat expands bodies has an
exception; thus Mariotte's law has an exception, and so forth. It is indeed at odds
with law as law to have exceptions; but what is at odds in this stems from the fact
that the empirical expression for the law is on the whole imperfect.

It is a matter of course that the cognition of particular laws must be the more perfect
the more one, by conditions and intermediate determinations, is able to derive them
from, or refer them under, a corresponding fundamental law; whereas they take on a
merely empirical character when no such deduction or reduction is possible. By
reduction and deduction the particular laws become rendered \textit{a~priori}. The more
general --- that is, simpler --- the condition is which the fundamental law entails, the
more easily it goes with the rendering-\textit{a~priori} of the particular laws; the more
involved and hidden the conditions are which the fundamental law itself contains, the
more difficult become the tasks of research. How clear and simple is the comprehensive
fundamental law of mechanical physics, the law of gravitation, in comparison with the
specific fundamental laws of chemical physics, especially those of magnetism and
electricity! While the former rests upon a single property common to all ponderable
matter, the latter are bound to a complex of derived properties and hidden conditions
that allow themselves to be seen through only with difficulty. The laws of mechanical
physics have then also been rendered \textit{a~priori} by a quite different measure than
the laws of chemical physics.

What this means can be seen from the relation between Kepler's and Newton's different
standpoints. The Keplerian reflection abstracts the laws from cases at hand and secures
them by empirical induction; the Newtonian method, on the other hand, sets out to
cognize the causes from their effects, to determine the particular on the foundation of
the general, and to comprehend the laws in their fulfillment. On the way of
objectification a grand advance! Kepler speaks only of the areas described in the
planetary orbits; but Newton proves that the law for the areas must hold for any motion
whatever about a fixed center, since it is in general determinative for the objective
concept of central motion. Kepler holds immediately to the factual, that the planetary
orbits are ellipses, but, bound to immediate facts, he is not able to state the cause;
from Newton's principles, on the other hand, the cause can be seen in the effect, in so
far as it can be seen that not merely the ellipse, but in general the conic section, is a
naturally necessary consequence, an express causal consequence, of the peculiar law of
attraction $T = \dfrac{\mu}{r^2}$, and that conversely this law is a necessary
consequence of motion in a conic section. Kepler confines himself to setting up the
three laws beside one another without any deduction whatever; but according to Newton's
principles the one law can, when one sets out from the general and introduces
conditions, be deduced from the other. That the deduction must constantly be attached to
given conditions reminds us that the proofs, as regards the physical, are neither merely
\textit{a~priori} nor merely empirical, but that they, in the nature of the case, rest
upon an exactly articulated reciprocal determination of experience and thought. It is
not a speculative constructing \textit{a~priori}; it is a thorough scientific
rendering-\textit{a~priori} of which we here speak.

It is otherwise in the doctrine of magnetism and electricity. Even if one, as has
earlier happened, set out from the presupposition that some bodies are magnetic, others
electric, one would yet not possibly be able, by setting out from the fundamental
phenomena corresponding to magnetic and electric activity, to apprehend these in such a
way that from them, by inserting conditions, one could derive and explain the particular
phenomena. What is magnetism? An expression of part-dividing, polar functions. An
attraction according to the law of gravity presupposes at least two terms, two masses or
parts of mass; nonetheless there is in this no proper part-division, the one mass is
like the other, and every difference in properties, even in size, is, generally
speaking, foreign to the essence of gravity. It is otherwise with the attraction between
two magnets; each of them has its poles, the unlike poles attract, the like ones repel
each other; attraction is conditioned by repulsion and conversely. In the relation of
the poles the one mass is not merely outside the other, but the otherness is accented
twice, the one pole is qualitatively opposite to the other; the south pole has its
definite other in the north pole, and precisely therefore the attraction is mutual; the
south pole, on the other hand, has only a formal other in another south pole; the
uniformity is at odds with the principle of part-division, and the repulsion mutual. How
originally the opposite terms presuppose each other is seen in fact from this, that one
can indeed procure magnets with many poles, but not a magnet that has only one pole.
What is electricity? The analogy with magnetism is striking on account of the polar
opposition; but as regards the inner active element, the likeness must rather seem to
belong to the phenomenal than to the essence. For magnetism is, as ``force'' and
``property,'' bound to the magnetic body, while electricity ``can freely pass over from
the one body to the other,'' and by this passing-over shows itself as analogous with
heat. There is, as regards conduction, so remarkable a connection between electricity
and heat that one has found, for the metals as well as for their alloys, the same
relations between their conductivity for heat as for their electrical conductivity. If
heat had poles, it would lie near at hand to apprehend electricity as polarized heat;
but since no fact points to there being a positive heat-pole that attracted a negative
and repelled another positive, since by bound and freed heat is understood something
quite other than by bound and free electricity, since heat, in other words, is only
quantitatively, not qualitatively, different in itself, it is seen from this that there
must be a fundamental difference between electricity and heat.

To this comes the fact that the magnetic and electric properties are communicable and
variable, so that bodies under different conditions can become now magnetic, now
diamagnetic, now again non-magnetic, or now positively, now negatively electric, now
again non-electric. These are remarkable insights that are hereby opened to us into the
molecular changes. Immediately seen, it presents itself as though the magnetic force sat
in each of the magnet's poles, while in the so-called equilibrium-region there was no
such force; but the part-division is pervasive. By breaking a bar-magnet in the middle,
one gets not two half magnets, each with its pole, but two whole ones, each with two
poles. The part is a complete whole, and the division can be continued down to the
elementary singularities: every magnet is a totality of elementary magnets. The ground
that the active poles form themselves toward the end-faces lies in this, that the
elementary magnets in the equilibrium-region equalize and cancel each other's effects,
whereby there arises a peculiar opposition, conditioned by the circumstances, between
bound and free magnetism. --- How deeply electricity is implicated in the peculiarity of
the different bodies can already be seen from the manner in which it is awakened. It is
awakened by friction, but the friction precisely makes the bodily differences
recognizable; it is a characteristic expression of the difference between glass and
sealing-wax that the two bodies, by being rubbed with the same substance, wool or silk,
become, the first positive, the latter negative. It is awakened by distribution. If a
body A is charged, for example, with positive electricity, it will through an insulating
interval act on another body B in such a way that by distribution in B there arises
negative electricity on the parts of the surface nearest to A (induced electricity of
the first kind), but on the surface farthest from A positive electricity (induced
electricity of the second kind). To this corresponds again the stronger or weaker
attraction, for since the distribution is greatest in the good conductors, the
attraction to these must also be strongest. The peculiar difference, determined by
electricity, between good, half-good, and bad conductors, insulators, which is so
characteristic of part-division, heat will not merely alter, but at higher degrees even
efface; one knows no body that remains an insulator at white heat. But when heat acts not
inhibitingly but conditioningly, the part-division is so marked that even mere contact is
enough to awaken the electricity.

That the specific phenomena of magnetism and electricity are sheer fulfillments of laws,
of this no one can doubt; but the relation between the individual and the general is so
articulated and therefore so involved that the researcher, with all his art, every
moment loses the thread in the great labyrinth. On what does the reciprocal action
between the electric and the magnetic in electromagnetism properly rest? Ørsted's
discovery sprang forth from a fact like a blaze in the dark; but where is the ground of
the phenomenon, in what conditions lies the active element itself? What is the flowing
element in the electric wire? Why will the north end of the magnetic needle make a
deflection toward the east when it is situated above a current from south to north, and
on the other hand a deflection toward the west when it is itself situated below the
current? What is the ground that two parallel current-conductors attract each other when
the currents go in the same direction, but repel each other when they go in opposite
directions? Ampère's theory, that every magnet can be regarded as a system of electric
current-circuits, solenoids, and magnetism thus regarded as a peculiar phenomenon of
electricity, seems, remarkably enough, to fit the earth's magnetism, but not the bodily
magnets, since it cannot be assumed that electric currents should be able to maintain
themselves, even in good conductors, without either giving off or receiving any work.

Since a rendering-\textit{a~priori} can, in the nature of the case, advance far more
rapidly in mechanical physics than in chemical, it might at first glance look as though
the system of real-concepts in the former lay far nearer to the realities than in the
latter; but the relation is precisely the reverse. The laws of mechanical physics have a
far more abstract stamp than is the case with the laws of chemical physics. The more
deeply the laws themselves enter into the concretions --- that is, the more richly the
specific quality of the terms and the specification of the relations reflect each other
--- the more clearly it will show itself that the laws themselves are the
universally-necessary expression of the real-concepts. When it nonetheless looks as
though the physical-mechanical were in itself more comprehensible than the
physical-chemical, then the ground is simply this, that the physical concretions are
anti-logical, and that the anti-logical must always become the more prominent the more
involved the concretions are, while empirical cognition, in so far as it is in truth
scientific, however great its interest in facts may be, is always referred to the more
abstract and thereby to the logical.

\section{Chemistry and Mineralogy}
\label{sec:15}

\noindent In all the phenomena of part-division and polar opposition hitherto
considered, the bodily terms yet still keep outside one another: the energies are
exchanged, tension is transposed into living force, and living force into tension,
heat into electricity, and electricity into work; but under all this the bodies in
question yet maintain their relative independence. Only in chemism does the
part-dividing activity become so pervasive that the terms surrender their
independence in their affinity: the concrete-bodily existence is resolved in order
to be produced, and produced in order again to be resolved. Just as mechanical
physics forms the transition to chemical physics, so chemical physics in turn forms
the transition to chemistry.

\emph{Chemistry.} What is in a voltaic pile the originally acting cause? Is it the
electric current that is the cause of the chemical activity, or is it conversely the
chemical activity that is the cause of the current, or is there here perhaps a
peculiar, equal cooperation of both factors? The last is indeed, on science's present
stage, the most probable. By the voltaic pile water has been separated. Like water,
all composite substances are separated, provided, be it noted, that they conduct
electricity and are at least in part in a melted or dissolved state. The electric law
says that the weights of the substances separated by the same quantity of electricity
are related as these substances' equivalent numbers. As physics strives from all
sides to give an account of the molecular states of bodies and the possible changes
these undergo, one has explained to oneself the separation of water's constituents,
which can indeed each be collected and measured in a voltameter, in the following
ingenious manner: When hydrogen is combined with oxygen into water, then by the
intimate contact between the minute parts the oxygen-atoms become negatively and the
hydrogen-atoms positively electric; but on account of the uniform distribution of the
minute parts of oxygen and hydrogen the combination shows no free electricity. There
now lies --- one may imagine --- in a given line a certain number of water-parts
between the two poles of a galvanic chain. Each water-part consists of one atom of
oxygen and a double-atom of hydrogen. The positive pole acts on the first water-part
in such a way that the oxygen-atom, which is negative, is attracted, the
hydrogen-atom, which is positive, on the other hand repelled. But the first
water-part acts in turn on the second, this on the third, and so forth. When now the
attraction that the positive pole exercises on the oxygen-atom in the first
water-part is sufficiently great, this is torn away from its hydrogen-atom; but the
separated hydrogen-atom combines immediately thereupon with the oxygen in the second
water-part; the hydrogen-atom thereby separated out further draws the oxygen from the
third water-part, and so forth. Throughout the whole stretch there takes place a
continuous separation and composition of the water; only at the poles themselves can
its constituents become free.

The inner connection between electricity and chemistry is clearly shown in the
electrochemical theory. Just as zinc and copper in contact with each other acquire
opposite electricities, so, according to the theory, the atoms of two elements
acquire opposite electricities. That the atoms are not in and for themselves electric,
but first become so by mutual contact, is a presupposition of representation that
explains the fact that one and the same substance can, according as the corresponding
term of the combination is, be now positive, now negative. Thus sulfur in combination
with oxygen forms the electropositive, but in combination with hydrogen the
electronegative term. By the mutual contact of the fundamental parts the electricity
becomes bound, so that the body is everywhere in its inner composition equally both
positive and negative; but the part-division repeats itself in ever more concrete
forms. When the body produced by the combination is, under the requisite conditions,
brought into chemical relation to other bodies, the neutrality ceases. The elements
combine in pairs into binary combinations, combinations of the first order; the binary
combinations again combine in pairs into combinations of the second order and form
salts, in which the one of the combinations, the acid, is the electronegative, the
other, the base, the electropositive term; salts are again combined into double salts,
and so forth. The affinities have different strengths.

It is, however, not enough to speak in general of stronger and weaker affinities; here
a measure is required for the absolute magnitude of the affinities: this measure is
the development of heat. Thus we come to the thermochemical theory. In the same
chemical process an equally great quantity of heat is always developed when the outer
conditions are the same. The development of heat that accompanies the formation is
always equally great, whether the combination or the separation happens at once or
gradually. By the separation of a chemical combination there is always produced just
as much heat as is developed by the formation of the combination. When one wishes to
accomplish a work, the quantity of work must be greater than, or at least as great as,
the work. If one wishes to separate a chemical combination by the quantity of work,
here quantity of heat, that the formation of another chemical combination is found to
require, then this quantity of work must be greater than that which is demanded for
the separation of the first combination.

More closely viewed, it is yet not enough either to have a measure for the affinities;
for what must in truth awaken astonishment --- the strength of affinity by no means
corresponds directly to the relation in which the substances are saturated. Chlorine
and hydrogen have a great affinity to each other, chlorine and carbon a comparatively
slight one. Nonetheless one hydrogen-atom is saturated with one chlorine-atom, while
one carbon-atom is fully saturated only when it has combined with four chlorine-atoms.
So there arises the doctrine of atomicity, the doctrine of ``valence,''
``affinivalence,'' ``quantivalence''; whence it is seen that the method of chemistry
must ramify into different sub-methods more or less dependent on one another. Since one
atom of chlorine combines with one atom of hydrogen into hydrochloric-acid gas, one
atom of oxygen with two atoms of hydrogen into water, one atom of nitrogen with three
atoms of hydrogen into ammonia, one atom of carbon with four atoms of hydrogen into
light carburetted hydrogen, then, accordingly, the chlorine as well as the hydrogen is
called one-atomed (monovalent), the oxygen two-atomed (divalent), the nitrogen
three-atomed (trivalent), and the carbon four-atomed (tetravalent). In the combination
$\mathrm{CH}$ the carbon thus has three, in $\mathrm{CH_2}$ two, in $\mathrm{CH_3}$ one
atomicity free; this casts a light on the isomeric and polymeric formations that are so
frequent among the carbon-combinations.

The chemical processes are body-forming activities, both with respect to separation
and to composition processes of concrete bodiliness. The analysis of concretion must
therefore take up into itself so great a manifold of disparate elementary analyses that
one cannot possibly, with Comte, regard chemistry as fundamental in opposition to
mineralogy. Mineralogy and chemistry mutually condition each other, even though we must
grant that the latter has a broader basis than the former. No body-formation is, in the
physical sense, to be apprehended as a single effect of a single cause, but every body
arises, be it as educt or as product, from the cooperation of different causes. As it
goes with the formations, so also with the transformations. In the reciprocal action of
cohesion and heat, by which the states of aggregation are determined, the mechanical,
the dynamic, the magnetic, electric, and electrochemical can in the most varied ways
take part.

\emph{Mineralogy.} While the chemist looks chiefly to the body-forming activity, the
mineralogist begins immediately with the product of the formation, whereby at the same
time the state of aggregation and the crystal form conditioned by it come into
consideration. The crystal form is a physical result of a body's transition from
fluid to solid state, the minute parts being, in consequence of the body's chemical
character, deposited in a definite manner. That in this respect there can be no
question of a peculiar principle of crystallization is seen from the fact that at the
transition an amorphous or merely crystalline formation can just as well arise as a
fully formed crystal; in this respect the conditions reign. In order to bring bodies to
crystallize, one can employ different procedures --- dissolving, for example, in water,
melting and cooling, evaporation and condensation of vapors, according to the
constitution of the bodies in question; but the influence of the conditions is on the
whole known by the rule that the crystals become the larger and more beautiful the
more slowly and quietly the crystallization takes place.

Since the crystal form is formed of plane faces, such mathematical solids as the cube,
the square octahedron, the rhombohedron, and so forth, can be chosen as fundamental
forms; but what matters especially in the formation of a system is the number,
position, and size of the axes; how far in a group there are, for example, three
equally large axes, each single one perpendicular to the plane of the others; or only
two horizontal axes equal, the vertical, the principal axis, on the other hand greater
or smaller; or none of the three axes equal, and so forth. As regards the physical, two
disparate substances can crystallize in the same form (isomorphism) and one and the
same substance in different forms (dimorphism, polymorphism).

If by molecular constitution we understand the manner in which the minute parts in a
body are bounded, deposited, and combined, then one will also understand that here
there must be a peculiar connection between atomic constitution and crystal form. As
regards the single atoms corresponding to the elements, polymorphism may arise from the
fact that the same atoms, under several different crystallizations, deposit themselves
differently; on the other hand, isomorphism allows itself to be explained partly --- as
is sometimes one-sidedly emphasized --- simply from the fact that the atoms in two or
more different elements are equally large and equally formed, partly from the fact that
the atoms, though unequally large, yet in figure and size are so related that a group of
the smaller atoms can be thought congruent with the single larger one. As regards next
the atomic compositions corresponding to composite bodies, the number of possibilities
is the greater the more alternatives the compositions in question contain. In order for
two bodies to be able, in the proper sense of the word, to crystallize together, it
must indeed be presupposed that they have the same crystal form; but likeness in
crystal form is yet not without further ado a ground for their crystallizing together:
for this is required, as a rule, likeness in chemical composition. A crystal of sal
ammoniac, for example, has the same geometrical form as a crystal of alum; nonetheless
they crystallize separated from each other in one and the same liquid; but if we look
at the constituents, then the sal-ammoniac molecule has after all only two, the alum on
the other hand thirty composite atoms.

If one now reckons among minerals not merely the inorganic substances consisting of
homogeneous parts that form the solid part of the earth-globe, but even several
substances of organic origin occurring in the earth, such as coal, brown coal, amber,
and infusorial silica, then one can with reason say that mineralogy and chemistry are
sciences of concretion in a peculiar sense, sciences that take up the results of the
foregoing and unite everything into a penetrating cognition of concrete bodiliness. As
a mineral the body has, besides its chemical and crystallographic distinguishing mark,
both mechanical, optical, thermal, and electrical properties; a mineral system that
merely holds to outer, ``natural-historical'' marks does not satisfy; a system that
takes account only of the mineral's chemical composition is also one-sided: here, as in
the ``mixed'' system, account must everywhere be taken of the concrete.

As regards the scientific worth of the method of concretion, two one-sided ideals
present themselves: according to the one, science aims simply at essential insight;
according to the other, it aims chiefly at the acquisition of knowledge and
acquaintance with real things. In the first respect one is tempted to set the logic of
absolute idealism higher than astronomy, which with its more abstract terms of relation
and more general cognition of law is yet empirical, indeed far higher than chemistry and
mineralogy, which with their many special investigations are quite lost in the empirical;
in the latter respect one will reject the doctrine of ideas, which does not enrich us
with knowledge, and set mineralogy, which with ever deeper penetration into the inner
structure of bodies procures for us much to know about their real constitution, not
merely above logic, but even above astronomy and general mechanics. Such ways of valuing
are, however, vague and arbitrary. The only true measure in the judgment of the
scientific in science must, as shown above, be sought in the relation between the
concept, the law, and the thing, or concrete bodiliness.

If it were possible to derive the laws of existence from pure concepts, from the laws
the general properties of existing things, and from these again real existence, then the
doctrine of ideas would indeed take up the other sciences into itself; but it cannot be
done. If it were possible that sensible observation and empirical induction could by
themselves lead to essential insight into the ground of things, then we should have to
grant those right who, with disdain for deeper-going analyses of concepts, hold
exclusively to facts; but it cannot be done. Only from the three starting-points, from
the real-concept, the law, and concrete bodiliness, can science work its way forward to
a grounded cognition of real truth. A development of concepts which neither has roots in
the real nor aims at definite cognition of actuality is, despite all virtuosity of
reflection, barren and unfruitful, and an empirical art of observation with endless
analyses lacks true scientific consciousness. Positivism is in its way just as one-sided
as objective absolute idealism in its; from both sides one carries on smuggling with
contraband wares: it is a kind of scientific customs-fraud. In the clear recognition of
the reciprocal relation between the logical in the analysis of real-concepts and the
empirical, the anti-logical, in the analyses of concrete bodiliness lies the scientific.

\section{Geognosy and Geology}
\label{sec:16}

\noindent As a whole with its parts the earth-body presents two different standpoints
for scientific apprehension. In so far as it is a matter of investigating the parts of
which the whole is composed, the standpoint is geognostic; in so far as one aims at a
presentation of the development of the whole in and with that of the parts, the
standpoint is geological. Even if one will not set geognosy and geology against each
other as two separate sciences, it is yet unmistakable that in the geognostic decisive
weight is laid upon actual research, while the geological, reminiscent of the geogonic,
falls chiefly under hypotheses and theories. The method of geognosy is analytic, that
of geology synthetic.

If we comprehend all parts of the doctrine of nature that belong to the investigation
of the shape and composition of the earth-body under the common name geognosy, then
geognosy, with astronomy, physical geography, chemical physics, chemistry, mineralogy,
and lithology, is a science so widely ramified that it may well be said to comprehend
the whole earthly-elementary nature. Can this science then, in the proper sense, teach
us what nature is? No! It investigates the things of nature, but not nature. It does
indeed teach that every peculiar body has its peculiar nature, and that the activity
under which the bodies are formed, grouped, changed, is a pervasive natural activity;
but nature itself in the disparate natures it does not investigate, and that for the
good reason that nature itself is neither a thing for us with properties nor a thing in
itself without properties, but a real-concept, or rather: the system of real-concepts
that, unceasingly by the outburst of power in law-determined connection (cf.\ \S\,11),
comes to existence in the things of nature.

Critically understood, the geognostic passes imperceptibly over into the geological;
but, essentially seen, the transition is dialectical. Here, namely, is the pervasive
difference, that geognosy, which begins with the parts and moves between dividing-up and
piecing-together, must at all points strictly hold to the physical apprehension of cause
and effect, of premise and result; whereas geology, which sets out from the whole and
looks to the advancing development, has difficulty in excluding the representation of
lower and higher stages of development, and therefore, in the name of progress, is at
every moment tempted to transpose results into ends. It is naturally only general and
indefinite hints that physics can give about the very first state of the earth-globe;
but already these hints are sufficient to show the relation between the geognostic and
the geological.

If we set out from the assumption that the moon must have been formed from the
earth-mass in a manner corresponding to that in which the latter was formed from the
sun-mass, then the consideration of the moon's formation as a represented result leads
back to the representation of the necessary presuppositions. The earth, pregnant with the
moon, must then have had a different volume than it now has; while the earth-globe's
radius is at present only 860 miles, it must in that time have been about 60,000 miles.
Before the moon came into the world, the maternal earth must then also have turned rather
slowly about its axis: the day-and-night, which is now only 24 hours, must then have
amounted to 27 days. Various things can further be adduced for the explanation of the
earth's flattening, condensation, increased rotational velocity, and so forth. With all
this the consideration moves in a circle of represented phenomena. The whole grounding
amounts to fitting represented presuppositions to a merely represented result; for the
moon's formation from the earth is for the observer only a represented, not a real,
event. But even if we assume that the event has so great a probability for it that it can
be regarded as a fact, is there hereby properly designated any progress? By no means ---
that is, in the geognostic sense. If here in this connection there is to be talk of a
progress of formation, the observer must slip in his own subjective representation. The
only thing that from a purely objective standpoint can be seen in all events, all
formations and changes, is the fulfillment of ``the unchangeable laws.'' But that, as
regards the fulfillment of the laws, no progress takes place is, from the side of physics,
incontrovertible: in the earth, the rotating ball of vapor, the laws were just as
completely fulfilled before as after the moon's formation.

The earth, it is said, was next a fire-fluid mass, and as the cooling proceeded, the heat
radiating into the universe, there began to form a solid crust on its surface, much as a
sheet of ice forms on water; but since ebb and flow, arising from the moon's and the
sun's attraction, called forth constant movements in the perfectly fluid earth-mass, this
crust-formation could not go on without continual and great disturbances. Many a time the
thin cover was burst, until the earth's constantly advancing cooling produced so mighty a
cover, consisting of baked-together stone-flakes, that this at least could not be burst to
pieces more often, but could indeed crack now at one place, now at another. By following
the description in detail and especially considering how the earth-crust as a first rough
draft is constantly transformed and enriched, in that it is broken through by eruptive
masses --- granites, porphyries, greenstones, and so forth --- which differ in composition
and specific gravity according as they come from a different depth; by seeing how the
waters, under their destructive influence on the constituents of the surface, themselves
work along on further formations; how the masses dissolved in the waters and afterward
deposited increase the thickness of the crust here and there, by this increase hinder the
radiation of heat and cause a part of the crust to be melted again from below, so that by
the heating there are formed crystalline schists, gneiss, mica-schist, hornblende-schist,
and the like; by, as said, absorbing ourselves in these first formations, we get
involuntarily the impression that here, in all this, something is being striven for; it
is for imagination --- as though we became witnesses to an incipient execution of a grand
building-plan; as though we got a glimpse into dim workshops where an invisible master
works at fashioning the material, covering over the abysses, laying stone-beams and
building mountains. But geognostically speaking --- it is a fantasy. Here nothing is
prepared; for here nothing is aimed at; the whole is a series of changes in which the laws
are fulfilled: it all goes on as naturally and simply as when a sheet of ice forms on
water in windy weather.

If the geological language sometimes takes a half-poetic flight, then geognosy comes at
once with a prosaic cooling. One opens to oneself the archives of the primeval world,
contemplates the old stone-monuments, speaks of picture-writing, of hieroglyphs whose
meaning it is the task to interpret. What then is the meaning of these cyclopean
masonries, these gigantic granite-blocks, this stacking-up of eurites, porphyries,
gneisses, greywackes, and whatever they are all called? ``Do not ask about it, but
research! Nature itself has no opinion about such things, and on men's opinions it does
not at all depend.'' --- But if there is no meaning in the dim writing, how then will one
be able to interpret it? ``Research, research, for the key lies hidden in the
connection.''

If we now investigate how the earth's development presents itself when we especially
regard it under the standpoint of the whole, then we come to a peculiar apprehension of
the relation between the scientific and the teleological. It must beforehand be expressly
emphasized that what we here apprehend as teleological has nothing to do with
supernaturalist representations of an outside-standing being's personal intentions or
arbitrary interventions in the connection of things, but turns exclusively upon natural
self-formation, upon endowment and end, real-concept and existence. The three grand
hypotheses --- the cosmological, about the origin of the planetary system; the geological,
about the origin of the earth-globe; and the biological, about the origin of living
organisms --- furnish examples of what we in the scientific sense must call natural
self-formations.

Physical astronomy presupposes that matter with its forces is ``eternal'' as space, and
then sets out: from the first being to comprehend the becoming and gradual formation that
the present must in the course of time have undergone. The first being is to be apprehended
as a state in which everything existing can be assumed to have found itself from the
beginning. By a comprehensive reduction we are then led back to a matter uniformly spread
in space, still indifferent, an extraordinarily thin but enormous mass of vapor, whose
diameter is estimated by some at 12,000 billion miles. The formation now begins with the
parts attracting one another, whereby there arise condensation, rotation, and heat. The
mass of vapor becomes globe-shaped, raised at the equator; the rotation increases, and the
centrifugal force grows: all natural. In the raised belt of the globe-mass there will then
gradually form rings, which loosen, separate off, and again roll together; the
rolled-together masses move about their axis and about the principal mass: here we have a
first draft of the system in all naturalness. What happened with the great mass can repeat
itself with the smaller; some of these can again set off one or more rings, these rings can
separate off and become moons, and so forth. --- It is not the merits and defects of the
hypothesis in physical respect that we are here to test: our task is to regard it in
relation to the teleological.

It is evident that, according to such a presupposition, no single phenomenon whatever can
be demonstrated that should especially bear witness to plan and intention. That the
planetary orbits have only a slight inclination, that the periods of revolution are
incommensurable --- $\dfrac{T}{t} = \dfrac{A^{\frac32}}{a^{\frac32}}$ ---, that the sun,
which with respect to the ``perturbations'' is indeed especially important, acquired so
preponderantly great a mass that it is about 700 times greater than the masses together
that move around it, is indeed a telling testimony of law-determined order, but by no means
any objective proof of a peculiar creator-wisdom. Only a single circumstance could seem to
point to such a wisdom. Of the three conic sections only one, namely the ellipse, is fit to
give a closed system of motions; the two others would on the other hand make the system
meaningless, in that the planets, when they move either in a parabolic or in a hyperbolic
branch, gradually remove themselves from the sun, without ever returning. If one will not
write it down to chance that there is reason in the solar system, then it seems undeniable
that precisely the peculiar initial velocity, by which the elliptical orbits are
conditioned, must presuppose a peculiar arrangement of wisdom and thus contain the answer
to Laplace's demand: \textit{Philosophe, montre-moi la main qui a jeté les Planètes sur la
tangente de leur orbite} [Philosopher, show me the hand that cast the planets upon the
tangent of their orbit]. But the answer resolves itself into a reasoning of ignorance.
Rational mechanics proves in an exactly satisfactory manner both why the parabola, with a
probability $= \dfrac{1}{\infty}$, has been excluded, and also what the consequence would
have been if there had been initial velocities that entailed the hyperbola: planets with
hyperbolic orbits would in primeval time have left the system, and we should thus retain
those that received a smaller initial velocity, fit for the ellipse.

Meanwhile there is here yet an obscurity that must be cleared up, an ambiguity we must see
to get removed. How does it really stand with the hypothesis of the eternal atoms? If they
are eternal, then they are absolutely independent. But between absolutely independent atoms
there can be no mutual relation of dependence. If the atoms have properties and forces by
which they can act upon one another, then they are not unchangeable, of course not eternal
either; if they are eternal, in the sense that they are because they are, then we may
attribute to them neither forces nor properties. Even though we cannot penetrate to ``the
innermost causes'' or explore the first beginning of things, we must yet set out to free
our representations of the original from manifest self-contradictions. That there is
something precarious about the hypothesis of the eternal atoms in the eternal primeval vapor,
physics has an easy time convincing itself of. Infinite space cannot have been filled with
atoms, for infinite space can never become filled. But suppose even that the immeasurable
universe has from eternity been filled with primeval vapor: how then should this primeval
vapor have come into motion? --- ``By attraction!'' Among the eternal atoms no attraction is
possible. --- ``Toward the center!'' There is no center. --- The longer one continues these
considerations, the more clearly it is seen that what is here sought is more than masses,
atoms, and forces, that it is nature itself, the eternal ground of existence. But what then
is properly that which we call nature? An original unity of concept and existence. A concept,
even an all-apprehending real-concept, cannot be nature, for the reality is lacking; a sum of
all kinds of things and realities nature can also not be, for the ground is lacking. Nature
must thus be the original unity of ground and reality in the form of ground. But in the
original ground, which has all conditions in itself, since it is the ground of all, the whole
sum of becoming existences has not merely its possibilities, but, in the essential sense, its
endowment.

The development of the planetary system is thus a development of endowment, its formation a
natural self-formation, a formation of an all-preforming natural endowment. The presupposition
thus determined is decisive for the apprehension of the geological hypothesis. If we
comprehend the formations, systems, and groups --- the greywacke group, the coal group, the
Zechstein group, the Triassic, Jurassic, Cretaceous, and Molasse groups --- which geology has
for the present arranged in great intermediate periods, under a general survey, then one can
hardly defend oneself against the fundamental representation that here, under the whole
physical development, there runs a driving stream of means, motives, and ends; but if one
looks sharply at the material causes, then there are none. Why did the glowing earth become
enveloped as in a stone-mantle with many folds? Because the layers in its surface cooled by
radiation of heat, and eruptive masses broke through the tilted layers. Why has it gotten a
sea? Because in the first, surrounding gas-space there were so great quantities of oxygen and
hydrogen, which entered into chemical combination. Why is there in the earth-crust so
extraordinary a spread of silicates, carbonates, alkalis, and so forth? Because the
corresponding elements were all present beforehand in the rotating ball of vapor. It is the
merely physical grounding that continually repeats itself. But the physical real-grounds are
only half grounds, which is immediately seen when one asks about the first essential
presuppositions. Why did the earth-mass originally contain the many disparate substances?
Here the grounding stands still.

The essential connection between the cosmological and the geological hypothesis here shows
itself in the clearest light. If the formation of the planetary system is a natural
self-formation, because it has proceeded from a natural endowment determined in the
real-concept, then the same must hold of the formation of the earth-body. The endowment of
the earth-body must then be said to be contained in the endowment of the solar system. All
development of nature is, essentially seen, a development of endowment, and every development
of endowment a development of ends, a teleological development. In order to determine what
and how much lies in the endowment, we should thus have to know beforehand the relative ends;
in order to determine the relative ends, we should have to know beforehand the last,
concluding end, the final goal; in order to determine the final goal, we should have to be
able to survey the course of the whole, but that is impossible. The cognition of ends becomes
only a human interpretation, a guess, if one will, which, as regards the reality of the
content, cannot make claim to objective-scientific validity. The teleological does not enrich
us with knowledge, but it conditions our insight and is, for the sake of the fundamental
categories, of decisive significance in science.

On behalf of knowledge and real cognition we must, with the researcher, hold to facts and
docilely follow his track, as he, working with particulars, penetrates deeper and deeper into
the real connection of things. But when the very most essential categories of grounding ---
endowment and motive, end and means --- are either overlooked or even disdained, it will be
impossible for the researcher to make a scientific distinction between secondary matter and
main matter, between essential and unessential; with the merely physical real-grounds
everything becomes, for the sake of the connection, equally essential and just therefore
equally unessential. The very moment the researcher declares anything whatever --- be it now a
substratum of stone-masses, formation of sea, upthrust of mountains, extension of mainland,
precipitation of lime, purification of atmosphere, whatever it is --- important and
significant, he slips in an end and lays in a motive.

From this is then finally understood the inner organic connection between the geological and
the biological hypothesis, which, apprehended and expressed in its whole compass by Darwin,
sets research in motion in order to come to its full right. Just as the endowment of the
earth-globe is contained in the endowment of the solar system, so the endowment of organic
life is essentially contained in that of the earth-globe. The organic formation is natural
formation, because it, teleologically understood, is essentially self-formation, a
self-formation that has kept pace with the development of the earth-globe.

\chapter{Teleological-Empirical Sciences}

\section{Teleology and Biology}
\label{sec:17}

\noindent In the endeavor to demonstrate a law-determined continuous transition from
inorganic to organic nature, the physically exact research is near to effacing the
boundary between the lifeless and the living. Quite naturally. The most essential
distinguishing mark of life is to be teleological; but for the teleological the
physically-exact research has no use; a teleological explanation cannot satisfy the
cognition that asks only about observable conditions and materially acting causes. A
natural consequence of this is that the researcher himself, instead of making clear to
himself the scientific element that lies in the concept of teleology, prefers to hold to
mythical representations, against which there must, in the name of science, be polemic. A
supernatural being who has breathed and perhaps still now and then breathes life into
matter does not help science forward; a life-substance for itself or a single and
peculiar vital force is nowhere to be discovered: it then lies near at hand to bring out
the likeness between the inorganic and organic formations so powerfully and many-sidedly
that the transition becomes imperceptible and the laws of nature the same.

The difference, it is said, between the organic and the inorganic, as regards the state
of aggregation, is only that the three states --- the solid, the fluid, the air- or
gas-formed --- which are so characteristic of the inorganic, are in the organic
formations combined into a solid-fluid form that imparts to them a peculiar
``imbibition-capacity.'' On account of this imbibition-capacity the organisms can be
distinguished from the crystals, with which they otherwise, as especially the lowest
skeleton-forms show, have a striking likeness. --- By the help of some limping analogies
and especially by reference to carbon's remarkable capacity to combine with the other
elements in the most varied relations, the researcher can at last persuade himself that it
has now finally been proved how the living is produced with ``absolute necessity'' from
the lifeless. Such an imagining can be quite innocent, provided only the researcher does
not become dogmatic; but should he get an attack of dogmatism and plainly assert that the
origin of life is hereby sufficiently explained, and that especially every attempt in a
teleological direction must be regarded as a sign of outdated cerebral spinnings: then
there must, in the name of science, be a protest.

Investigations of the teleological are seldom carried through with full consistency, in
that they are all too often encumbered with double-sided misunderstandings. The defenders
of teleology readily set out from the presupposition that it is the researcher as
researcher who needs a course in the doctrine of ends and means, that a thorough,
speculative-dogmatic doctrine of purpose, in combination with some miracles of creation
distilled into an ingenious explanation of a reliable tradition, would best help him
forward on the way of research and show him the bridge over the deep gulf from the
lifeless to the living. But the researcher rejects the proffered direction, and that for
the good reason that it does not at all suit his scientific task. The misunderstanding
turns about, in that the researcher, who admittedly has no use for the teleological, feels
obliged to drive teleology out of science by physically exact proofs. And what does he then
prove? Nothing that in the least touches the concept of teleology. He sets out from the
presupposition that what demonstrably proceeds as a result of physical causes and happens
with conditioned necessity cannot possibly be teleological; but how great this
misunderstanding is can be seen from the analysis of real-concepts. He urges that the
imagining of the teleological hampers research, in that one, especially in biology and
physiology, finds it more convenient to make use of turns of phrase about ``the vital
force,'' ``the vital power,'' ``the life-principle,'' ``the Idea,'' and the like, instead
of working with the anatomy of the tissues and the chemistry of carbon-combinations. But
that teleology will by no means exempt the researcher from work follows simply from the
nature of the case.

That teleology and biology are connected concepts makes itself involuntarily known in the
characteristic slogan: struggle for life; it is a teleological watchword. Despite the
analogies that can otherwise be demonstrated between inorganic and organic bodies, the
differences are yet so striking that they cannot possibly escape common sense; they make
themselves immediately known by characteristic words and turns of phrase. That inorganic
bodies, especially the solid, make resistance against what would penetrate into them and
change the connection of the parts, is undeniable, and one can so far regard every thing
as an individual that, in relation to its forces, maintains its independence; but no one
will derive the inorganic individual's resistance against dissolution from an inner urge to
self-preservation. No one will assert that the ice from an inner urge to hold itself fast
defends itself against the heat, or that the water feels a quiet longing to crystallize.
The inorganic individuals, which punctually fulfill the law of action and reaction, are
wholly indifferent to the states in which they find themselves and the changes they
undergo; it is as indifferent to the cliff whether it preserves its bearing or is
shattered by the lightning as it is indifferent to the earth-crust whether it cracks in
the sun or is overflowed by the sea: this complete indifference is precisely the
distinguishing mark of the lifeless. But that the self-maintenance of the organic
individuals must be regarded from a quite different standpoint, science itself inculcates,
in that it sets the struggle for life as an irrefutable condition for all that is living.
Self-preservation is hereby apprehended as an express end, that is, as the task nature has
set every living individual.

This task impresses on all organic functions a peculiar stamp, so that they, despite
physical likenesses, essentially distinguish themselves from corresponding inorganic ones.
A crystal can grow as well as a plant and an animal; but that organic growth has a quite
different connection than the inorganic is seen in nutrition. The plant grows on account of
nutrition, but the crystal by a stacking-up of parts; its growth is without nutrition.
Nutrition, which unites in itself metabolism with assimilation and the maintenance of a
typical form, designates a teleological cycle; nutrition takes place not merely so that,
but expressly in order that, the organic individual can sustain life.

If we ask about the material causes and conditions on which the nourishing activity
depends, then the researcher is right when he point by point refers to the physical
constitution, deposition, and transposition of the parts, the molecules; but he is not
right when he, by analogies drawn from the inorganic, supposes he can ground the
self-continuing activity on a continued cancellation of mechanical and chemical
equilibrium and thereby explain away the teleological. The examples with ``diosmosis''
prove precisely the opposite of what they should. For they prove, when one holds strictly
to the inorganic, that the equalization stops after a shorter or longer lapse of time, so
that the movement stops; but what was to be proved is precisely how it comes about that the
activity is constantly maintained. What is adduced about the imbibition-capacity, greater
tension in the relation of the molecules, the share of sunlight in the separation of
carbonic acid, and so forth, belongs indeed among the physical conditions; but as regards
the phenomenon of nutrition itself, the grounds of explanation supported on this are so
weak that the researcher, if he had any prospect of attaining others and better ones,
would doubtless himself be the first to lay bare their weakness. Here something is lacking,
but what is it, and whence does it come?

What is it? Life itself! In order, if possible, to make it fully evident what this means,
we must above all see to secure for ourselves a clear and unambiguous usage. There is
activity in the inorganic as well as in the organic; if now there is any difference in
principle between lifeless and living, then there must be a corresponding difference
between inorganic activity and life-activity, and if this difference is to become quite
distinct to us, we must seek to get it designated as sharply as possible. We therefore
begin by making a distinction between substance-activity and self-activity, understanding
by substance-activity every expression of force that proceeds from substance as substance,
while by self-activity we think of functions that have their starting-point not in a
particular substance, but in an ideal causality corresponding to substance-conditions.

What now are the criteria decisive for all substance-activity? If one represents to oneself
two parts of mass $m_1$ and $m_2$ acting on each other, and lets the law according to which
the activity takes place be denoted by $f$, the distance by $u$, the activity by $V$, then
the equation $V = m_1 m_2\, f(u)$ can be set up as symbol and fundamental formula for all
possible substance-activity. By application of the so-called potential functions and the
relation of the ``potential'' to ``living force,'' it can be exactly proved that $f(u)$ can
be anything whatever, $u$ vary between the boundaries $0$ and $\infty$, the number of parts
of mass be as great, and their combinations as different, as one will, noting each time the
difference between movable and immovable parts of mass, even the parts the potential
concerns, and still the symbol of substance-activity --- the activity $f(u)$ being
repelling or attracting, mechanical, magnetic, electric, chemical, in short, of whatever
kind one will --- will keep itself unchanged. In order to come to cognition of what is
properly to be understood by substance-activity, we therefore need only subject the
fundamental formula to a logical analysis.

From the formula $V = m_1 m_2\, f(u)$, when it is developed with respect to positions,
projections, and so forth, there then appears: $\alpha$) that matter under all conditions
acts according to the law of inertia: $m_1$ acts on $m_2$, not in consequence of any
self-determination or inner prompting, but solely because $m_2$ acts on $m_1$; it is
self-less activity; $\beta$) how far there is in a system of active masses or molecules
rest or motion is decided not by any self-governing unity, either in the whole or in any
single part, but solely by the resultant of the cooperating forces; $\gamma$) the system's
transition from motion to rest, or conversely, has its ground not in the system itself, but
depends on outward causes, outward conditions; $\delta$) the active terms are sheer
substance-subjects, and what is peculiar in the activity is referred to some property or
other characteristic of the substance in question, in mechanical physics to gravity, in
chemistry to affinity. If this is now clear, then we need not continue the analysis further
in order to convince ourselves that substance-activity, however many conditions and
combinations one may slip in, cannot possibly be transformed into life-activity. Here
plainly something is lacking, and it is this something that we shall now characterize as the
self-active.

What is self-activity, and by what do we know it? It lies near at hand to point to an
immediate consciousness of willing and will; for without regard to what reality may be
attributed to our willing and self-determining, it is incontrovertible that we have in this
a form of activity that is diametrically opposite to the material. What we call our self is
no substance existing for itself, and the activity proceeding from the self no
substance-activity. The active self cannot, like a part of mass, exist as a term in
opposition to its activity, for it has no immediate subsistence; its subsistence is an
ideal acting, and this acting an unbroken reflecting. From its ideal center the activity
passes over upon matter and from matter back to the starting-point; self-activity is a
circular movement, in the course of which the opposition of center and periphery constantly
mirrors the unity.

Let the self-determining element, $s$, be surrounded by material points, $a$, $b$, $c$; then
the relation will not be such that $s$ can act only in so far as it immediately reacts
against an influence. The reciprocal action takes place, on the contrary, in such a way that
$s$, affected by $a$, can thereby be moved to act on $b$, or, if it affects $a$, then act
now attractingly, now repellingly, according as the state is in which it finds itself.
Self-activity rests upon the fact that $s$ acts according to its condition, and therefore
essentially has the motive of its activity in itself. A system of material parts gives only
a unity of connection; only when the self-determining element comes in does there show
itself a governing unity, a self-unity.

That a governing unity, a form-giving self-unity, is traced everywhere in all living
organisms and organic formations, in the moneron and the cytode as well as in man's
self-conscious soul, is so striking a fact that no physiologist will attempt to get it
explained away, unless it happens in despair at not being able to explain it.

Whence comes the forming, governing unity, the ideal causality, the acting self-unity? From
the system of real-concepts, from the endowment of life. This answer must of course be
rejected by those who are beforehand ensnared in the fixed idea that active real-concepts
with essential endowments do not belong in the order of nature. But even if one grants the
validity of a general endowment of life, the matter is in itself yet dim and enigmatic. How
is one here, without smuggling in unscientific presuppositions or remaining standing at the
general, and for that reason empty, reference to the relation between real-concept and
existence, to be able to form any definite representation of the relation of the self-active
to substance?

Here there are several hypotheses to choose between. One can, for example, set out from a
life-hypothesis analogous with the light-hypothesis, and assume that just as light has its
real possibility in a medium different from the ponderable masses, so life has its
endowment in a corresponding medium spread everywhere. When light is kindled, it happens
under the condition that the light-medium is set in activity. That the ideal self-activity
must be conditioned by a corresponding substance-activity, all who know the laws of nature
will indeed grant without further ado. But the self-active may very well be conditioned by
substance, without therefore being a substance-subject. That self-activity sets in rests
then on the fact that the mechanical, chemical-physical, and chemical in the
substance-activity sets the life-medium in vibration: under the vibrations of the
life-medium starting-points for self-activity form themselves. In a molecule $A$ let the
mass-atoms be $a$, $b$, $c$, and so forth, the ether-atoms $\alpha$, $\beta$, $\gamma$
\ldots; the mass-atoms will attract one another, but the ether-atoms repel them and remove
them from one another. If now there is no self-activity in the molecule, then the tension
can be equalized physically, that is, the molecule is lifeless. If, on the other hand,
substance-activity has set the life-medium in motion, then there will have formed a
multitude of selfhood-points $s$, $s$, $s$ between the atoms; these selfhood-points may
each for itself be so weak that their activity approaches a differential activity: they
will yet bring it about that the molecule becomes living. If now a living cell consists of
living molecules, then the governing unity arises by the fact that $s$, $s$, $s$ are not
merely in reciprocal action with the material parts, but under this reciprocal action at
the same time in such a relation to one another that they form a resultant $S$, which
differs from every resultant of material forces in this, that it, as selfhood-resultant,
governs its components and reduces them to servants of its peculiar demands. If the organism
consists of several cells, then $S$, $S$ will go in under a higher resultant, which in
teleological self-activity will direct the division of labor, and so forth.

If one will prefer another hypothesis and assume that it is the atoms themselves which, in
consequence of a duplicity in the essence of matter, now hold to a mechanical and chemical
standpoint as immediate substance-subjects subject to the law of inertia, now again,
especially in carbon-combinations, raise themselves to furnishing selfhood-resultants, then
this can happen only on the express condition that a scientific account is given of where
the inertia ceases and where the self-activity begins, as well as of how it comes about that
the same substance-parts can at once preserve their inertia and at the same time function as
moments in the organizing self-activity.

If one will finally give up all hypotheses, appeal to the restriction of knowledge, and
declare the biological problem simply insoluble, then this is undeniably an expedient; and
research can undisturbed go its quiet way. But to abolish the factor of self-activity and
explain away the problem by the help of hollow analogies and physical ambiguities is by no
means admissible. The fundamental formula $m_1 m_2\, f(u)$ will continue to bear witness to
the peculiarly bounded sphere of substance-activity, and thereby to the qualitative
difference between lifeless and living.

Life-activity is neither substance-activity for itself nor self-activity for itself, but a
pervasive unity of both. Every organic function is at once a substance-function and a
selfhood-function. From this standpoint it holds that ``the brain thinks, just as the
stomach digests''; yet it must, not merely in the name of the brain and the stomach, but
above all in the name of science, be most emphatically brought out that the relation between
the substance-moment and the selfhood-moment in the organic double-functions is articulated
down to the smallest particulars.

As a researching science physiology cannot directly make use of the teleological; it is and
remains an empirical science. But deny or overlook the significance of teleology it cannot
either, for physiology is in essential kinship with biology. Physiology cannot set out to
make biology impossible; but biology and teleology are necessarily connected concepts.

\section{Botany and Zoology}
\label{sec:18}

\noindent Although botany, which investigates the plant kingdom and its products, is an
empirical science just as well as geognosy, which investigates the strata, there is yet
here, precisely on account of the difference of material, an essential difference between
empiricism and empiricism. Geognosy's material is simply physical, botany's is organic,
hence teleological. Even if one abolishes teleology ten times over, the simplest
description of a plant with root, stem, stem-leaves, sepals, petals, stamens, and so
forth, will involuntarily demand an explanation of these organs' use for the whole
organism, and thereby betray the teleological. The excuse that we humans are obliged to
apply teleology where it does not at all fit takes a strange form in a science. The
subdivisions of general botany --- plant anatomy, morphology, plant physiology, plant
chemistry, and so forth --- can indeed be treated separately in opposition to the
doctrine of systems; but it is yet characteristic to see how systematic botany, which
takes up the results of the particular investigations, leads by the empirical way deeper
and deeper into teleological questions.

One begins immediately by taking the typical forms as they are found. For use in the
merely heuretic classification one seizes upon any mark, even the most constant, quite
at random. The Linnaean system (the sexual system), which is built on the number and
other relations of the organs of fertilization, furnishes a striking example of this.
One classifies plants as if they were minerals. But the plant carries on a peculiar
husbandry, and continued observation teaches that to this corresponds a development of
form that shows on what stage of formation the plant stands. The ``artificial system,''
which is formed by isolating, through abstraction, certain marks from the organism's
other peculiarities, becomes thus unsatisfactory. According to the sexual system plants
that are otherwise utterly unlike each other can come to stand quite near each other when
they happen, with respect to the chosen mark, to be alike; thus, for example, a barberry
(\textit{Berberis}) and a snowdrop (\textit{Galanthus}) both belong to the sixth class,
an ash (\textit{Fraxinus}) and a glasswort (\textit{Salicornia}) both to the second
class, although nature has, by a highly different manner of organization, separated the
barberry far from the snowdrop, the ash far from the glasswort. On the other hand, plants
that in all else show themselves mutually related can become far removed from each other
when it happens that they are different with respect to the chosen distinguishing mark.
Thus, according to the Linnaean system, agrimony (\textit{Agrimonia}) belongs to the
eleventh, meadowsweet (\textit{Spiraea}) to the twelfth, lady's mantle
(\textit{Alchemilla}) to the fourth, and burnet (\textit{Poterium}) to the twenty-first
class, although according to their natural peculiarity they are all members of one and the
same family, namely the rose family.

According to the natural system it is a whole totality of organic forms that must claim
the attention. If we regard from this standpoint the two great, comprehensive groups,
cryptogams (the flowerless plants) and phanerogams (flowering plants), more closely, then
the teleological character of the development of form also shows itself. There form, for
example, spores and seeds, not merely so that, but expressly in order that new specimens
can arise. If the cryptogamous spore-case (\textit{sporangium}) is set against the
phanerogamous seed-case (\textit{pericarpium}), then one might, by a superficial
consideration, easily be misled into confusing the lower with the higher. Thus the mosses
have a spore-capsule formed of a large complex of cells. The spore-capsule opens across,
and in not a few cases in such a way that a spring-ring (\textit{elater}) located between
the lid and the capsule throws off the lid; the so-called peristome has now a single, now
a double wreath of teeth, and so forth; here is a rich equipment. In comparison with such
richly equipped spore-capsules the seed-case of many a phanerogamous plant must,
especially when this seed-case, as in \textit{Lathyrus latifolius}, remains standing at
the form of an ordinary leaf, undeniably take a poor appearance and seem to us both simple
and insignificant. It is otherwise when we look at the peculiar endowment and the
morphological foundation corresponding to it. The morphological foundation coincides with
the element upon which the formation chiefly rests. Such an element is, as regards the
formation of the spore-capsule, the cell; as regards the seed-case, it is the leaf. If for
the organism we call the cell the primary element, a formation-element of the first order,
then we must call the leaf --- the formation in which the cell-element is reduced to a
moment --- a higher element, an element of the second order. However high the formation of
the spore-capsule may rise, its inner structure yet rests only upon cells and aggregations
of cells, whereas the lowest seed-case is originally a leaf. The lowest phanerogamous
fruit-leaf thus stands, in virtue of the fundamental form, higher than the highest
phanerogamous spore-capsule.

Despite surprising advances in the anatomy and chemistry of plants, it will yet not be
possible to explain the activity of form-giving by physical-chemical causes alone. The
fact that chemistry on its present stage can artificially produce combinations that are
otherwise formed only in the living organism is insufficient; for against this stand yet
other facts. It is indeed notorious that two chemical products which resemble each other in
each and all can have arisen by very different ways. It is moreover incontrovertible that
organic products produced by art do not have organic form; should they nonetheless be
produced in the selfsame manner as in the living organism, then one would be obliged to
assume two disparate forces, namely, besides the force in the affinities that produced the
product, a peculiar force that produced the organic form.

It is vain to appeal to ``the peculiar conditions''; for one of two: either the conditions
are physical-chemical, and then the fundamental formula $m_1 m_2\, f(u)$ (\S\,17) shows
that everything falls under the lifeless; or a certain self-activity must be understood,
and then we have the teleological. That the teleological makes no breach in the
unchangeable nature of the substances, the affinities, and the laws follows simply from
the fact that the self-active must precisely be presupposed to belong essentially to this
nature. That outward conditions can act not merely favoringly and modifyingly, but even
inhibitingly and disturbingly upon the organic activity of form-giving, is a matter of
course when we consider that self-activity, not here and there alone, but term by term and
point by point, is conditioned by substance-activity.

He who would have a great and clear impression of the teleological activity of
form-giving must not remain standing at an elementary comparison between the crystal and
the plant, for from this there appears only what no one can deny --- that physical causes
and conditions are everywhere at work; he must go a step further, he must consider the
life-courses of plants; crystals have no life-course! He must consider what Goethe calls
``die Metamorphose der Pflanzen'' [the metamorphosis of plants]; he must convince himself
that all the different parts of a phanerogamous plant, from the cotyledon to the pistil,
arise by the fact that one and the same element, at the different heights of the axis, puts
on a different shape, in consequence of the whole economy. What at one stage is seed-leaf
(cotyledon) is at another foliage-leaf, at a third bract, at a fourth sepal, at a fifth
petal, at a sixth stamen, at a seventh fruit-leaf. --- If yet a further proof is demanded
that the plant's parts arise not by finite piecing-together, but are developed from within
out of one self-differentiating endowment, then we need only recall the manner in which
especially Schleiden has tested and carried through what Goethe intuited. Instead of merely
comparing the phenomena in their immediate existence, Schleiden asks also about their
coming-to-be, goes back to the original, compares, for example, a cotyledon in its first
appearance with the first appearances of other leaves on the same plant, and then finds
that all the leaves from the beginning resemble each other, that they are all originally
of one form, each a ``wart-shaped protrusion'' below the tip of the axis.

In zoology the teleological becomes, in the nature of the case, still more prominent. If
it is from a physical standpoint impossible to demonstrate any properly decisive
difference between lifeless and living, then it is surely just as impossible to find the
boundary-line where the plant kingdom leaves off and proper animal life begins. The
intervening province of the ``protists'' is as yet hardly sufficiently investigated. That
the marks usually adduced for the animal organism --- sensation and voluntary motion ---
do not satisfy the researcher is understandable. If it is already difficult to come to
agreement about what the named expressions properly mean, then it is still more difficult
to secure their application by unambiguous facts. An aphoristic setting-up of such
criteria will always have its precarious features. Yet the teleological mode of view
presents in this respect practical hints.

Have plants a soul? This very ancient question has in the more recent time returned with
reinforced emphasis, and acquires a peculiar significance from the fact that it contains a
summons to an ever deeper-going account of the original relation between life and soul. But
when one, by referring to ``activity of formation, healing power of nature, reflex-motions,
instinct, and the impulse toward beauty,'' has exhausted everything that can speak for
attributing to the plant a soul, then one must yet halt at a question that can be explained
only teleologically. Why does plant life, with its otherwise rich development of form, not
go in a soul-ward direction? The concentration to free inwardness in sensation and
self-determination, which is the distinguishing mark of soul-life, is wholly lacking in the
plant. Plant life, which on the whole is so endowed that it cannot go beyond the boundary of
material organizing, must, with this quality-boundary lying in the endowment, remain
standing at a life without soul. If a living being is to be able to raise itself above the
selfhood streaming in immediate vegetating to a self reflected into itself, then life itself
must from the endowment be soul-endowed.

The two categories which, in comparison between plant and animal, are confused in an all-too
superficial manner are degree and measure. The plant, it is said, has sensation and
consciousness as well as the animal, but naturally in far weaker degree. It then comes to
look as though the supposed plant-soul could work its way imperceptibly up to a higher and
higher degree of conscious living, for for mere gradations there is no fixed boundary; it
looks as though the plant had the capacity for consciousness, but merely lacked the right,
corresponding organs --- nerves, ganglia, instruments of sense and motion. But that the
relation is precisely the reverse lies in the teleological character of the organic
functions.

Physically seen, it is said that the animal senses and moves merely because it has nerves;
but teleologically understood it holds, on the other hand, that the animal has nerves
because it is endowed for sensing and moving. In the higher animals, in which the
distinguishing marks become distinctly expressed, this is seen, as is well known, already
in the first stages of embryonic formation. Cuvier's proposition, that it is the form of the
nervous system that determines the whole animal form, will always maintain its validity. For
the one-sided, skewed, or radiate animals (\textit{Animalia mollusca}) without real limbs
the scattered, irregular nervous system is fitting; the ventral-ganglion system for the
one-sided, jointed animals (\textit{A.\ articulata}) without inner skeleton, with many or no
limbs; the spinal-cord system for the one-sided, unjointed animals (\textit{A.\
vertebrata}), whose bone-structure encloses the spinal cord, the axis of the nervous system,
and supports four internally jointed limbs; it is the pervasive main thought in systematics,
whatever modifications may be made with the schematic.

It is a grand advance that biology has made in the most recent time, in that one has fully
become conscious that botany and zoology cannot be cultivated as isolated fields. Earlier
these fields touched each other only when low, ambiguous forms awakened uncertainty about
whether it was a plant or an animal that one had before one in the specimen in question. The
``protists'' have not merely widened the investigations; but, what for scientific
consciousness is still more important, the investigations have led to a methodical
connection between zoology, protistics, and botany (Ernst Haeckel: \textit{Generelle
Morphologie der Organismen}. First Volume.\ Berlin 1866.\ p.\ 191 ff.). As regards the
chemical character, it is among other things brought out that the animals are distinguished
by the formation of nitrogenous acids (uric acid, hippuric acid, inosinic acid, and so
forth), while the nitrogen-free acids play a greater role in the plants; that the animals'
nitrogenous bases (urea, creatine, leucine, and so forth) are weakly alkaline and of an
almost constant composition, whereas the reverse is the case with the plants; that the
plants preferentially, by reduction and synthesis of quite simple into highly composite
combinations, bind heat and develop little mechanical work, while the animal organism
conversely, by oxidation and analysis of highly composite combinations, forms simpler ones,
frees heat, and develops mechanical work. ``The animal is an oxidation-organism.''

The circumstance that the chief sum of the chemical processes in the animal body aims
especially at developing a quantity of living force --- namely by oxidizing the more involved
and loose carbon-combinations and transforming them into simple and firm combinations such as
water, carbonic acid, ammonia --- explains that there can be forces enough to set muscles and
nerves in motion. But when the physicist then has to determine more closely how it really
stands with ``the peculiar molecular motion of the representation,'' then he becomes
unintelligible. ``Only'' --- it is said (op.\ cit.\ p.\ 234) --- ``when representations are
awakened in the albumin-molecules of the ganglion-cells can we speak of an animal's soul, and
we then designate the representation that arises under awakening in the ganglion-cells by the
centripetal threads as sensation, the representation, on the other hand, that is produced by
awakening in the centrifugal threads of the ganglion-cells as will. The animal soul's dimmest
and highest function, which is most difficult to grasp, is the formation of thought, which
consists in representations and arises under conduction, but probably always by a highly
complicated reciprocal action of numerous centrifugal and centripetal awakenings.''

We have here an example of how immediately the scientific and the meaningless can touch each
other. It is organic functions, double-functions, that are here described. If we hold to the
substance-activity, then the description is indeed as apt and exact as it can be on science's
present stage. If, on the other hand, we seek in the functions that which properly concerns
``the representation,'' ``the sensation,'' ``the will,'' ``the thought,'' namely the
self-activity, then the description is completely meaningless. What is representation,
sensation, will, thought? Whence do these words draw their meaning? Plainly from our own inner
experience. I know what representation is, for I can represent; I know what sensation is, for I
can sense; I know what will is, for I can will; I know what a thought is, for I can think; I
know what life is, what it is to live, for I live. We have repeated this ``I know'' in order to
inculcate what the manifold considerations of substance continually bring to oblivion --- that
the representing, sensing, willing, living subject is not a substance-subject, but a
selfhood-subject, that the ideal activity which here shows itself, however conditioned by
substance it may be, and however weak the consciousness that accompanies it, is yet not a
substance-activity, but a self-activity.

If we hold fast the starting-point for the apprehension and understanding of what we call
self-activity --- for that in all sensing, however dim, weak, and unconscious it may be, there
is a self-acting, can be expressly demonstrated --- then such determinations as are drawn from
the clear consciousness will also be able to be carried over to the dim, indeed, with peculiar
limitation, to the unconscious life. That functions such as sensing, representing, thinking,
willing must belong to the organic functions just as well as respiration, blood-circulation,
digestion, and so forth, the physiologist will indeed, in the name of physics, most
emphatically maintain. Well then, we grant him this; but what is the consequence? If the
physiologist will assert that in thinking and willing there is no moment of self-activity, no
trace of teleology to be discovered: then he must excuse us for finding him unintelligible; for
the unintelligible shows itself at once when one has leaped over the understanding. If, on the
other hand, he grants that in thinking and willing there is an appearance of something
self-active, wherever it may come from, then it follows from this that in the other organic
functions there must also be a moment of self-activity, however hidden it may be.

The question of the self-active hangs most closely together with the question of the
teleological, this again with the question of endowment, and this again with the question of
the stages of life's development, of the organic types, of genera and species.

\section{The Biological Doctrine of Descent}
\label{sec:19}

\noindent No natural researcher can make immediate use of the teleological, and no
researching biologist can overlook the real-concept of teleology without going astray in
the biological labyrinths. Telling proofs of this are to be had in the investigations of
the emergence of organic life on earth. The older researchers set out from the dogma of
creation, and assumed without further ado that plants, animals, and men had from the
beginning arisen by a supernatural --- that is, miraculous --- act of creation, that each
genus and each species had, in consequence of the Creator's plan, its unchangeable
character, which maintained itself in propagation. But gradually, as science advanced, one
found oneself embarrassed by the miraculous presuppositions. Experience showed that
especially the species had no fixed boundaries, since there were varieties. Reflection
taught that the assumption of supernatural acts of creation not only explained nothing,
since they were themselves inexplicable, and it became more and more troublesome to
ascertain the Creator's intention, especially with the parasites, but that a mixing-in of
the miraculous was simply dangerous, since it stood in manifest conflict with the
fundamental principle of the unchangeableness of the laws. One sought to hush up, fade out,
and veil the actual miracle of creation, as far as it could be done; but even if one
appealed directly to nature's own creative forces, these had yet to take on a rather
supernatural appearance, so long as one had not yet found the thread of continuity in the
course of the earth's development. According to the catastrophic-dramatic presentation of
geological upheavals that extinguished all life on earth, one was indeed continually in
embarrassment with beginnings afresh, and had to take refuge in new, extraordinary acts of
creation. Only with Lyell's ``Principles'' did the whole mode of view undergo an essential
change. Nature makes no leaps, and everything in nature must go on naturally --- these
were scientific thoughts that now in earnest found entrance. What Lamarck and several others
had spelled out came to full natural-historical expression in Darwin.

As every teleological consideration is resolutely given up, the theory of descent is then
worked out in such a way that the living organism, regarded as a physical product of the
mechanical and chemical activity of general natural forces, presupposes no inner, peculiar
endowment, but only a sum of outer, favorable conditions. It is now carbon-combinations and
chiefly albuminous substances that take the place of the miracle and, by a peculiar plastic
capacity, bring about the otherwise inexplicable primeval formation, archegony. Archegony
is, as the protists show, not originally bound to any definite form; typical endowments
have on the whole no significance, unless by endowment one will understand the result of
mechanical work that can, from a preceding organism, pass by propagation as inheritance to a
succeeding one. Life works its way forward with the same kind of necessity as that with
which a broad river that has sufficient fall moves. Just as the river arranges its bed, its
bends and ramifications according to the physical constitution of the soil, without any
inner striving or urge toward a definite outlet, so the forward-striving life arranges its
forms and types according to the outer circumstances, the favorable or unfavorable
conditions under which it works. It strives and struggles; but that its striving is without
goal follows simply from the fact that the teleological is abolished. When from a
teleological standpoint it looks as though life, in a system of typical endowments with
lower and higher ends, unfolded itself according to a plan of higher wisdom, then this is,
regarded from a physical standpoint, an imagining. General life stands in its forms lower
and higher according to the same law of limitation as the different lower and higher
water-level in the river. That there is otherwise great difference between the streams of
life and of water goes without saying when we look at ``the absolute necessity
(ἀνάγκη).''

In this natural-scientific style the hypothesis of descent has, in a truly grand manner,
been thoroughly worked through and equipped with an extraordinary apparatus of learning by
E.\ Haeckel. The protists --- monera, protoplasts, diatoms, and so forth --- form the
autogonic foundation; from them descend, in the one direction the plants, in the other the
animals. The genealogical tree of the plants is pursued through all ramifications from
archephytes and florideae up to the phanerogamous gymnosperms and angiosperms; the
pedigree of the animals reaches from the plant-animals (\textit{Coelenterata}) and the
star-animals (\textit{Echinodermata}) through arthropods (\textit{Articulata}) and mollusks
to the vertebrates. The vertebrates descend especially from the worms, and man, the most
excellent vertebrate, from a primeval form related to the apes.

Embryology, paleontology, and comparative anatomy, in combination with the laws of quality
(the theory of selection), the laws of inheritance, the conservative and the progressive,
the laws for direct and indirect adaptation, and so forth, give on the whole the impression
that the doctrine of descent is scientifically fruitful. It is indeed characteristic that
embryos of the most different species within the same class have at the beginning the same
appearance, that the embryos of mammals, birds, salamanders, and snakes in the earliest
states resemble each other not merely on the whole, but also in the manner in which the
single parts are developed, that the likeness goes so far that they can hardly be
distinguished from one another except by the help of size; it is characteristic that even
the human fetus, with certain modifications, passes through developments of form that not
only recall the kinship with nearly standing animal species, but even point back to the
protists; it is indeed characteristic that series of paleontological formations upward must
furnish commentaries in the interpretation of the now-living forms. That with respect to
rudimentary, abortive, stunted, miscarried, atrophic organs there can even be set up a
``dysteleology,'' a doctrine of unpurposiveness, is a confirmation of the unbroken
cooperation of the physical causes and conditions. From a protistic beginning the emergence
of man is prepared step by step; we see how our invertebrate forefathers raise themselves up
eight stages to those endowed with vertebrae. The sea-squirt (\textit{Ascidia}) on the
eighth and the lancelet (\textit{Amphioxus lanceolatus}), the lowest representative of the
vertebrates, on the ninth stage, form the transition to the monorrhines, those of our
kinsmen that have, though only a distant, likeness to the now-living lampreys and
cyclostomes; thence the way goes through the fishes, the monotremes, the marsupials, the
half-apes, and the narrow-nosed apes with tails to the man-apes on the twentieth stage and
the ape-men on the twenty-first, until we finally reach the genuine, speaking men on the
twenty-second stage, and systematically survey the whole ``natural creation.''

In so far as one holds exclusively to the objective standpoint of empirical research, one is
indeed obliged to grant this hypothesis of descent the same scientific significance for
natural history as we grant the Copernican hypothesis for astronomy. Modifications,
limitations, extensions, correctives will by no means diminish its worth, but on the
contrary prove its fruitfulness. When we nevertheless express it as our conviction that the
task set is not solved, then it happens only in order to call attention to the fact that a
task of this nature cannot possibly be solved in a physically objective manner. Here are
lacking explanations of the mechanical forces that unbrokenly move the life-stream from
below upward, of the particular conditions under which the knots are tied on the thread of
continuous development, of the characteristic limitation within certain more or less
constant principal forms, while the species vary and pass over into one another, and so
forth. Why precisely these organic types and not others? Such questions the researcher can
indeed parry by other questions: Why precisely these crystal types with these combinations?
Why these mountain types, indeed these typical cloud formations? To offer him a teleological
explanation does not help; it explains nothing to him; this must be granted. But when the
researcher then goes a step further, scoffs at the teleological, and asserts that the
physically objective research has the whole scientific content of the task in its power,
then we declare to him plainly that he is in an error.

How rash it is to set the biological task on the same level as the astronomical and assure us
that the Darwinian theory can ground the phenomena of life by the help of physical causes and
general laws of nature in the same manner as Newton has grounded the astronomical phenomena,
can without difficulty be seen. In order to explain the moon's motion about the earth, Newton
needed only to hold fast two decisive conditions: attraction toward the center and motion
along the tangent. The matter is indeed by no means hereby settled; for the relations are
involved. The formula for the moon's longitude presented by Delaunay in \textit{Connaissance
des temps} for 1865 takes up no fewer than 45 printed pages, and yet this formula takes
account only of the perturbations due to the sun. Despite such involved relations the
astronomer does not, however, resort to evasions such as that we naturally ``do not know the
innermost causes''; no, he makes it in great, simple strokes evident --- so evident that a
schoolboy can understand it --- how the motion in a curved line must arise, and how it can be
modified by outward conditions. From the simple first we can then get a glimpse into the
involved and understand the fundamental thought, even though we cannot follow the analyses.
But how would it go with the hypothesis of descent, if the many disparate laws of development
were to be derived, even but in part, from general laws of nature? The organic function is not
explained; autogony is not explained; the peculiar impulse of formation is not explained. In
order to be rid of the teleological, one has reduced the doctrine of selection with its
consequences to a plain and simple doctrine of utility, without even becoming aware that the
doctrine of utility is precisely a doctrine of purpose valid in the finite, hence a teleology.

But --- one might here object --- what help is this continual pointing to the teleological,
when it is yet plainly granted that teleology cannot be used in physically objective research?
We answer: physically objective research, indispensable as it is in science, does not exhaust
the scientific. The scientific in biology is two-parted: the one part belongs to the
empirical, the other to the \textit{a~priori}. Such a bipartition does not occur in the
real-empirical sciences. Astronomy, chemistry, mineralogy, geognosy have nothing belonging to
their tasks outside themselves; they are researching sciences, and the tasks are tasks of
research: everything falls under substance-activity. It is otherwise with biology, whose
object is organic life; the organic functions are double-sided, and the biological task must
as a consequence necessarily get a double-sided character.

What here can easily confuse is precisely the circumstance that the doubling requires two
disparate investigations, empirical research and \textit{a~priori} analysis; indeed, what is
more, that the disparate investigations even belong under disparate sciences, the empirical
research under physics, the \textit{a~priori}, or rather \textit{a~priori}-zing, reflection
under philosophy. We can see it in definite examples. If the talk is of autogony, the
self-formation of organic life, the physicist will one-sidedly fix his attention on the
peculiarity of the substances and combinations and apprehend the life-activity exclusively as
a peculiarly conditioned substance-activity; the \textit{a~priori}-reflecting person will with
similar one-sidedness fix his attention on the self-active, cursorily overlook the
substance-conditions, and regard the whole formation as a work of the self-being's forming
capacity. If it should sometime, which is yet perhaps possible, succeed for the natural
researcher to produce a living cell, then he will see in this an irrefutable proof of the
assertion that life is a physical resultant of the cooperation of physical forces; whereas
the \textit{a~priori}-reflecting person will urge that the conditions brought about by art
have only made possible a breakthrough of the hidden principle of selfhood. If the talk is of
an organic type and of the drive that makes itself known in the type's formation, then the
empiricist will hold exclusively to consequences of material presuppositions, and in the
so-called endowment see only a sum of conditions brought together piecemeal, which according
to circumstances can be changed; while the \textit{a~priori}-reflecting person will be tempted
to hypostatize the type as a finished idea-form, an ideal paradigm from which there can be no
deviation. If man is compared with the lowest animals, and it is a matter of a biological
valuation of the intermediate stages life has had to run through before it reached the stage
of completion --- that is, if man is the completion --- then the modes of view will likewise
diverge. The physicist will indeed not be able to avoid speaking of inner impulse of
formation, impulse of self-preservation, a struggle for life that leads to the perfecting of
certain organs, an influence of habit, a cunning of instinct, and so forth; but when it
hereby takes on the appearance that he is speaking teleologically, he is at once ready to
disavow any thought of end and to deny the figurative expressions in honor of physical forces
and mechanical causes. The philosophizing person, who \textit{a~priori} is bent on getting the
development of endowment, the self-activity, the forward-driving natural will, and thereby the
teleological set in full light, readily has a system of wisdom, in which man stands as final
goal, ready to hand; but unfortunately the importunate physics disturbs him: his beautiful
system of wisdom is here and there, in an all-too-striking manner, broken through by
``dysteleological'' facts.

So long as the two modes of view set themselves dogmatically against each other, the
scientific consciousness must find itself in doubt and lack of clarity. It is impossible to
refute the physical arguments, and it is, when one will not renounce all sound sense, just as
impossible to withdraw from the impression of something teleological: what shall we, on
behalf of science, now do? Resolutely reject both the physical and the philosophical
dogmatism, and hold to the fundamental determination of that doubleness of life-activity which
is admittedly seen in the organic function. The reconciliation between the two conflicting
modes of view will then be brought about by physics, chiefly physiology, apprehending the
life-function from the side of substance-activity, while conversely rational biology,
especially psychology, apprehends the self-activity in such a way that it is constantly seen
to be conditioned by substance-activity. By this method teleology can acquire scientific
significance, and the otherwise insoluble problem of the relation between life and organism,
soul and body, will in continued approximation attain a relative solution. The theory of
descent, far from losing its strict natural-scientific stamp, will, by going in under a
teleological illumination, win the recognition to which the grandeur of its researches gives
it title; the natural-historical kinship between the animal and man will be fully understood:
the doctrine of the physical laws of life will be reconciled with the immediate consciousness
of life itself.

\section{Anthropology and Psychology}
\label{sec:20}

\noindent Man is not an animal; this principal proposition ought --- in so far, be it
noted, as the matter of principle is regarded --- to stand in science just as fast as the
one that the animal as animal is not a plant. By being taken up into the animal-nature the
plant-nature is transformed; likewise the animal-nature undergoes a pervasive
transformation by being taken up into the human nature. That it nevertheless comes more
naturally to present-day consciousness to refer man wholly and entirely under the animal
kingdom than the animals under the plant kingdom is understandable from the prevailing
realism. For the transformation that plant life undergoes by being taken up into animal
life is of such a constitution that it allows itself to be demonstrated in the outward, in
the evident reshaping of the organs and functions, whereas the transformation that takes
place with animal life, in so far as it constitutes an integrating part of human life,
belongs more to the invisible than to the visible.

When one begins with the proboscis monkey, the chimpanzee, the gorilla, and thence passes
over to the Papuan, the Tasmanian, to the Australian and the African, in order through the
American and Mongolian human race to reach up to the Indo-German, the Caucasian man
(\textit{Homo iranus}); when one considers that there are still, in southern Asia and
eastern Africa, wild tribes that have no concept of morality, family life, marriage, but
live in flocks just like the apes, climb in the trees and feed on fruits; when one learns
from philology that none of the Australian languages can count to five inclusive; when
physiology and comparative anatomy expressly show how intimately the soul-element hangs
together with the bodily, how much labor of life it has cost in the course of millennia
before the cerebral activity could become articulated as in the higher animals, indeed,
what is more, before an ape-brain could be transformed into a properly usable human brain,
then one gets indeed an incontrovertible impression of continued developments and gradual
transitions.

If we settle into peace with the gradualness in the transitions, then it is just not
difficult to get the human nature to go up into the animal-nature and to reduce the
difference between the animal's and man's consciousness to a difference of degree. If
physiology can really establish that self-consciousness is a physical product of ``cerebral
oscillations,'' a portion of consciousness that especially becomes man's lot on account of
certain disposable cells in the brain: then all that has been said in the foregoing about
self-activity and teleology must be completely retracted. The ape's lively mimicry and
restlessly playing eyes clearly make known that it is attentive to the things around it,
that it --- in physiological terminology --- has cerebral oscillations enough for
``objective consciousness.'' What is still lacking here in all these oscillations? Only a
single turn, that consciousness could be turned inward and become attentive to itself. Were
it to happen sometime as by a lucky chance with a favorably situated specimen, then we
should have here so well-disposed an ape-brain that the choice of quality, with ``utility,''
though ``without end,'' could get occasion to form from it --- even if it cost millions of
years --- a self-conscious human brain. But if we ask --- when self-activity departs, can
the brain then think? Is any man really able to become conscious of his I as a
substance-subject, of his thought-activity as a substance-activity? Were the brain really
the subject of consciousness, how then should one understand such everyday expressions as my,
your, his brain? In order to be able to say: my brain, I must expressly have consciousness of
the brain. If it is now to hold physiologically that it is only the brain that has
consciousness, then it becomes difficult to ascertain whether it is then properly I that am
brain, or the brain that is I. It is a sign of mental confusion to confuse oneself with one's
brain; he who says: I know, I am conscious, that it is not I, but my brain, that is
conscious, must be on the way to the madhouse.

As psychology recognizes the necessity of the organic-physical conditions for development, it
must, for the sake of self-activity, apprehend the development as a development of endowment
and the boundary between animal and man as a boundary of endowment, a boundary of quality.
That both in a human and in an animal development of consciousness there must be countless
gradations, everyone sees. In the child's soul consciousness is dreaming; an idiot can in
development of understanding stand just as many degrees below a clever dog as a Papuan below
a civilized European. Nonetheless the relation between higher animal-understanding and lower
human understanding is of another quality than the relation between higher and lower degree
of human understanding, because consciousness itself is another. Why can the dog's
understanding never be developed into human understanding? To seek the ground in a less
articulated brain, in the lack of language and suitable language-organs, only betrays
misunderstanding; for the ground lies in the bounded endowment. Man has not reached the
higher intelligence because he has gotten a more differentiated brain; he has on the contrary
gotten his brain differentiated to a higher degree because he is endowed for a higher
intelligence. The soul's activity is not a direct consequence of the body's organization;
but the organization is a consequence of the soul's organizing self-activity. The soul
organizes its body, in that it organizes itself in consequence of its peculiar endowment of
life. The animal's soul is endowed only for sense-consciousness with corresponding
understanding; man's soul is endowed for self-consciousness: in this consists the essential
boundary. The self-conscious self, the human I, is neither ``a real of its own'' nor ``an
empty semblance''; it is the most essential, the most substantial unity in which the human
soul comes to ideal existence, an unambiguous testimony of man's marked opposition to the
animal. To be I and to be animal exclude each other; the proposition: I am an animal, is just
as contradictory as the proposition: the circle is a square.

That the doctrine of the development of endowment here set forth by no means conflicts with
the theory of descent can easily be seen. In presupposing a system of lower and higher
endowments, we by no means assume that these endowments should lie independently beside one
another, that the one's stepping-out into existence should be independent of the others'. On
the contrary! If these endowments essentially constitute a system, then they must not merely
develop out of a comprehensive unity of endowment, but even in definite developmental series
condition one another in such a way that the preceding are, as it were, pregnant with the
succeeding: the series of existence, regarded in physical-empirical light, must correspond to
a descent.

What development of endowment means is seen best of all from what we call giftedness, the
possession of particular natural gifts. Because man as man is endowed for becoming conscious
of himself, it does not follow without further ado that self-consciousness, with the whole
inward life corresponding to it, should as by a spiritual sunrise suddenly attain equal
clarity in all. It lies precisely in the essence of self-activity that the soul-life must
work its way forward, that the working-power can be different, and the yield poorer or richer,
according as the inner and outer conditions are unfavorable or favorable; but the endowment
itself will always play a principal role. One is, for example, gifted as an artist, another as
a general. On what now does the giftedness rest? On the brain; for that an artist must have an
artist's brain, a general a general's brain with corresponding nervous system, and so forth,
of this surely no one will doubt. But on what does the brain's peculiar disposition rest? On a
lucky chance! No, chance and science do not agree in this; so one seeks to ground the lucky
concurrence of favorable conditions by pointing to the parents, to the laws of heredity, to
the family, to the tribe, to the state of culture, and so forth --- all correct, provided we
only do not forget the infinite system of real-concepts in which the finite conditions have
their grounds and essential connection (cf.\ \S\,11). And yet there is one thing that is
lacking, the most important, the most decisive, namely the peculiar natural-thought that from
the first germ comprehends the finite conditions of the organization into moments in the
modifying unity of self-activity. It is on this that giftedness essentially rests.

Within the general-human endowment there are a multitude of different endowments, peculiar to
tribes, families, and individuals, with lower and higher natural gifts; the advantages that a
preceding generation, by its energy of development, leaves as inheritance to a following one,
can exercise a considerable influence without therefore making the particular giftedness
superfluous. It is thus understandable that there can be human races, real men qualitatively
different from the animals, whose giftedness is so slight that they perhaps will never at any
time be able to raise themselves up to a higher stage of civilization; indeed it is possible
that the tribes that now bear the culture may sometime be supplanted by tribes which, by
continued improvement of the brain, will be more richly equipped both to continue the handed
down and to begin afresh with something new and original. Only this must be held fast, that
just as little as a moderately gifted individual can work his way up to becoming a genius, so
little can a moderately gifted tribe, by a sum of inherited improvements of the brain, without
a breath of a higher spirit of giftedness, work its way up to becoming a productive bearer of
culture. The inherited and the acquired ought, as regards development, not to be reckoned too
low; but a sum of conditions, be it ever so great, is yet not the original cause: only when
the conditions become moments for the nature-powerful self-unity of the endowment do we have
the higher giftedness.

A still more concrete distinguishing mark of qualitative difference between animal and man is
conscience. It holds here, as with consciousness, that one neither misjudge nor overvalue the
analogies, but, amid the gradations, get an eye for the quality. Conscience is an instinct of
self-preservation of a peculiar kind; involuntarily it expresses itself as a feeling of the
individual's dissatisfaction with itself. Although the animal has no self-consciousness, the
higher animal has yet an immediate conscious sensation of itself and can thereby furnish an
analogon to conscience. A well-trained dog, for example, can give touching signs of
dissatisfaction with itself when it has done wrong; it is not merely the beating it fears, no,
it is its master's displeasure. It is ashamed, it expresses regret, it begs forgiveness: it
really seems as though it had conscience. And yet there is the same distance between this
analogon of conscience and conscience as between the animal's sense-consciousness and man's
self-consciousness. As the self is, so must the instinct of self-preservation be. In
conscience self-consciousness goes, in an immediate, naturally necessary manner, together with
the existing self. From conscience it is seen that man's I is more than a rainbow-gleam of
thinking cerebral oscillations, more than the formal center of ``apperceptions,'' of senses,
representations, thoughts. There is in conscience an involuntary, nature-powerful indication
of what offends and destroys, and what satisfies and sustains, the self-conscious self-being:
in conscience, apprehended as ideal nature-power, the self is in earnest.

When one, from the standpoint of a reflected-in ethics, apprehends conscience without further
ado as an innate sense for right and wrong, good and evil, a holy voice, and so forth, then a
misunderstanding can easily hide in this. If one determines the good as that which sustains,
the evil as that which destroys, the self, then one can indeed say that it is conscience that
in the highest instance must be judge over good and evil. But, seen from the side of nature,
man's immediate impression of what is good and what is evil is conditioned by the stage of
development on which consciousness stands. What the one's conscience recognizes as good, the
other's can recognize as evil; the judgment is so subjective that the only thing on which it
can be said unconditionally to depend is an inner self-preservation. In the natural sense the
instinct of conscience is, among all instincts, the most egoistic. The wild chieftain murders
and scalps his enemies. He has perhaps no word for conscience, and yet it is to him a matter
of conscience; if he did not do it, he would seem to himself to be wretched and cowardly and
contemptible in his tribe. When a Jacques Clément and a François Ravaillac commit regicide, it
is in the Jesuit sense for the sake of conscience that they can earn eternal blessedness. There
is in reality nothing so repulsive, so dreadful and inhuman, that it may not be a prompting of
a human conscience.

It is a hard contradiction that what constitutes the nature-kernel in man's self can express
itself as the most revolting unnaturalness; and yet this is precisely the right sign of man's
difference from the animal. In and with the self there is posited in immediate
self-consciousness an unconditioned. The unconditioned expresses itself naturally and yet
belongs not to nature, for it belongs to spirit. The contradiction is that human nature does
not permit man to settle into peace in a kind of good-natured animal life's natural state.
Conscience is no soft-boiled darling, but a bitter radical. To break through the hard shell the
nestling's beak gets a sharp point, which falls off when the shell is broken through; in a
corresponding manner the instinct of conscience in the naturally self-conscious self has, so to
speak, a cutting barb, which first falls off when the hard shell of rawness is broken through,
and the true human nature, endowed for the ideal of the good, dawns upon consciousness.

A just as telling sign of the essential difference between animal and man we have in
personality. The animal-nature brings it only to the point of setting an ensouled individuality
in existing individuals, specimens of the species; the human individual is more than a specimen
of the species; it is a self-valid individuality determined for right and freedom, a
personality. There is still in modern civilization a remnant of barbarism that science must
most powerfully counteract. While sensitive souls, in touching recognition of the hypothesis of
descent, have already brought it so far that they plainly declare the animals to be ``our
brothers and sisters,'' the less sentimental ones --- as the continually re-emerging slave-trade
bears witness --- have not yet been able to bring it to the point of discovering the personal in
the individuality of the negroes and the colored peoples. Of course, if the struggle for life
literally goes according to the laws of mechanics, it might seem as though ``the choice of
quality'' with ``the theory of utility'' gave them a certificate for standing on the heights of
researching science; but science is not for slave-traders. The truth-seeking researchers will,
when they reflect on the practically-consistent application of a physical-mechanical development
of life, which is to rest solely on the right of the stronger, yet surely grant that in the
determination of the original character of the life-functions an essential factor has been
forgotten.

Self-activity, the inner self-determination with clear consciousness of the goal of
development, is the fundamental factor of personality. Personality has its nature-ground in
individuality, but its principle in freedom. A nature-powerful individuality is yet no
personality. Individuality is something innate, but real personality is something
self-acquired. Nor is personality without further ado the same as character; fixed, prominent
marks of individuality give character, whether these marks and lineaments otherwise belong to
the innate or to that which has been adapted by the will: animal individuals too have
character. The demand of personality is that natural promptings, drives, and passions be
subordinated to the will, worked through, appropriated, and transformed in freedom. In
opposition to the animal, man is in principle, according to his endowment --- personality.
Man's personality is no immediate natural gift; it is a gift of spirit, a task of freedom: man
is essentially spirit.


%%% PART THREE: THE SCIENCE OF SPIRIT %%%
\part{The Science of Spirit}

\chapter{The Human Spirit}

\section{The Natural Ground of Spirit: Necessity and Freedom}
\label{sec:21en}

\noindent Life, soul, and spirit have in the natural ground a common original endowment.
The manner in which the difference steps out into particular existences rests upon
essential boundaries of endowment. Plant life exhausts itself in organic activity of
form-giving --- nutrition, growth, and propagation --- without reaching the independence
that, in selfhood, presupposes a self. For the animal organism the endowment of life has a
surplus of ideality which, in inner unity of self-sensation, sensation, and voluntary
motion, forms a soul. The ideal liberation of an essential self, which dawns in the
animal's soul, gains, in virtue of a peculiar endowment, clarity, independence, and
liberated inwardness in that form of self-understanding, self-working, and perfectibility
which marks the unambiguous superiority of the human soul over the animal soul. On the free
inwardness depends the consciousness of the ideal: only in the light of ideality do the
universally valid laws come to light. As thoughts, concepts, ideas, with great questions
about the ground of things, about the worth of life and the goal of existence, force
themselves forward in the free interior, the human soul can apprehend its essence as
spirit.

For the cognition of spirit as spirit science knows no other way than the one through man's
soul; every attempt to reach higher up by choosing a higher starting-point shows itself at
once as unscientific. By going up in a balloon one can get a free survey of a city, a
region, a landscape, and see the whole in bird's-eye perspective; but the spirit of science
knows no balloon in which one can raise oneself above the scientific and, with the supposed
higher presupposition, swoop down into the regions of research wherever one pleases. Shall
we then, in the name of science, declare the human spirit to be spirit itself in absolute
originality? By no means! How the relation between the human spirit and the absolute spirit
must, according to our mode of view, be apprehended, can be seen from the foregoing.

In order to reflect on the foundation of science, it is necessary to go back to the original
unity of being and knowing that is expressed in the Idea of Truth (cf.\ \S\,11). To the
subjective proposition: self-being is in the absolute sense one with self-knowing, there
corresponds the objective one: the being of other is indissolubly one with the knowing of
other. It is not difficult to see what in this is indefinite because it is indeterminable,
and what is incontrovertible, determinate, because it is an absolutely indispensable.
Indefinite and indeterminable is the manner in which the eternal self that bears all being
and knowing is to be fixed and held fast in inner independence. Attempt it by thinking,
representing, intellectual intuiting, however one will, it is and remains inconceivable.
Every personification, however supple, of the infinite self is to science an offense.
Likewise every attempt at a closer determination of how the inductive and distinctive, the
simultaneous and yet successive knowing of all is to be distinctly thought, however much
dialectic is lavished upon it, is wholly unfruitful. The only thing that can be attained is
a thorough rejection of the critique of understanding which, by setting the inductive and
distinctive, the simultaneous and successive outwardly against each other and crying out
about contradiction, would dogmatically abolish the presupposition of knowing's unity with
being. On the other hand, it stands unshakably fast that thinking without knowing is a
mystical fantasy, knowing without self-knowing a logical abstraction, concepts without
comprehending and comprehending without a comprehending self-unity a meaningless assumption.

As the unity-expression for all reality, all that is actual in existence, we have the
category Power. To ask whence power comes is just as idle as to ask whence being comes. That
about which alone one can ask is the original form under which power must show itself when it
is to be thought as power. In the circle of changing phenomena power shows itself from
countless sides as power in other, that is, as the world of realities. Since now, as
reflection makes clear, the phenomenon is one with and yet opposite to the essence, namely
the other-essence, and the other-essence both different from and yet grounded in the
self-essence: so the power in other, the broken power, the power of manifoldness, must be
grounded in the self-power, the sole power, which does not subsist outside or beside the
manifoldness of the realities, but has all power in other as a real-moment in its infinite
self-confirmation. That the self-power is ideal and, in the form of ideality, one with
self-knowing, lies in the fact that it is infinitely reflected; that the power in other is
real and, in the form of reality, identical with the existing thing, rests on the fact that
it has the immediate for its outburst. In the apprehension of this too we must distinguish
between something that remains indefinite and something that has a claim to scientific
recognition.

The fundamental relation between the real power of manifoldness and the ideal unity of
self-power does not allow itself to be made distinct by any articulating analysis whatever.
How should it be possible for sensing beings, who objectify in impotence, to reach insight
into the manner in which the creating self-being objectifies in omnipotence? The infinite
reciprocal determination between the ideal world of real-concepts and the real world of
things does not allow itself to be referred under a finite understanding's particular How.
But the decisive principal points of the presupposition can nonetheless stand luminously
clear before consciousness. That the logical in the real-concepts and the anti-logical in
things must punctually correspond to each other is more than a logical assertion; all the
sciences treated in the foregoing, the empirical as well as the \textit{a~priori}, have
furnished proofs. The interpretation of the reality of matter that follows from this is, if
we do not err, the only one which, in full agreement with physics itself, can bring the
materialist myth of ``matter and force'' out of science.

It can be granted that matter with its forces is from eternity, for the power in other is
just as old as the self-power; but it cannot be granted that ``matter and force'' constitute
the absolutely original, for the power in other --- the existence of matter --- has from
eternity its origin from the self-power and thus belongs to the derived, the grounded. It can
be granted that the natural elements, when the division goes to the outermost boundary of
reality, are existing atoms, ``ether-atoms,'' ``mass-atoms,'' and that these are of eternal
origin; for all finite existence is subject to discreteness, and the atoms furnish the
existential expression for the real discreteness of the power of manifoldness. But it cannot
be granted that the discreteness should have so absolutely exclusive a character that the
existence and constitution of the one atom thereby kept itself in a contingent relation to
the constitution and existence of the other. Imagine the swarm of atoms as one will; augment
the forces or lay them into the properties; begin with homogeneous atoms, hydrogen-atoms for
example, in order to derive disparate ones by the help of outward, contingent conditions;
resolve the atoms into forces, the forces into a fundamental force, in order again to get
several forces out of the one force; make oneself, in short, as profound, acute, subtle as
one will, with outward representations, hypotheses, and reasonings: the whole is nonsense and
remains nonsense.

In order to come to clarity about what essentially belongs to the foundation of science,
there is, so far as we can see, only one single way out: namely, to lay the identity --- and,
of course, also the opposition --- of self-power and power in other into the original
presupposition. The physically acting causes, which from the side of reality correspond to
the nature of things, are not either weakened or overlooked, but much rather recognized,
indeed first fully recognized, when it is seen that they, from the side of ideality, have
their ground in determining real-concepts. What the real-concepts are in the unity of
self-power, that is precisely what, in the manifoldness of the power of otherness, expresses
itself through outward causes, substances and forces, and outward conditions. That everything
in the world of things --- the living organisms as well as crystals and strata --- is subject
to conditioned necessity by no means conflicts with, but becomes precisely intelligible from,
the fundamental presupposition that the whole is an endowment of self-power, a real
development in the circle of facts and laws; thus, essentially seen, a development of
endowment, and the development of endowment an infinite satisfaction of the demand of
selfhood, hence teleological.

These fundamental determinations of Knowledge and Power in original self-unity are, however,
by no means sufficient to secure for us the fundamental concept of the absolute spirit; and
to let imagination fill out the emptiness and enliven the pale features by a personification
is only to encumber science with ambiguities. With respect to spirit's relation to nature
there are indeed profound investigations carried on concerning immanence and transcendence;
but the transcendence does not reach out beyond the dialectical reflection of the self-unity's
essential difference from the manifoldness-content with which it is identical. How unattainable
spirit's reflection into itself, its deep inwardness, its absolute independence is for human
knowledge shows itself clearly when the relation between necessity and freedom is more closely
investigated.

The comprehensive principal propositions, decisive for the foundation, about self-knowing and
knowing of other, self-objectification in objectification of other, self-power and power in
other are at once propositions of necessity and propositions of freedom. That they are
propositions of necessity no one shall be able to deny. If self-knowing, self-objectification,
self-power are thought away, then knowing of other, objectification of other, power in other
is simply impossible; self-knowing is thus necessary for knowing of other, self-objectification
necessary for objectification of other, self-power necessary for power in other. But the
reverse holds too. Without knowing of other there can be no self-knowing, and the same holds
with respect to the objectifying power. Knowing of other is thus necessary for self-knowing,
objectification of other necessary for self-objectification, the power in other necessary for
the self-power. From this it can be clearly seen both what necessity is and how the absolute
necessity stands related to the relative.

We have in the relation two terms, a subjective $S$ and an objective $O$. Neither of these terms
is in and for itself necessary; they are necessary only for each other. But how decisive this
necessity is, is seen from the fact that when the one term vanishes, the other vanishes too:
they must either both be or not be. The difference between absolute and relative necessity rests
upon the difference between the infinite and the finite, the original and the derived. In so far
as $S$ is infinite as being itself, and $O$ an existence in infinite becoming, we have infinity
in two forms, that of eternity and that of time. $S$ is infinite in the form of eternity, for it
circles in itself through $O$; $O$ is infinite in the form of time, or rather of temporality; it
is in essence without beginning and end, and yet every existing term in the endless series is
finite, and every series concluded in existence has in time a beginning and an end. From this it
is seen that both terms, although they are in time, so to speak, equally old, are yet not equally
original. $S$ is from eternity the absolutely original, $O$ the derived.

This difference cannot without further ado be carried over to necessity. If we let $S$ represent
spirit, $O$ nature, then one cannot, in consequence of our foregoing, say that only spirit is
absolutely necessary because it is the original, nature on the other hand merely relatively
necessary because it is the derived; for $O$, the derived, is, as we have seen, in the relation
itself just as absolutely necessary as $S$, the original. If, on the other hand, we hold to the
finite terms existing in $O$, and let the one furnish the premise, the other the consequence,
then one can with truth say that their connection expresses a relative necessity.

A series of conditions $B_0$, $B_1$, $B_2$, \ldots $B_n$ ($n = \infty$) will, when we look away
from the connection, give us an infinite multitude of contingent terms; in consequence of the
connection of conditions, on the other hand, we have a series of conditioning terms, necessitated
by the conditioned ones, that is, a series of terms of conditioned necessity. If $B_0$ is, then
$B_1$ must necessarily be. The necessity of $B_1$ is then conditioned by $B_0$'s contingently
existing. Thus the contingent and the necessary continue to alternate with each other up to
$B_n$, which, in consequence of the infinitely many preceding conditions, should be infinite,
absolutely necessary. But that there is no greater necessity for $B_n$ than for $B_1$ is evident
when one considers that the whole chain of conditionally necessary terms hangs on the one
contingently existing $B_0$. The conditioning and the conditioned are, despite the connection,
outward to each other; on this outwardness rest the finite relations, on the finite relations
again rests that reciprocal determination of the contingent and the necessary which characterizes
relative necessity.

The absolute necessity can just as little be without conditions as the relative; but while the
conditions in relative necessity are outward, they are in the absolute on the contrary inward.
The existence of $O$ is conditioned by $S$, and the being of $S$ conditioned by $O$; but since the
existence of $O$ is inseparable from the being of $S$, $S$ has the condition for its being in
itself, and since the being of $S$ is likewise inseparable from the existence of $O$, $O$ too has
the condition in itself. It is on account of this entering of the conditions into the unity of the
relation that we call the necessity unconditioned.

That the absolute relation of necessity is at the same time a relation of freedom follows simply
from the essence of selfhood, of subjectivity. Whether we apprehend $S$ as knowing or
objectifying, unconditionally mastering its real objects, the determining activity must
necessarily go together with a self-determining: unity of knowing and self-determining is
essentially freedom. In a development of the concept of freedom it is highly misleading to begin
immediately with will and freedom of choice. Will and freedom of choice are, in the proper sense,
categories that cannot be taken up into the foundation of science; they belong, strictly speaking,
under arbitrariness, but the arbitrary is so subjective that it withdraws from every possible
objective knowledge. If $O$ were posited arbitrarily by $S$, then necessity would be denied, and
the relation of $S$ and $O$ would not allow itself to be determined scientifically. The non-being
of objectivity would then be just as possible as its existence; in other words: objectivity, the
objective world, would become quite contingent. In such a contingency all nature and natural
ground has vanished.

In knowing self-determining, freedom can be united with absolute necessity only when
arbitrariness is excluded; a higher concept of freedom is unattainable in science. Is now the
concept of spirit and freedom hereby exhausted? By no means! The presupposition that $S$ with
absolute necessity produces $O$ contains no closer determination of how $O$ is constituted. $O$,
the real-concept of objectivity, of the objective world, can indeed, by an outburst of power, be
thought realized in different worlds: $O_1$, $O_2$, $O_3$, and so forth, and the anti-logical
impress on each of the existing totalities a peculiar stamp. That this nature, precisely the world
in which we live, is not the only possible one, and consequently not absolutely necessary, physics
and ``positive philosophy'' confirm in an indirect manner by holding exclusively to facts, that
is, to the factually given.

If we regard from this standpoint the unity of freedom with the absolute necessity, then we
discover several possibilities. It is indeed possible that the absolute spirit originally had a
choice between several worlds and with free decision produced precisely this one, in which we find
ourselves. Yes, it is possible; but this possibility science cannot use; for of such a thing there
is no objective knowledge. It is indeed possible that this our world, in preference to other
possible ones, has been realized, not in consequence of spirit's free choice, but by a kind of
chance; a world had to be produced, and so it became contingently this one. But such an assumption
has no significance in science either. Far from making the original self-determining recognizable
and thereby reconciling freedom with necessity, it much rather sets out to set chance in the place
of arbitrariness and to exalt blind facticity to the rank of ``absolute necessity,'' ἀνάγκη.
There remains then only to grant a gulf of possibilities between $O$ as real-concept and the
realization of $O$, the real, objective world.

It is the relation of necessity between the all-determining self-unity and the objective
determinations of the world, found expressed in real facts, to which we must hold. What lies before
us is an order of things with atoms and masses, forces and laws. These realities are inconceivable
to us without corresponding real-concepts; the real-concepts are impossible without a corresponding
comprehending; but such a comprehending is not thinkable without an all-comprehending,
all-objectifying self-unity. In the infinite comprehending and objectifying of power there lies a
self-determining, and in this freedom; but the objectifying takes place with absolute necessity, and
we thus come, in our cognition of freedom, not a hair's breadth outside necessity.

We cannot speak of spirit and nature as though spirit were outside nature or nature outside spirit.
The unity here treated is a self-unity that essentially belongs to nature, and without which nature
cannot possibly be. So far self-knowing is then also nature-knowing, self-comprehending
nature-comprehending, self-power nature-power; but, be it noted, not from the standpoint of reality,
but from that of ideality. Ideality and reality are indeed two sides of the same thing, but these
two sides are also two totalities in one totality. The ideal totality is the spirit of nature, the
real totality is real nature. The self-unity is thus the unity of spirit, self-knowing the knowing
of spirit, the self-power the power of spirit. Whether one will call this view theism, naturalism,
dualism, or monism is quite indifferent. The point is exact determinations that can and must be
taken up into the foundation of the science of spirit. The question on which everything depends is
only this: how far can science reach in the determination of the fundamental relation between concept
and reality, self-knowing and knowing of other, self-power and power in other, or, otherwise
expressed, between freedom and necessity, spirit and nature.

A proposition such as this: nature has its origin from spirit, because spirit has its ground in
nature, shows how impossible it is to determine the concept nature in such a way that the word gets
only a single meaning. Those who lay the whole of nature into matter and force are obliged to smuggle
in so many secondary determinations of the logical, the intelligible, the rational, indeed, when we
look more closely, even of secret teleological striving, that the concepts ``matter and force'' bulge
and swell and are near to bursting. Those who make nature one with the sum and connection of the
things of nature can easily be convinced that there is a difference, indeed an opposition, that they
would keep hidden in the immediate unity, but which yet everywhere crops up and speaks of a lack of
reflection. Those who assert that life is a ``mechanical'' product of lifeless nature have a
self-activity hidden in nature, much as one who does tricks with the pocket-compass and the little
ends has a magnet hidden in his sleeve. The word nature is ambiguous and must necessarily be so,
because nature itself is a double-sided being. On the one side, spirit in self-knowing and
self-determining cannot possibly be spirit without having a ground; but the ground, with all spirit's
real possibilities, must in its first immediate form be ground for spirit's nature, hence
nature-ground. On the other side, nature in its realities as real nature with matter and force cannot
possibly be objective without being objectified. If it now owes its objective existence to an original
objectification, then one cannot get free of spirit by the assertion that nature is objective because
it has objectified itself, for the selfhood in this ``itself'' is precisely spirit: nature has its
origin from spirit.

Under this, and only under this, presupposition is it intelligible that nature, with its endowments
stamped by spirit's self-determining, can in the lifeless begin with masses and forces, and yet in a
continuous development from the lifeless through the living and the ensouled organisms end with a
development of organically-conditioned spirit, that is, with the human spirit.

\section{The Self-Development of the Human Spirit: History}
\label{sec:22en}

\noindent As a special science history is not merely, like anthropology, a
teleological-empirical science, but becomes also, since its essential content is a content
of spirit, especially to be determined as an ideal-empirical science. While the facts of
nature, with the stamp of conditioned necessity, always lie ready for closer scrutiny, so
that the doctrine of nature can always have its object present, history is on the other hand
referred to the past and cannot, however carefully the research is carried through, reach
immediate autopsy. Although the historian can build on a foundation of natural laws, his
material is yet so modified by spiritual freedom and crisscrossed by arbitrary endeavors
that it cannot possibly be treated simply as a product of conditioned necessity. The
historical facts can then not be proved in the same manner as the physical and
natural-historical ones either. The proof of any past event whatever is, strictly speaking,
only a proof of probability, and presupposes a certain faith in source-writings, documents,
and reliable authorities, a historical faith (\textit{fides historica}). The critique has a
far easier time awakening doubt about the single constituents, form, and connection of what
is handed down than securing for itself the positive kernel of the content. If we ask what
it then properly is that upholds the scientific authority of history and prevents skepticism
from gnawing through and devouring its foundation, then we must especially note the
following:

Probability has many degrees, and the historical proofs of probability would indeed hover in
the air if they had no substratum that, instead of a mere probability, gave apodictic
certainty. That the present has a past is more than probable, it is decisively certain; that
between the life and life-development of the preceding and the succeeding generations there
must be an essential connection cannot be doubted either; the law is here at once a natural
law and a law of history. One cannot indeed, from states and events in the present, calculate
with mathematical exactness which states and events must belong to a definite past, and
traditions do not help when one will not believe in them at all; nonetheless the impression of
tradition will, in great and distinct strokes, furnish maxima of probability that can be
regarded as the ineradicable radicals of historical certainty. About these radicals the
degrees of probability for the single events then group themselves in such a way that the
whole and its particulars melt together in trustworthy agreement; the historical presentation
will then be like a painting with foreground, background, and middle ground, and by the
distribution of light and shadow give a living, colorful, characteristic picture of the spirit
of the past.

The spirit of the past is always different from that of the present, but therefore it is yet
not two, or, according to divisions of the past, several spirits that we have to do with; the
spirit different in the different times is the human spirit.

The scientific in historical knowledge will then especially be recognizable by the
apprehension of the relation in which the different developmental stages of spirit are set to
spirit itself. As an empirically researching science history must lay decisive weight upon
actuality, hence upon the phenomena of actuality in which the spiritual forces of the race have
made themselves known, and can in no way content itself with reflected forms of general
development of spirit. The strife that hence arises between learned history and the philosophy
of history seems, however, as regards the principles, to be capable of being settled without
difficulty. The historian holds strictly to what has happened, absorbs himself in events,
actions, circumstances, and takes along even the smallest particulars, because he knows that
the real spirit can come to view only through all this; the ensouled whole consists of sheer
particulars: \textit{les détails sont l'âme de l'histoire} [the details are the soul of
history]. The learned historian is tempted to regard the philosophy of history more as a
fantastic attempt to hypostatize the general spirit and let it act in the place of the
individuals than as a scientific demonstration of the spiritual yield brought about by the
common exertions of the individuals. --- The fundamental question, whether it is really spirit
or the human individual that is the acting subject in history, allows itself, however, to be
answered in such a way that learned history and the philosophy of history each get their
scientific task. It is the relation between real-concept and existence that again plays the
principal role. Seen from the side of existence, it is really the individuals that act; but if
there were nothing else in the action than the immediately individual, then the actions would
sink down to drive-like expressions of life: there would be events in quantity, but no history.

In order for a deed to deserve the name of action, it must proceed from a willing self that is
conscious of its meaning and end. It is not deeds, but actions, that in the essential sense
produce history. The animal cannot act; only man can act, for man is spirit. If we begin with
the race on the very lowest stage, then we find hordes of individuals in whom self-consciousness
dawns and shines in strong flashes. Let the outward likeness to an animal flock be ever so
great, the qualitative difference will yet come to view. The raw human individuals that hunt,
fish, struggle with the animals and with one another, have yet in their essence not merely a
part of human nature, but the whole human nature as original endowment. Late generations,
which precisely by having advanced on the way of development feel themselves oppressed by the
piecework of civilization, look with longing back upon the simple first and dream of a golden
age, a state of bliss, when nature was spirit, and spirit nature without parceling-out. It is
indeed a dream, but a dream of the waking spirit, a dream of historical significance. History
begins not with a golden age, but rather with a stone age; and yet the ideal that a race
assigns to the golden age furnishes a sure measure of the height of spirit on which it stands,
for it is the human essence with its original endowment that comes to view in the reflection of
the golden age.

The mirror-image of spirit that represents the beginning is in reality the future goal that, by
the exertion of forces under unceasing struggle and labor, is to be reached in approximation:
the movement in history is more than ``mechanical,'' it is teleological. The teleological
becomes, however, overscientific and just therefore unscientific when one would, in the course
of events, demonstrate the governing providence and make history into a theodicy. Here too it
holds that development of ends is a development of endowment, only that the human spirit takes
freedom along in conscious, nature-bounded self-development. We shall now see how the
development can take place, in that every man who participates in it is at once individual
subject and spirit-subject.

We let then the human hordes with which history is to begin stir in all naturalness, without
order of society, without marriages, without moral or civil ordinances. What will now happen?
The individuals flock together because they need one another mutually; the animals do that too.
Each individual is driven by natural egoism to take what it can get, and thereby comes into
strife with the others. Among the animals likewise. In the flock there will arise a difference
between stronger and weaker; the stronger will rule, the weaker must subordinate themselves, and
so forth. These analogies between the human and the animal can be carried very far, provided one
will be content with the superficial. But if we take the human essence more sharply into view,
then there shows itself at once the peculiar element that stems from spirit. Language, the free
common designation of sense-images, the ideal embodiment of representations and thoughts, begins
to form itself. In and by language the immediately individual is transformed into something
common and generally comprehensible, and this common, generally comprehensible element, in which
the individuals understand one another, forms a common consciousness with the first daybreak of
spirit. The single one's consciousness of itself, of its own needs and its own demands, is from
now on inseparably united with a consciousness of others' needs and others' demands. The
individual, however egoistic it may be, can no longer isolate its immediate drive-will from
regard for the others' will; there begins to form a common will: the first, weak intimation of
will-of-spirit expresses itself in the flock, and the flock is transformed into a society of men.
The elementary society may in the first draft be raw and loose, but the impulse of formation
stirs in spirit, and fruitful germs shoot up from the ground of endowment. An ordered family
life, a kind of legal distribution of mine and thine, corresponding to working-powers and works,
will obtrude as of itself; society is founded, the deeds of life ordered, laws given, before one
can yet name any founder or orderer or lawgiver: whence then does it come? One thinks back, one
guesses, recounts, and reflects; legend names now this one, now that one, but all the figures from
the first beginning are veiled in mist: they have likeness to men, and are yet superhuman.

When the first ground is laid, the struggle can begin in earnest in the life of society. The many
wills arm themselves; it looks like a war of all against all. But war alone does not build up
society, it brings only destruction; there must, under the war and the enmity, hide something
positive and sympathetic that can reduce the struggle to a condition, bring reconciliation, and
fill the activity with richer content of life. The struggling individuals must gather into groups,
each of which has its positive end. But whether we now regard the contending groups singly or all
together, we see how the single wills are everywhere obliged to subordinate themselves to a ruling
general will. Every horde must have its leader and feel the necessity of obeying him and following
him. One might perhaps suppose that many single wills were hereby immediately melted together into
one single individual single will, namely the leader's will; but that this is not so can easily be
seen. It is, after all, a feeling, a common desire, a common goal, that unites the horde and the
leader. If the leader merely followed his egoistic self-will, he could not possibly compel the
horde to follow him. The strife can end with victory, and the victor become tyrant; but in
subjugating the single wills the tyrant does not annihilate them, he only forces them in under the
general will. Although the general will thus compelled goes up into his egoistic self-will and can
be misused in the most shameful manner, he must yet fear its power, for it can suddenly turn against
him and crush him. A state of society can for a time be such that all the single wills are at bottom
in revolt, and yet the general will is obedient; and conversely such that the egoism in the
individuals is inclined to peace, and yet the general will wills insurrection.

The general will is will-of-spirit; it is not in the individual wills, but in the will-of-spirit,
that we must seek the proper driving-wheel of historical events. From this it can be discerned how
it stands in history with the multitude and the single one, when the question is of giving an age
its direction, its upsurge, and spiritual stamp. Those who assert that it is always the single ones,
the heroes, the champions, the geniuses, who accomplish the historical, are so far right, in that it
is in such richly equipped individual subjects that the spirit-subject, the active, acting self-being
of the general will, can be said to be centralized. Those who lay decisive weight on the popular
sense, on national spirit and national will, are from their standpoint also right, in so far as
national will is essentially general will, and the national spirit the bearer of the general will.
But the will of the multitude is usually spiritless; one cannot rely on it as a general will, for at
bottom it is only a sum of contingent individual wills held together by the mood of the moment; nor
can one regard it as a complex of merely egoistic single wills, for a breath of spirit from the
general does pass over it, and it can as drive-will bear something essential: the misuse of the
ambiguous will of the multitude, which under the most varied forms --- to the advantage now of
monarchy and hierarchy, now of republic and democracy, and always to the advantage of tyranny ---
repeats itself in history, is a distortion of spirit, a lasting cause of the peoples' sufferings.

From the general will apprehended as will-of-spirit one will then also be able to understand the
relation between the divine and the human in history. From the standpoint of the science of spirit
it becomes necessary to consider this relation more closely. In the circle of phenomena the matter
presents itself thus. In consequence of its divine right the human race regards itself as standing
under the special guardianship of the Deity; it is the pervasive theme that among the different
peoples is varied in the most varied ways in myths and legends. That the folk-religions and the
peoples' natural peculiarity on the whole correspond to each other, that the gods as bearers of
culture go before on the way of development, and that every people which has reached the boundaries
of the old gods' spiritual reach must either live on in the repetition of a dead tradition or be
visited by higher, truer, richer gods, in order through these leaders to come further forward, is
a historical fact. What do these god-revelations mean, and how do they stand related to the human
spirit? For the folk-consciousness it presents itself as though these gods came from without, or
rather from above, from heaven down to men. When the critical understanding begins to dissolve the
representations of the divine beings' outward independence, the folk-faith is in danger. On the
other hand, the folk-faith is willing to declare the old gods to be idols, false gods, when it has
found the true God, the God of whose revelations no one can doubt. A confirmation of the truth of
these revelations man can indeed find in his own interior, his own heart; but this is yet only a
derived testimony: the original proof, decisive for the folk-faith, is historical. Such a
historical proof cannot, however, in the nature of the case, be given: here is a problem of the
science of spirit.

That the figures of the gods come from within and yet present themselves as though they came from
without, so that revelations of every kind take place within the human spirit's own boundaries,
points to a peculiar tension between all that breathes in the overflowing content of the endowment
and that which the spirit-poor actuality offers. Man is under this tension higher than himself,
stronger than himself, richer than himself; what stirs in him, thinks in him, speaks to him and sets
itself before his soul's eye, is indeed really the divine from the ground of the self-being, and the
revelation has so far its truth. But imagination does not permit reflection to come forward with its
interpretation; it sets in its way an insurmountable gulf between the divine and the human, in order
again to fill out the gulf in its way; and the power of imagination is irresistible: it works hidden
in deep feelings, speaking images, words, and thoughts. While man immediately and naturally
apprehends the single will as his will, and despite all egoism must recognize its boundary and
impotence, he on the other hand lays the great, heavy tasks of the general will trustingly in the
hand of the mighty gods, believes in their wisdom, brings them his sacrifices, and relies on their
protection.

When history, then, is apprehended as the phenomenal time-expression for the self-development of the
human spirit, the thought of divine will, divine help, and divine governance is therein included;
but that the development must be a self-development, and that this self-development, which takes up
moments of relative necessity and relative freedom, subsists together with natural barriers and
outward hindrances, is just as certain as that spirit itself is the goal of the development. Society
with its organization, its labors, products, and relative ends is the realm of actuality in which
the personal formation of human individuality takes place, and this realm of actuality is in the
name of the general will a realm of spirit. It is only a fragment of history that we can survey; but
already this fragment bears sufficient witness to that reciprocal action of nature and spirit on
which the historical advances essentially rest.

Three principal stages are already passed. The unclear unity of natural and spiritual, human and
divine, that on the whole characterizes the life of antiquity, marks the first stage; the sharp
separation and strongly prominent opposition of nature and spirit, sensible and supersensible,
worldly and otherworldly, which from so many sides was lived through and worked over in the Middle
Ages, marks the second stage. The present's involved task of maintaining the unity without effacing
the opposition is indeed still very far from its solution, but the center-point is yet on this third
stage already found; for the center-point about which all movements, all endeavors --- religious and
worldly, material and spiritual, practical and theoretical --- essentially turn is man himself,
the personal man. We can with respect to this divide history into ages: the age of myths, the ancient
history; the age of churchliness, the history of the Middle Ages; and the age of the rights of man,
of humanity, the more recent history.

From the standpoint of the science of spirit it is especially of importance to investigate how
posterity in each generation apprehends its relation to the past. We do not here think of the more
or less superficial one-sidednesses, that some are inclined to disparage their own time and through
the spectacles of imagination see sheer splendor, sheer examples of the noble, the good, the exalted
in bygone times; while others, on the contrary, taken with the advantages of the present, disdain
bygone times' customs and institutions to such a degree that they are not able to penetrate into the
spirit and essence of the past. Nor shall we dwell on the apprehension of those who, without an eye
for the teleological in the life of freedom, lose themselves exclusively in a pragmatic presentation
of the irrefutable necessity with which a following time's conditions are bred from the conditions in
a preceding one. That on which it essentially depends is a true valuation of the inheritance that the
past leaves to posterity; on this rests the true insight into what has abiding worth, and what must
disappear because it, in merit as in error, belongs only to the changeable. The Middle Ages, with
their one church-faith that alone makes blessed, had little or no understanding of the great past,
which it had yet cost centuries of colossal labor to build, the substratum of civilization that was
indispensable for its elementary-personal life of society. It regarded the spirit of antiquity as
the spirit of darkness, while it, with ascetic suspicion of the natural, supposed itself to live in
the supernatural light of revelation. It had gotten its eye open for the opposition of natural and
spiritual, for the will-of-spirit's supremacy over the immediate drive-will; but it misjudged
spirit, because it disdained nature; it lived in a spirit-consuming enmity between natural and
supernatural, because it was not able to grasp the abiding element in antiquity's view, that nature
is eternally the ground of spirit.

The more recent time has a sharp eye for the shadow-sides of the Middle Ages; out of offense at the
critically notorious enmity between the flesh and the spirit, it is in the act of skipping over the
Middle Ages, declaring them an age of darkening, and in the name of human nature, humanity, and
truth seeking the original points of attachment in classical antiquity. This ``modern''
one-sidedness is in the historical sense just as great as the medieval. The point is a deepening of
man's nature, of the human spirit; it is unmistakably a great advance that the stage of consciousness
is already reached on which it is clearly recognized that the human spirit, through everything --- the
highest and the lowest, the supernatural and the natural --- has in the last instance to do with
itself, with itself alone. But the more one absorbs oneself in tasks belonging here, the more clearly
it will be seen that man's essence is dualistic, that the unity which is to bring peace and
reconciliation between the conflicting elements lies deeper than classical antiquity could reach.
Precisely for the sake of unity, of personal unity, the human spirit must, when in full
self-understanding it is to come to itself, go through the absolute spirit; but this way is not
historical, it is over-historical.

\chapter{The Self-Liberation of the Human Spirit}

\section{Self-Liberation in the Medium of Fantasy: Aesthetics}
\label{sec:23en}

\noindent The goal of spirit's self-development is its self-liberation. Whatever advantages
the historical advances bring with them, the roots of evil and good are yet so intimately
grown together that pressure and want and suffering must follow the race on all the stages of
history. Can the human spirit now really be said to feel these pressures, these wants and
sufferings? This question must, in consequence of our foregoing, be answered with yes. If
spirit as spirit could not suffer, then society could not suffer either; and yet it is the
sufferings of society, not the sufferings of single individuals, that history properly treats.
That the single individual can suffer no one doubts, although physiology has trouble enough
ascertaining what the suffering subject properly is. Even if one sees the impossibility of a
ganglion's, as a substance-subject, being able to suffer, one must yet grant that there can,
physiologically speaking, be just as many suffering subjects as there are ganglia in an
organism. When psychology nonetheless must assert that it is the soul in its self-unity as
soul that suffers, then this assertion by no means conflicts with the physiological
apprehension. The many immediately suffering subject-points are essence-moments, reflected in
the comprehensive unity of subjectivity, the suffering soul. In a corresponding manner it
stands with the suffering of society. Immediately seen, the individuals are the suffering
subjects, just as they, as regards the actions, are the willing subjects; but that which
brings it about that in the relations of society, through the many egoistic single wills, there
develops a ruling general will which has its essential and therefore reflected subjectivity in
spirit, the same also brings it about that out of the many individual and disparate sufferings
there arises a common suffering, a many-sidedly reflected suffering of society, which has its
subjectivity in the suffering spirit. From this it can be seen --- what a well-understood
concord of physiology and psychology sufficiently shows --- that life-subject, soul-subject,
spirit-subjectivity are not physically existing central links, but actual resultants of
nature-conditioned selfhood-functions, with the distinguishing feature whereby they
distinguish themselves from all physical resultants, that they react with weaker or stronger
centralized self-determination.

It goes with the suffering spirit as with the acting; it has in reality peculiar organs.
Although all individuals are, according to the possibility of endowment, organs for the acting
and suffering spirit, it is yet far from being the case that they are so all in equal degree.
The individuals have different endowment; the particular natural gifts are the peculiar
presuppositions of spirit's development: the suffering spirit is just as essentially
conditioned by natural gifts as the acting. Since spirit as spirit is self-activity, its
suffering can indeed not be without acting; but the distinguishing feature is that all its
acting is imaginary. It lets the real be what it is, and creates for itself a world of its own.
Like a migratory bird it flies away from the grey, cold lands where the flowers die and men
toil under a hard necessity; it knows warmer, brighter regions, the land where the wonders
dwell, where beauty plays at the abyss of horrors, and the bird of freedom sings in the tree of
life. Fantasy is born of suffering spirit.

What is fantasy? Intuition-capacity, image-forming activity, the power of imagination is not
yet fantasy. Intuition coincides immediately with the object; the power of imagination takes up
the contingently given and transforms it contingently; but fantasy assimilates what it takes
up, and shows itself in its transformations inspired by mood, immediately moved by end. That
inspiration and spirit essentially belong in fantasy's products is seen from the fact that it,
in all its arbitrariness, must always follow a higher law. Fantasy works creatively, in that it
transforms the natural, not into something unnatural, but into a new nature, enriched and
purified in form; in the same moment as it, by introducing the improbable, indeed the
impossible, makes a manifest breach of the laws of actuality, it knows how to grasp the essence
of things and lay bare the inner in such a way that the improbable yet becomes probable, and the
prosaically impossible is hidden by richer possibilities. Fantasy teaches us to make a
difference between real and essential truth. When it seems as though it, in the treatment of
stuff and material, disdained the real and followed only its own promptings, then it happens
solely for the reason that in the phenomena of actuality there is so much lacking, so much
diffuse, contingent, and superfluous, that the phenomena only imperfectly correspond to the
essence: it is thus in the name of essence and truth that fantasy sets its own phenomena in the
place of the real ones, and gives the imagined so deceptive a likeness to actuality itself that
we are enchanted by the illusion.

Fantasy works originally without all arbitrariness; one can with truth say that it is
immediately moved by spirit, in so far as content and form in its products correspond to each
other. The fruitful fantasy, moved by spirit, does not then fall outside feeling, thought, and
will either; although fantasy is contemplative, it stands in the most intimate connection with
the life of feeling, and although the world of images is its proper element, it yet holds a
wealth of thoughts. That willing and striving are not foreign to its essence either is seen
from its connection with drive and passion. Like every other function of consciousness, fantasy
begins in the unconscious, but has the peculiar feature that the objects it produces can set
themselves in full illumination without its becoming conscious of them as its own products; it
is the magical element in fantasy.

In what sense fantasy can be said to be the organ for spirit's self-liberation is recognizable
from the mythical gods. One can speak of a world-pressure that has rested heavily on the human
mind, of a surplus of ideality in the endowment of human nature, of secret anxiety before
invisible powers, of intimations, of fear and hope in the fermenting, unclear interior; one can
go into the particular and bring out the psychically-natural urge to personify the natural
forces and to find the human feelings, thoughts, and passions mirrored in great nature; one can,
in dim historical traces, demonstrate how strong impressions of existence have set fantasy in
motion and occasioned an unconscious, half-conscious god-poetry, how this poetry, which began
piecemeal, gradually, by added invention and re-invention, in relation to local surroundings and
the peculiarity of the tribes, rounded itself off into different cycles with different stamp. But
when all this is brought out, discussed, and clarified, one has yet only reached a
parceling-out of that which has its origin from original unity. It is spirit, the suffering
spirit pressed together in the breast of the human race, that has given itself air in an
outburst of power. The high figures of the gods are only products of fantasy, in so far as they
are spirit's own self-productions, that is, its revelations. The difference between lower and
higher revelations is then to be seen from this standpoint as an unmistakable sign of spirit's
lower and higher consciousness of itself, its goal, and its freedom.

It is not, either from the aesthetic or from the ethical side, pure ideals that come to view in
the primitive god-revelations. It is the concrete spirit which, far from beginning as abstract
spirituality, with all the marks of naturalness raises itself up from its ground and begins its
self-liberation; there are monsters among the old gods; even those especially guided by human
spirit are not pure models of virtue. The freer these gods move, the more clearly it is seen
that they go on the leading-strings of necessity, that they are under fate. Nonetheless they
open to consciousness a wide horizon, open it to a higher light, and show themselves surrounded
by the nimbus of beauty, power, and freedom. As the human consciousness, by the help of the
gods, makes progress, the gods themselves make progress; the formless disappears or is forced
into the background; the superhuman figures become gradually calmer, in that they thereby become
more spiritual, more rational, more ingenious; compare, for example, Aeschylus with Homer.

With the unconscious, half-conscious illusion as presupposition, reflection can little by little
work its way forward to higher reflection. Thus arise art and poetry with the many ideals of
life. The assumption that the ideals are exclusively determined by \textit{a~priori} independent
ideas, which demand a realization in the material over which finitude disposes, is abstract and
one-sided. The concrete apprehension of the ideal always bears the stamp of the ideal character
of the age in which it comes to be; it is the spirit suffering in actuality which, by a peculiar
tension of overreaching demand and inhibiting barriers, seeks and finds the expression of
liberation in the intuited ideal.

``No one is blessed who has not seen Zeus at Olympia''; this outburst of Greek enthusiasm says
more than diffuse explanations. To understand it thoroughly, one must be a Greek. The deepest
urge the Greek mind could then hold was the urge to see the gods. The figure of the god was, in
a wonderful fusion of human and superhuman, the perfect revelation of man himself: illusion and
yet truth, essential truth! If one considers now that what could in stillness press on the Greek
mind was precisely a lack of clarity about the free, completed individuality, that this lack of
clarity in the mind was like an underworld darkness, then one can approximately understand that
Phidias's Zeus must work liberatingly, that the sight of the father of gods and men in full
majesty must be a blessedness. In a corresponding manner it stands with the church ideals in the
Middle Ages. To see heaven open and ``God's Mother,'' the apostles, and all the saints gathered
about the throne was a grace that fell to the lot of only a few; but all felt the urge for such a
revelation. If images of saints, especially images of Mary, were to replace the lack, then the
spirit of Catholicism had to recognize itself in these products. The common people's raw
superstition could indeed use old, smoke-blackened caricatures just as well, indeed perhaps even
better than works of art; but the more spiritual the sense was, the greater was the urge for
inspired images. The art of painting stands in just as intimate a connection with Catholic
Christianity as sculpture with Greek paganism. The Scholastics had long-lasting disputes about
the explanation of how Mary as a virgin could give birth to Christ; it had to be comprehensible,
but no one, not even the Pope himself, could comprehend it. It was painful, not for the
simple-hearted believers, for they believed in the Church's authority --- whatever it was without
concept --- but for the thinkers, especially for those who thought, in Plato's spirit, of a
humane reconciliation between the natural and the supernatural. What a liberation when art at
last reached the point where it, in the conscious form of illusion, was able to reconcile the
otherwise irreconcilable. Raphael's Madonna! Yes, she is in truth mother and the pure virgin; so
human and yet divine, so sensuously beautiful and yet so exalted; one must be a humanist and a
Catholic fully to grasp the blessedness in this revelation. We can scarcely designate a Leo X's
rapture more aptly than by recalling the words with which Goethe ends his Faust: \textit{»Das
Unbeschreibliche, hier ist's gethan, das ewig Weibliche zieht uns hinan«} [The indescribable,
here it is done; the eternal feminine draws us upward].

Poetry, the art of the word, which can hold a whole worldview and partly directly, partly
indirectly express itself in thoughts and reflections, shows in a comprehensive manner how the
relation between the ideal and the real varies with each age's particular demands. In poetry,
especially in the fairy tale, we have a striking example of what it means that fantasy is, even
in the general sense, the medium for the self-liberation of suffering spirit. The further
development and scientific grounding of the principal propositions here expressed must of course
be sought in aesthetics; but the question decisive for the theory of science is this: in what
sense can aesthetics properly be said to be a science? In so far as it is built on a historical,
especially on an art-historical foundation and confines itself to dividing and characterizing
the different kinds of art, there can be no doubt of its scientific title. Nor can anything be
objected against aesthetics subjecting the different works of art to a critical valuation; it
thereby sets them forth in a clearer light; it demonstrates what the artist has intended and what
he has attained; it draws attention not merely in general, but even in the finest transitions,
particulars, and shadings, to how a single work of art can mirror the physiognomy of an age with
all characteristic traits; it teaches us that the mirroring and the mirrored correspond precisely
to each other, because both products are subject to the law of reciprocal action: the age's
spirit produces art and poetry in a peculiar style, and art and poetry act in turn determiningly
on the age. But aesthetics cannot content itself with art-historical and literary-historical
reflexes; if it will not be like the moon, which shines only with a borrowed light, then it must
have a light in itself. This light must necessarily have the idea of art and poetry for its
original source, and the question then becomes how far aesthetics is able to give a scientific
account of this idea and its relation to spirit.

According to logic the idea is a syllogism and its essence teleological; but in the syllogism here
spoken of the mediate has, of course, disappeared in parceling-out and piecing-together: the
extremes interpenetrate each other in the unity of the middle. The aesthetic extremes are essence
and phenomenon, and the idea of beauty the content-filled unity. Every syllogism rests upon a
subjectivity that accomplishes the syllogism; this subjectivity is in the individual sense the
artist, the poet, but in the essential, substantial sense the human spirit. From this standpoint
the investigations must be conducted.

One has psychologically asked what kind of drive it properly is to which the origin of art must
chiefly be referred, the drive of imitation or the drive to free, independent production?
Sculpture and painting seem to speak for the first, architecture and music for the latter of
these assumptions. Yet here too it shows itself that the boundary between imitation and
independent production is not critical, but dialectical. A flower, a landscape can be reproduced
in a picture with photographic exactness, but such an imitation is not art. The extremes of the
idea, essence and phenomenon, the ideal and the real, are lacking. But, one asks, how can such
extremes come to view in such natural objects as a landscape or a flower? The artist shall, it is
said, hold strictly to given nature, hence to actuality; but --- it is added --- he shall yet
also, by ingenious apprehension, idealize nature. On the one-sided emphasis with which now the
first, now the latter of these demands is brought out rests the opposition between realism and
idealism in art. Thus: the extremes are here. The enclosing unity of the middle, which the idea
demands, can be attained neither by increasing the ideal to the real, nor by letting the one
extreme weaken and consume the other: the unity of the middle is a unity of duality. The ideal
comes to its right only when it shows itself as natural actuality's own inner, concrete essence;
the real acquires the significance of true reality only by becoming a phenomenon that is worked
through, filled out, and shone through by its inner.

That the unity of the syllogism in all its immediacy is infinitely reflected appears from the
fact that it is by no means a presentation of the pure, blank idea of beauty to which art and
poetry confine themselves, but that actuality as elementary actuality, in bold, daring, sappy
strokes with a multitude of particulars and apparent contingencies, must, in consequence of the
idea itself, be taken along. The idea is like the light of day, which does not directly show
itself, but the objects. Thus not merely the beautiful, but also the ugly, not merely the
exalted, but also the petty, the ridiculous, not merely the harmonious, but also the
disharmonious, the conflicting, irreconcilable, indeed the horrors of conflicts, can find a
place in the artistic presentation. When one nonetheless demands, and with full right, a
resolution of the dissonances, a rounding-off and harmony in the whole that satisfies the sense
of beauty and awakens a contemplative pleasure, then one understands, consciously or
unconsciously, that the aesthetic idea-syllogism must necessarily be teleological. Only the
governing end of the idea justifies the use of the means; the feeling in which the reflection of
end and means finds its immediate expression is the mood.

In so far, then, as we hold to the general form of the syllogism, everything goes scientifically
from the hand; but when the question of the subjectivity that in a definite case at hand can with
truth be said to be competent to judge the extremes and accomplish the act of syllogism, then
difficulties arise in quantity. It lies nearest to hold to individual subjects and appeal to a
general sense of beauty. We have then on the one side the artist who produces the work of art, on
the other the many who step up to enjoy and judge the product. If it now happens so fortunately
that the work of art in each and all satisfies the demands of those who enjoy, then this could
indeed easily be apprehended as a factual proof that the demands of the idea were in all ways
satisfied, and that no other proof could be required. The scientific character of aesthetics
would then confine itself to the modest task of explaining to the artist and his admirers what
it properly is they feel, have felt, and ought to feel, a task whose solution can have its worth,
in so far as it must always please everyone to see his emotions set forth in the daylight of
consciousness and received with applause. We call it nonetheless a modest task, for suppose that
the aesthetician attributed to the artist and his public feelings they would not acknowledge, so
that the growing explanation rested on a misunderstanding: can the aesthetician then prove that
his apprehension is the true one? Let us assume that the relation is another, that in the public
there is a party that praises, another that censures the work of art; the aesthetician with his
critique is then to mediate between them. But where is the firm foundation that secures a
starting-point? It does not stand with the self-evident in aesthetics as in logic, in
mathematics, or in physics. One can with a certain tact reckon on a general estimation and, by
the help of assertoric judgments passed off as apodictic, seek to get as many on one's side as
possible; but the conduct of the proof is just as little scientific as if one called together a
general assembly, set forth one's opinion in the form of an assertion, and let those in favor
raise their hands in the air.

So long as individual subjectivity is to be the one valid thing, every individual has an
unmistakable right to appeal to his feeling and his estimation; here holds the well-known saying:
\textit{chacun son goût} [each to his own taste], and the discussion between the many who
disagree ends with the negative result: \textit{de gustibus non disputandum} [there is no
disputing about tastes]. To wish to let the idea itself judge between the contending parties is
vain, since everyone can apprehend the idea in his own way. The only way out, when aesthetics is
to be maintained as a science, is to set out from impressions of nature and appeal to the human
spirit; we shall in brief elucidate what this means.

An unconditioned \textit{a~priori} idea of beauty, resting on itself, does not exist. All forms
of beauty are, essentially seen, teleologically apprehended natural forms; this holds of the
inorganic as well as of the organic. The representation of an objective beauty of nature is just
as subjective as the representation of an ugliness of nature. The separation of the beautiful and
the unbeautiful in nature rests upon a judging subjectivity that makes a demand. Although the
demand is subjective, it is yet not arbitrary; it has its objective ground in nature itself. Here
it is presupposed that a fundamental form can be more perfectly expressed in one natural product
than in another, the species more purely and completely realized in one specimen than in another;
the subjective demand is to see the species, the typical form, realized as perfectly as possible
in its specimen. Such a teleological demand has no validity in physics. The opposition between the
ideal and the real reduces itself in physics, as we have seen, to an opposition between
real-concept and existing thing. The existing thing is a realized concept, and the concept is
realized in the thing as perfectly as the conditions have allowed. In aesthetics, on the other
hand, the real-concepts are apprehended not as concepts, but as typical forms with a stamp of
endowment that can be realized in a more imperfect or more perfect manner. From the standpoint of
physics the Australian aborigine's real-concept must be said to be realized with the same
necessity in the Australian as the European's real-concept in the European; here there is no
other measure for the perfect than necessity. When aesthetics, on the other hand, declares that
the European is more beautiful than the Australian, then it happens not on account of individual
feeling or an immediate estimation that is beforehand taken and captivated by the typical
advantages of its own race; but it happens because it is here presupposed that the fundamental
type of the human race is more perfectly realized in the European's than in the negro's figure.

Can this now be proved? Not in the manner of physics. If the physical mode of proof is to hold as
the only true scientific one, then aesthetics is no science; but if subjectivity, with its
teleological demands, is allowed to come along --- not as individual subjectivity in all its
contingency, but as the subjectivity that is on a level with the whole nature-conditioned
development of spirit in history and with the suffering spirit's self-liberation in fantasy ---
then the whole horizon changes. The general subjectivity, the highest judge of beautiful and
unbeautiful and moreover of the essential significance of art and poetry, becomes thus the human
spirit itself. This judging subjectivity is not an empty abstraction, a phantom one can turn and
twist at pleasure; it has concrete existence, flesh and blood, in the individuals who through the
particular breathe in the general, and whose aesthetic judgments become more and more objective
in proportion as those who judge become, in spirit itself, deeper and deeper subjective.

\section{Self-Liberation in the Medium of Will: Ethics}
\label{sec:24en}

\noindent What subjectivity is, and what significance the question of the reality of the
subjective must have, also for science, is striking when one looks at the will. From a
physiological objective standpoint it is impossible to explain what will is, and just as
impossible to deny that there really is something called will. Since now the will is and
continues to be as a fact admitted by all, whatever one may otherwise judge of it, there must,
although the will is the most subjective of all, be able to be an objective knowledge of the
will. What then is here in objective knowledge self-evident? We answer: the will presupposes a
willing subject, an act of will, and an object upon which it goes out.

That the willing subject can be neither a single substance-subject nor a sum of mutually
attracting and repelling substance-subjects is, so far as we can see, sufficiently demonstrated
in the foregoing. It is vain for the physiologists to plague their brains with willing
substance-subjects, for the day they find them, the theologians will already have found the
dogmatic-speculative proof for the over-speculative resurrection of the dead. It stands with the
willing as with the knowing subjectivity: it is a self-evident, psychological fact. On this
presupposition a scientific explanation can be laid out. That the will as will must originally be
free follows from the fact that every act of will, essentially seen, is an act of
self-determination, and conscious self-determination is in its ground freedom. With its organic
conditions the finite will is surrounded by barriers, which it can only in consequence of a
representation of alternation partially cancel by preferring the one case to the other, that is,
by choosing. On freedom of choice rests arbitrariness. In a state of unconsciousness
self-determination can express itself only involuntarily; in such involuntary expressions as
when a fetus in the mother's womb moves its limbs there is neither will nor freedom. In freedom
of choice and the corresponding arbitrariness we have the first elementary indications of the
life of freedom.

There is in freedom of choice a peculiar union of free and unfree. Let A and B be two
alternatives between which a choice must be made; then there are properly three cases possible;
for the will has not merely a choice between A and B, but it has also a choice between choosing
and not choosing. Although the last alternative seems to cancel itself and become an empty
semblance, it yet plays an essential role in the life of freedom, indeed a far more essential
one than one would at first glance suppose. If there were only a question of finitely definite
alternatives such as either to shoot oneself or another, then an outward superior power could
always compel a choice, and by such compulsion not merely restrict, but even annihilate,
freedom; but the choice between choosing or not choosing the outwardly prescribed no power can
compel. This choice rests solely upon the will's unity with the willing subject, and in this
consists the unforfeitable freedom. The firmness with which subjectivity is able, despite the
highest degree of compulsion and temptation, to hold its will enclosed in itself and hold back
every outburst of outward-going choosing or acting, is the greatest test of strength to which
the energy of the will can be subjected; but it is at the same time of so subjective a kind that
it withdraws from every objective control.

But however high the energy of self-determination is raised to a power, even if it is thought
omnipotent, it is yet impossible for it to free itself from motives. No one can will without
willing something; even to exclude every regard for the outward objective, in order merely to
will one's own will, is to will something. Since now the self-determining will cannot move itself
without being moved by the something that is willed, the understanding must from the objective
side, with its dissuading and advising grounds, evaluate the objects of will, and from the
subjective side illuminate their relation to feeling, desire, inclination, and passion. Thus the
will stands surrounded by the different motive-grounds, that is, by the motives. When the
understanding, under all-sided deliberation, has weighed the one motive against the other, by
what then is the will moved? By the strongest motive, one answers. This answer is correct; but
it chiefly matters to see how it properly comes about that one of several motives becomes the
strongest. If the cold understanding were allowed to rule, then the motive would draw its
strength from the essential worth of the object of will for the willing subject. But desire,
passion, the momentary lust can cast their lot in the scale and give another motive the
preponderance. Could now the will, as a disinterested spectator, keep itself quiet and passive
under the alternating struggle of the motives, then it would blindly be dependent on the outcome
of the struggle; it would become a prey for the victorious motive. But it is not so. The
self-determining will is from the beginning in on the deliberation, and takes so lively a part in
the alternating struggle of the motives that it chiefly depends on itself which motive is to
become the strongest. It is thus not enough to say: the will is determined by its motive; one
must go a step further and add: the will itself is involved in determining the motive by which
it is determined.

From the reciprocal action between motive and will is explained a manifold of phenomena in the
life of freedom that would otherwise be unintelligible: the will's dependence and independence
both on impressions from the surroundings and on the nature-ground of individuality, countless
gradations of strength and weakness, guilt and excuse, accountability and unaccountability.
Pervasive is especially the opposition between what we may call the individual's peripheral will
and its central will. The peripheral will varies with the changing impressions and prefers to
follow the momentary lust; the central will demands the self's agreement with itself in each and
all, a consistency of character, a carried-through plan. From this it will be evident how
difficult, indeed impossible, it is for the individual will to bring its elements of life into
harmony and give freedom shape in flesh and blood. The antique virtue would succeed as little for
the Stoic as for the Hedonist; the strife between the individual and the general ended in
self-contradiction and screaming disharmony. When, on a far later stage of development, one came
to consciousness that the will of reason must unconditionally master the drive-will (Kant), the
doctrine of duty became so intolerably dualistic that drive and reason united in revolt.

The negative critique has then undertaken to free individuality, not merely from the supernatural
God, but even from the eternity of the self and the good as an idea valid in itself. Although this
critique has brought several things to light that science will indeed take \textit{ad notam}
[into account], it is yet a peculiar matter with the negative critique. It cannot possibly be the
earnest of the critique that the enlightened one who emancipates himself from religion should also
emancipate himself from ethics; on the contrary! one attacks precisely the religious in order to
purify the ethical. Instead of making, for example, justice dependent on the representation of
God, the critique makes our representation of God dependent on our idea of justice; instead of the
religious article of faith ``God is just,'' it sets its ethically scientific proposition ``Justice
is divine.''

The idea as idea was thus to ground the individual's ethical self-liberation. But how difficult it
is for the idea of the good to maintain its objective validity when the eternal self-being of
spirit is denied, and human subjectivity held exclusively from the standpoint of temporality and
mortality, is easy to see. The idea of the good is an idea of personality that has not and cannot
have unconditioned validity in any abstract objective form whatever. In and for itself there is
nothing absolute, either good or evil. Only in and for the absolutely Good is the good absolute;
if the Good is not itself eternal, then there is no eternity in the good. Such ethical
propositions as: ``the good is divine,'' ``justice is divine,'' must then either seek their
grounding in the religious: ``God is good,'' ``God is the Good,'' ``God is just,'' ``God is the
Just,'' or, as positivism has shown, transform themselves into turns of phrase, and the result
become that there is no absolute at all.

Also from a purely human standpoint it is seen that the idea of the good must be an idea of
personality; it is seen from obligation. Let the idea for itself be eternal, human subjectivity
temporal: how then can a subjectivity that is in its essence temporal, mortal, perishable
acknowledge demands of life that call for an eternity which does not concern it, an existence in
which it cannot itself become a partaker? Could the self-consciousness confined to the present life
speak to the idea in human fashion, then it would have to say: go you with your principle of virtue
and your ideal commandments of duty to the gods, and teach the Immortals to do the good for the
good's sake, for that eternal beings live for the eternal is in order; but leave us mortals in
peace! Or is it perhaps not just as unreasonable that a mortal will make the eternal his highest,
as if an immortal took it into his head to live for sheer temporal things.

Thou shalt not kill! This commandment, which in a culture-world's rational ethics must be reckoned
among the fundamental ones, loses its unconditioned validity when subjectivity recognizes itself
only as temporal. Why may the one man not murder the other? From the standpoint of negative
reflection the answer must run thus: where murder and killing are practiced unpunished, no society
can subsist, no development of spirit take place. Only --- it is said --- when the development of
spirit has reached its highest, the self-conscious subjectivity recognizes its own mortality; only
on the stage of enlightenment will man see that religion's represented God is the eternal idea of
ethics, and in a far deeper sense apprehend the obligation to unconditioned obedience to the idea,
it is said from the side of the idea's defenders. But how self-contradictory such a defense, despite
its apparent thoroughness, must always be, is easily seen on a little closer deliberation. The
obligation to observe the commandment: Thou shalt not kill! can in a civilized society indeed,
without express regard for religion's doctrine of God and eternal life, be inculcated from so many
sides, grasped and felt in so many ways, that moral cultivation in union with the civil laws can
--- so long as the order of society is maintained --- furnish sufficient guarantee for life and
property. But in this there is yet contained no rational answer to the question on which it
properly depends.

It can indeed, so long as there is tolerable order, present itself as though the individuals with
the clear but perishable self-consciousness set the idea as the unconditioned; and should there be
individuals who will not in the name of the idea bow to the inviolable right of the existing whole,
they would soon succumb to the general law and disappear. But although the law can master the
disobedience of single individuals, so long as the decisive majority of the temporal
self-consciousnesses is prepared to risk all for the unconditioned demands of the eternal idea, it
will yet, when society is only temporal, gradually dawn on all individuals how impossible it is to
respect man's eternal, inviolable right and be conscious that in man's self there is nothing
eternal. So then all is temporal, the good as well as the evil, the noblest as well as the basest,
right as well as wrong; each contends against wrong for what he calls his right: what shall society
now do?

It must gradually dawn on those contending in society that the single one's right cannot possibly
rest on his own demands; for if the single one had without further ado the right to drive all his
demands through, then the last remnant of the state of society would have to disappear: right as
real right exists only in real society. And yet the existing society has no right either merely
because it exists. If an order of society had the right to subsist unchanged, the state of right
would with that same disappear; society thus has the right to subsist only in so far as the
individuals have the right to transform and develop the existing according to their justified
demands. The right of society, right in itself, holds chiefly to the existing, while right as the
individuals' right, the moral right, holds chiefly to an advancing change, an ever further
development of the endowment of the existing. There must in society at every time be a definite
order for the working classes, but this order must in turn be able to change in relation to the
need of the individuals. The clearer the consciousness is of the demands of the existing in all
directions --- for the great multitude, which is chiefly referred to living immediately by the
works of its hands; for the well-to-do classes, which by their powers and cultivation set the great
works in motion and carry them out, both for the welfare of those nearer and for that of the whole,
in such a way that they serve both the individuals' and society's ends; and finally for the
position of those who, by particular endowment and cultivation, are called to promote the aim of
enlightenment for all estates in different measure, and thereby the higher development, both
theoretical and practical --- the clearer will society know how to maintain its right and its
bearing. With regard to all this, the individuals must nonetheless preserve an original right to
raise objections, without which the freedom, the full freshness of presence, which alone is able to
keep life in the whole movement, would be lost.

All this cannot possibly be kept in such a course that the nearest, practical ends are promoted,
without spirit being involved --- the spirit that explains both the individuals' and society's
right and points out the way of development. What kind of spirit this is, and whence it comes, has
been shown in the foregoing. It goes through the individuals and reaches its development, in that
the individual subjects become serving moments, bearers of the general subjectivity. As the
individuals of society, each according to his estate, his position, and particular giftedness,
develop and bear the subjectivity corresponding to spirit, these individuals furnish the common
foundation on which the development of spirit, which promotes and secures the welfare of the whole,
essentially rests.

On the difference of the working individuals there indeed rests that activity of spirit on which
both the welfare of the whole and that of the single one is essentially grounded; but when it comes
so far in a society that the conflicts are in a higher sense recognized as conflicts of spirit, then
the laws that govern these conflicts --- the worker's as well as the employer's, the activity that
stirs in the possessing and the acquiring, and that which comes to view in the energy of the
theoretically cognizing and the artistically forming --- will strive forward to a common unity, a
higher goal for the freedom of spirit and of individuals, a freedom recognized by all, since it
proceeds from all, and finds with all its ever richer and ever more essential abode.

As free activity is more and more recognized as the inviolable right of spirit, a right that
extends to all and on given conditions satisfies the demands of the whole society and of the
single ones, the single complaints will sound as spirit's complaints, and the single punishments
be felt as spirit's punishments; but spirit as knowing, self-conscious spirit will sink itself
deeper and raise itself higher, in proportion as it gains more and more entrance both into society
and into the single one, into the lowly and poor as into the richest and most gifted.

If we single out the different gifts of spirit, it lies in the nature of the case that they must
all allow themselves to be referred to the natural endowment of individuals. The common spirit
that is in a society must work its way forward out of gifted individuals, and these in turn be
helped forward by spirit. Although each individual begins immediately with its own subject, it can
yet in the main not avoid furnishing subjectivity for the life of spirit; it is this common
subjectivity, to which the individuals practically and theoretically contribute, and from which
they in turn draw their nourishment. The ethical law that determines the development of the spirit
of individuals and the spirit in society must be regarded as the fundamental law of development;
but the natural conditions for the development of the fundamental law are so manifold that they, so
to speak, lead to disguises, to determinations of objective laws, to concessions of subjective
demands, under which the life of spirit must contendingly move forward.

Here then arises the question whether the human spirit, which in the proper sense is the spirit of
knowledge and science, is also in the same sense the spirit of life, of faith, and of hope; whether
it can with truth be said that the religious, which began so superhumanly, more and more reshapes
itself into the human itself, and enters into one with the human spirit, or whether it should lie
in the human spirit itself, the more completely and deeply it develops according to its own essence,
the more inwardly to feel and grasp its need for a life of spirit on higher conditions, so that an
absolute spirit, which is not merely the unity of knowledge and power that is truth's, but the
personal love's unity of knowledge and power, should furnish that on which salvation rests. We
shall in the name of science consider this spirit, which, if it has validity, must stand infinitely
high above science itself and all human knowledge.

\chapter{The Absolute Spirit}

\section{Science and the Absolute Spirit}
\label{sec:25en}

\noindent That the absolute spirit cannot possibly give itself concrete form or finite shape in
science must, in consequence of our foregoing, be regarded as settled; but, of course, so long as
such a spirit nonetheless finds belief among men, it cannot wholly withdraw itself from science's
critique.

As regards the human spirit, it is impossible for it to be in everything subject to science. The
scientific control to which the human spirit is subject can never keep pace with the development of
life itself, but must as a rule follow behind. Let the prescriptions a society can bid its
individuals to follow be ever so universally valid and worthy of recognition, they will yet from
the individual standpoint only fix certain boundaries within which the life of single individuals
moves with unviolated freedom. In the manner in which the promptings are expressed and followed we
have essentially the true, half-true, and untrue indication of the demands of the life of spirit.
But when these demands are subjected to a comprehensive critique, life has in reality moved
further, and practically passed its judgment, in order itself one day to be thoroughly judged, when
its time is over.

If, then, the question arises of the relation between the abiding and the vanishing in life and
action, the answer will be that what immediately concerns the individual is perishable; only that
which as the general belongs to spirit can approach the imperishable. Skill in using firearms will,
for example, belong to the perishable individuality; the use, on the other hand, that can be made
of it in the fatherland's service will be less perishable; but love itself, the sacrifice, will be
the abiding, the imperishable in spirit, the gleam of immortality that rests over mortal
individuals. Skill in acquiring the goods of life can distinguish certain individualities specially
disposed for it, but only in so far as these goods are used for the common good, and thereby acquire
a spiritual worth, do they cast a gleam of immortality over the acquirer. Knowledge and insights
glorify the individual who possesses them; but only when they are used in higher service, for needed
enlightenment and true cognition of the goal of life, do they receive spirit's consecration to
imperishable honor.

How does it now stand? The question of the in-truth eternal life must decide the matter. Is there
then in this anything eternal? By no means, neither for the individuals nor for spirit. That the
individuals themselves have no share in the eternal follows simply from the fact that they must all
die, and that death wipes out consciousness, so that spirit must take refuge in new, perishable
individuals. But the spirit of perishable individuals, though it can wander enriched from generation
to generation, is itself perishable, and can only through the succeeding individuals pass its
judgment on the preceding. Can such a judgment now be satisfactory in each and all? By no means! It
is satisfactory only for those who to such a degree have their life in the perishable that the
semblance of imperishability that spirit can spread over it satisfies them. So they die, the bad
and the good, the happy and the unhappy, and death, which lifts consciousness away, makes them all
equal. The gleam of justice that posterity can spread over the deceased is, with all possible
striving after truth and thoroughness, yet so fleeting and vanishing that in many highly important
respects it had best be absent; it also sinks with time more and more, and is at last buried in
oblivion.

One can then, in the name of science, not be offended, indeed not even astonished, that there are
still, among the knowing --- even though they carefully follow the way of scientific development and
grant the results already gained great and general validity --- those who, with respect to life's
personal task, seek and find their proper home in the realm of the absolutely personal spirit. Here
there can naturally be no question of breaking off the scientific deliberations at any arbitrary
point whatever, since science must on the contrary feel itself called to investigate everything that
in one way or another falls under consciousness. A doctrine with overscientific fundamental
propositions science cannot possibly leave standing untested; to wish, by an arbitrary dictum, to
withdraw certain things from science's investigation is a proof of misjudgment not merely of
scientific truth, but of all truth. It is otherwise when science itself is to judge. It does not
judge arbitrarily; on the contrary, it summons everything before its tribunal in order to
investigate what it can really judge, and what it is that by its nature must withdraw from its
judgment.

Can then science in spirit and truth pronounce any decisive judgment on the supernatural being in
and for itself? In consequence of all that has been shown in the foregoing, the assertion of this
must be regarded as invalid. Science is in the strictest sense able neither to deny nor to confirm
the validity of the supernatural in itself; but as to the manner in which its outward appearance
takes place in this natural world, it reserves to itself a right to judge. It turns with
irrefutable critique against the Old Testament in the doctrine of creation, of paradise, and of the
manner in which it is lost. Where the boundary between the natural and the supernatural lies it can
indeed not directly demonstrate; for the natural itself has, as we have seen, marks of the
supernatural. The power in other rests upon the self-power, and the self-power is just as well
supernatural as natural. On the other hand, the relation of creation must expressly be said to be
distorted by the one-sided appearance of the supernatural. Not to speak of paradise's miraculous
combination of supernatural and natural, there is here, to name the relation between earth, moon,
and sun, a relation that science must be able to judge, if it is to be able to judge anything,
something unwarranted in the creation-doctrine of the first book of Moses. If one will not grant
science the right to undertake its astronomical fundamental change, then scientific authority falls
to the ground with that same; science then gets no other right than Scholasticism had in the Middle
Ages. If, on the other hand, one grants science its full right when the question is of the
astronomical, on what ground will one then be able to deny it the same right when the talk is of
geological development? That the theory is not yet finished science itself sees most clearly of all;
but if one but grants its methodical procedure objective validity, then it is with that same
conceded that it neither can nor will submit to biblical representations. And now the last grand
attempt with the descent of the higher animals from the lower and the descent of men from the
animals! Can anyone really suppose that this can be combated with scientific weapons in favor of
the biblical-religious apprehension, or by some ingenious modification be brought in under the
old-biblical view? Vain! --- With the outward reality of the supernatural in the New Testament it
goes exactly the same way: the scientific apprehension stands at all outward miracles sharply
against the religious.

What is now to be done about this? To wish, in the name of religion, to declare war on science is
by no means admissible; everyone who has any concept of the nature of knowledge sees this. To wish
to hold to a religion that has taken on a scientific form, or to a science adapted to religious
demands, is not admissible either; for such a bastard must from both sides, from religion's as well
as from science's, be seriously combated. What then is here to be done? Should there be anything
other than to demonstrate the religious as the generally-religious, which has its own principle of
cognition, to let the in-truth scientific go its way unhindered, and from both sides to guard
against unwarranted encroachments.

Religion and science have in their temporal development inherited a mutual inclination to weaken and
harm each other reciprocally. In the Middle Ages the papal religion obtained, by continued struggle
against worldly resistance, at last such a preponderance over the powers of the world that it not
merely practically, but even theoretically dominated everything; had it in addition been competent
in spirit and truth, there is absolutely no doubt that it would have struck every opponent to the
ground. The resistance that was essentially to weaken its power came in its beginning not chiefly
from the outward, but essentially from the inward. To the inward resistance the outward gradually
joined itself; how decisive the outward was is seen not merely in the lands where the state power
directly protected Protestantism, but also there where, especially as in France, the papal power
more and more had to bow to the royal power. With the carrying-through of the Reformation, science,
with its original right of investigation, began to press forward into the foreground. Philology,
history, and philosophy struck out ever further in free investigations; at last the doctrine of
nature came to full consciousness and led the way on the paths already broken. By the Revolution
the freedom of peoples made, as well as it could, common cause with scientific freedom, and a union
of the nature of both is recognizable enough on the whole.

But the freedom in the outward has left the inward freedom considerably behind. While there is
struggle in all directions for political and social freedom, the question of inner personal freedom,
the inner, true life of freedom, which religion alone can give, has remained comparatively little
developed. To raise up again the sunken freedom, a reawakening of the old church-faith will be just
as insufficient as a union of scientific and worldly social endeavors. A reawakening of the old
church-faith will stand in manifest conflict with the development of spirit on the whole; and a
restriction to the merely scientific and social freedom will, when nothing else is at issue, end by
bringing humanity into revolt. The Socialists have already begun to show the example. If the true,
original life of freedom is not to perish in humanity's struggle to acquire the conditions of
freedom, then it must breathe in religion and live in the faith in the absolute Spirit of freedom,
which in truth is able to do what no relative spirit is able to do, to free every believing soul.
The absolute spirit does not free man by breaking through the connection of life and leading him
away from real experience and real development of society's tasks; it does indeed raise
consciousness high above the tasks of science and art, of politics and national life, but, be it
noted, not in order to give up these tasks to whoever wants them, but in order to set them forth in
their full and true illumination, and in the name of eternity to hasten their temporal solution.
Every judgment of the relative worth of the solution is naturally left to the relative spirits; but
as the eternal truth shines through the relative, the powers of the time first acquire the nobility
and the stamp of truth which they greatly need.

That the absolute spirit cannot possibly go directly out upon anything relative for its own sake is
evident; but from this it by no means follows that it should be excluded from the relative. If it
could really be proved that the relative stood in contradiction to the absolute, then we should have
to grant positivism right, which effaces the absolute in order to maintain the relative. But unless
we are completely blinded, it is precisely the absolute that is relative in the relative: the
absolute spirit must thus bear the relative, must, in sinking itself therein, at the same time raise
itself thereabove, and thereby save the human spirit from death and dissolution.

\section{Philosophy of Religion}
\label{sec:26en}

\noindent Religion, which according to its originality is necessary for man, cannot be replaced
by science; it has and must have its own doctrine of cognition, its own presuppositions
independent of science. Yet it is science's right to revise these presuppositions, in order to
convince itself that they, with their overscientific character, contain nothing that essentially
conflicts with the nature of knowledge and its irrefutable consequences. Science thus gets the
double task: to secure for itself the exclusive right to revise the objective constituents of the
doctrine of religion, and to develop the fundamental thoughts of religious cognition in such a
way that they become generally comprehensible to all.

The relation to God determined by the miracle, which is religion's own, irrefutable
presupposition, science cannot appropriate, and must therefore recognize a boundary it cannot
overstep; but the boundary of knowledge cannot possibly coincide with the boundary of thought.
There must thus be an essential difference between falling outside knowledge and falling outside
thought. The opposition between the principle of knowledge and the principle of faith must be
expressly recognized. Science would not be able to apprehend and determine its principle of
knowledge if it were not, by the opposition between what can be known and what cannot be known,
able, even if only negatively, to characterize the principle of faith impenetrable to knowledge.
But if it is possible for science to separate the absolutely disparate principles (cf.\ Part One;
\S\,1 and \S\,6), then it is also possible for it to separate the thoughts proceeding from these
principles, possible to make a difference between thoughts of knowledge and thoughts of faith,
between concepts of knowledge and concepts of faith. With the peculiar task of clarifying the
relation of conflict between the disparate thoughts, the philosophy of religion is chiefly to be
determined as a boundary-science.

The philosophy of religion does not, however, remain standing at general boundary-determinations
for faith and knowledge; it cannot possibly let religion develop its content without letting a
whole system of concepts of faith come to development. But faith's religious content is, according
to the principle, an overscientific content; here then enters the contradiction that an
overscientific content is taken up into science. If the form of science were now really confused
with the content's true, original form and introduced and held fast instead of it, the
contradiction would indeed be insoluble. But this is by no means the case. Scientific reflection
has precisely the flexibility that it can bring the form of knowledge to reflect another and
opposite form. For just as the clear day can bring the eye to forget the gleam of light and look
at the objects, so reflection too can bring thought, in knowledge, to look at faith and forget
knowledge. The philosophy of religion is so far from forcing on the religious content a foreign
form that its development of form consists, on the contrary, in freeing it from the additions of
immediately individual representations, images, feelings, expressions of will, and so forth, with
which it in the life of faith itself is always grown together. In so far as reflection, both as
regards the form and as regards the content, by knowledge cancels knowledge, in that it sets out
to clarify, not a scientific, but an overscientific truth, the relation between truth and science
is essentially reversed, and the philosophy of religion is with respect to this to be
characterized as ``an inverted science.''

\begin{quote}
``In all direct sciences the principle of the grounding of the matter is at bottom one with the
general principle of knowledge, so that the grounding of the matter coincides point for point with
the science in question's own self-grounding. In the inverted science the grounding is, on the
contrary, inverted. The fundamental principle is here so far from going up into the principle of
knowledge that the two absolutely disparate principles, that of the miracle and that of reason,
everywhere repel each other and in their negative relation determine each other negatively. Of
this negative relation, with the corresponding reciprocal determinations, a clear knowledge can
indeed now be given; but the general knowledge that hovers over the opposed principles is on the
present stage no longer able to take the part of its own principle; the inverted science has
precisely the task of setting faith in force and bringing knowledge itself to bow before it. As it
goes with the grounding principle, so must it go with the grounding. The inverted science shall
neither conceal nor deny the universally valid grounds of knowledge, but it shall bring it to
clear insight how religion's self-confirmation precisely takes place by the fact that the grounds
of knowledge consistently resolve themselves and give place to faith's absolutely self-valid
grounds. By grounding the transposition and resolution of the grounds of knowledge into grounds of
faith, the philosophy of religion grounds its own scientific dissolution, and acquires on these
conditions significance as a science sacrificing itself upon religion, cancelling itself in the
doctrine of faith.'' (\textit{Philosophy of Religion}, pp.\ 8--11).
\end{quote}

From the standpoint of knowledge it must now be striking that science only gives itself up, in the
overscientific, in order to convince itself that herein is contained nothing that can dissolve its
foundation --- thus in order to return into its own circle and maintain its validity in all that
can be known. Its investigations begin continually afresh, in that it, both outwardly in the
extension and working-up of the material and inwardly in the continued testing of the principles,
unbrokenly continues its circular movement, condemns the superficial, supplements the one-sided,
and appropriates the grounded. No outward power can condemn science: it is, in the circular
movement, its own judge.

\bigskip

{\footnotesize\noindent\textsc{Errata} (page and line references to the 1880 Danish edition):
p.\ 98, l.\ 5 from above: ``\S\,1'' --- read: ``\S\,2.'' --- p.\ 107, l.\ 9: \textit{Mediationer
naae} --- read: \textit{Mediationer aldrig naae} [Mediations \emph{never} reach]. --- p.\ 157,
l.\ 3: \textit{Axler} --- read: \textit{Axer} [axes]. --- p.\ 160, l.\ 11: \textit{Vare} --- read:
\textit{Varer} [goods]. --- p.\ 196, l.\ 14: \textit{Lyélls} --- read: \textit{Lyells}. (These
corrections, listed at the end of the original, are already incorporated in the present text.)}

\end{document}
